\section{Closure-Liquid Crystal and Holographic Julia Boundary}
\label{sec:closure-liquid-crystal-holographic-julia-boundary}

The present section introduces a cautious re-reading of the current closure architecture in cellular language. It does \emph{not} claim that the framework is literally a condensed-matter crystal. Instead, it proposes one structured pilot in which the same closure-bearing carrier is treated as a bulk order parameter whose visible sectors appear as phase-sensitive ordered domains and whose Julia-type diagnostics are read as boundary holograms of that same order. The point of this re-reading is organizational: it compresses the current OP~3/OP~5 language into one local compatibility / defect-transport / boundary-readout picture without creating a second closure theory.

\begin{definition}[Closure-liquid-crystal cell]
\label{def:closure-liquid-crystal-cell}
A \emph{closure-liquid-crystal cell} is a quintuple
\[
  X=(K,\sigma,\rho,\phi,\delta)
\]
consisting of:
\begin{enumerate}[leftmargin=2em,label=(\roman*)]
  \item a closure-bearing carrier datum $K$;
  \item a shell/refinement datum $\sigma$;
  \item a branch/route datum $\rho$;
  \item a local phase/readout datum $\phi$;
  \item a local defect datum $\delta$.
\end{enumerate}
The intended reading is that $K$ records the bulk carrier order, while $(\sigma,\rho,\phi,\delta)$ records the visible local realization of that order in one readout domain.
\end{definition}

\begin{definition}[Local cell-compatibility equation]
\label{def:local-cell-compatibility-equation}
Let $X_n=(K_n,\sigma_n,\rho_n,\phi_n,\delta_n)$ and
$X_{n+1}=(K_{n+1},\sigma_{n+1},\rho_{n+1},\phi_{n+1},\delta_{n+1})$
be adjacent closure-liquid-crystal cells. Their local compatibility equation is written
\[
  \mathcal C_{\mathrm{cell}}(X_n,X_{n+1})=0
\]
and is required to mean that:
\begin{enumerate}[leftmargin=2em,label=(\roman*)]
  \item the carrier datum remains inside the same closure-bearing class;
  \item the branch/shell transition between the two cells is admissible;
  \item visible readout variation remains at screening level and does not yet generate a genuine shell defect witness.
\end{enumerate}
Thus the equation $\mathcal C_{\mathrm{cell}}=0$ is the local cellular form of admissible OP~3 passage.
\end{definition}

\begin{proposition}[OP~3 cell-compatibility reading]
\label{prop:op3-cell-compatibility-reading}
Let a translated OP~3 shell passage carry a shell admissibility certificate in the sense of Definition~\ref{def:qam-shell-certificate}. Then it determines two adjacent closure-liquid-crystal cells whose local compatibility equation holds:
\[
  \mathcal C_{\mathrm{cell}}(X_n,X_{n+1})=0.
\]
In particular, the current OP~3 certificate package may be read as a local cell-compatibility law on the same closure-bearing carrier medium.
\end{proposition}

\begin{proof}
A shell admissibility certificate records admissible branch passage, shell-admissible refinement, and closure-compatible readout on the translated OP~3 side. By Lemma~\ref{lem:qam-shell-certificate-closure} and Theorem~\ref{thm:qam-op3-shell-closure}, the source state and the branch-shell return state remain in the same closure-bearing class and differ only by an admissible shell passage together with readout-screening data. This is exactly the content encoded by $\mathcal C_{\mathrm{cell}}(X_n,X_{n+1})=0$ in Definition~\ref{def:local-cell-compatibility-equation}.
\end{proof}

\begin{definition}[Locally corrected crystal compatibility equation]
\label{def:locally-corrected-crystal-compatibility-equation}
Let $X_n$ and $X_{n+1}$ be adjacent closure-liquid-crystal cells in the closure-phase regime. Their \emph{locally corrected crystal compatibility equation} is written
\[
  \mathcal C_{\mathrm{cell}}^{\sharp}(X_n,X_{n+1})=0
\]
and is required to mean that:
\begin{enumerate}[leftmargin=2em,label=(\roman*)]
  \item the original compatibility equation $\mathcal C_{\mathrm{cell}}(X_n,X_{n+1})=0$ holds at closure/readout level;
  \item the local phase-side persistence defect stays below the active reinforced threshold;
  \item the associated Teichmüller-side persistence defect stays below the active reinforced threshold;
  \item the packet-side admissibility monitor remains inside the safe admissibility band.
\end{enumerate}
\end{definition}

\begin{remark}[Why the crystal equation must be corrected locally]
\label{rem:why-crystal-equation-must-be-corrected-locally}
The original crystal compatibility law is an OP~3-style admissible passage law. The present post-v3.32.0 local obstruction analysis shows that this is no longer enough for the active closure-phase trigger problem. The compatibility equation must therefore be tightened by a local persistence correction, not by changing the carrier class itself, but by sharpening the phase-side and Teichmüller-side persistence conditions on the second step.
\end{remark}

\begin{definition}[Local crystal persistence correction datum]
\label{def:local-crystal-persistence-correction-datum}
The \emph{local crystal persistence correction datum} at step $n$ is the quadruple
\[
  \mathfrak K_{\mathrm{cell}}^{\sharp}(n)
  :=
  (\eta_{\Theta}(n),\eta_T(n),\zeta_T(n),\lambda_T(n)),
\]
where:
\begin{enumerate}[leftmargin=2em,label=(\roman*)]
  \item $(\eta_{\Theta}(n),\eta_T(n))$ is the first local repair witness;
  \item $(\zeta_T(n),\lambda_T(n))$ is the local Teichmüller-side repair datum.
\end{enumerate}
\end{definition}

\begin{definition}[Corrected crystal two-step window]
\label{def:corrected-crystal-two-step-window}
A two-step local crystal window
\[
  X_n,\;X_{n+1}
\]
is called \emph{corrected} if the locally corrected crystal compatibility equation holds and the corresponding reinforced trigger package passes on both steps.
\end{definition}

\begin{proposition}[Teichmüller repair lifts the crystal equation from screening compatibility to persistence compatibility]
\label{prop:teichmueller-repair-lifts-crystal-equation}
Assume the first local repair witness and the first local Teichmüller-persistence repair theorem. Then the local crystal compatibility equation can be lifted from
\[
  \mathcal C_{\mathrm{cell}}(X_n,X_{n+1})=0
\]
to the corrected form
\[
  \mathcal C_{\mathrm{cell}}^{\sharp}(X_n,X_{n+1})=0
\]
on every corrected crystal two-step window.
\end{proposition}

\begin{proof}
The first local repair witness supplies the phase-side tightening required to repair the second-step drift, while the local Teichmüller-side repair datum supplies the remaining geometric persistence correction. If both corrections are applied and the reinforced trigger package passes on both steps, then the original admissible crystal passage is no longer merely screening-compatible but persistence-compatible on the active two-step window. This is exactly the meaning of the corrected equation $\mathcal C_{\mathrm{cell}}^{\sharp}(X_n,X_{n+1})=0$.
\end{proof}

\begin{theorem}[First local crystal-equation correction theorem]
\label{thm:first-local-crystal-equation-correction-theorem}
Assume:
\begin{enumerate}[leftmargin=2em,label=(\roman*)]
  \item the first local Theta-drift witness theorem applies;
  \item the first local phase-side repair theorem applies;
  \item the first local Teichmüller-persistence repair theorem applies.
\end{enumerate}
Then the present line supports a first local correction of the crystal equation.
More precisely:
\begin{enumerate}[leftmargin=2em,label=(\alph*)]
  \item the original crystal compatibility equation remains valid as the carrier/readout baseline;
  \item the active local persistence problem is repaired not by changing the carrier class, but by replacing $\mathcal C_{\mathrm{cell}}=0$ with the corrected equation $\mathcal C_{\mathrm{cell}}^{\sharp}=0$ on the active two-step window;
  \item the first crystal-equation correction is therefore a local persistence correction, not a new bulk crystal law.
\end{enumerate}
\end{theorem}

\begin{proof}
By the local obstruction analysis, the first failure is phase/geometry-prior rather than carrier-prior. The phase-side repair theorem supplies the first burden-side tightening, and the Teichmüller-persistence repair theorem supplies the remaining geometric tightening. Therefore the crystal equation is not replaced at bulk level; it is corrected locally by a persistence reinforcement on the active two-step window.
\end{proof}

\begin{corollary}[Current crystal-equation working rule]
\label{cor:current-crystal-equation-working-rule}
The current post-v3.32.0 line should not yet search first for a wholly new crystal equation. It should first search for the minimal local persistence correction that upgrades
\[
  \mathcal C_{\mathrm{cell}}=0
\]
to
\[
  \mathcal C_{\mathrm{cell}}^{\sharp}=0
\]
on the active closure-phase two-step window.
\end{corollary}

\begin{proof}
Immediate from Theorem~\ref{thm:first-local-crystal-equation-correction-theorem}.
\end{proof}


\begin{definition}[Local compact reinforcement pair]
\label{def:local-compact-reinforcement-pair}
The \emph{local compact reinforcement pair} extracted from the crystal persistence correction datum
\[
  \mathfrak K_{\mathrm{cell}}^{\sharp}(n)
  =
  (\eta_{\Theta}(n),\eta_T(n),\zeta_T(n),\lambda_T(n))
\]
is the pair
\[
  \mathfrak A_{\mathrm{loc}}^{\sharp}(n)
  :=
  (A_{\mathrm{bur}}(n),A_{\mathrm{geo}}(n))
\]
defined by
\[
  A_{\mathrm{bur}}(n):=\eta_{\Theta}(n),
  \qquad
  A_{\mathrm{geo}}(n):=\max\{\eta_T(n),\zeta_T(n),\lambda_T(n)\}.
\]
\end{definition}

\begin{remark}[Meaning of the compact pair]
\label{rem:meaning-of-the-compact-pair}
The first parameter records the burden-side phase reinforcement required to stop the second-step Theta drift.
The second parameter records the total geometric persistence reinforcement required on the same two-step window.
Thus the original four repair numbers are compressed into one burden-side and one geometric reinforcement scale.
\end{remark}

\begin{definition}[One-parameter shadow reinforcement]
\label{def:one-parameter-shadow-reinforcement}
The \emph{one-parameter shadow reinforcement} of the local compact reinforcement pair is
\[
  A_{\mathrm{tot}}(n)
  :=
  \max\{A_{\mathrm{bur}}(n),A_{\mathrm{geo}}(n)\}.
\]
\end{definition}

\begin{remark}[Why the one-parameter shadow is secondary]
\label{rem:why-the-one-parameter-shadow-is-secondary}
The scalar shadow $A_{\mathrm{tot}}(n)$ is useful only as a coarse local size readout.
It should not replace the two-parameter pair, because the current post-v3.32.0 line still distinguishes between burden-side repair and geometric persistence repair.
\end{remark}

\begin{definition}[Burden-dominant compact reinforcement]
\label{def:burden-dominant-compact-reinforcement}
The local compact reinforcement pair is called \emph{burden-dominant} if
\[
  A_{\mathrm{bur}}(n)\ge A_{\mathrm{geo}}(n).
\]
It is called \emph{geometrically dominant} if
\[
  A_{\mathrm{geo}}(n)>A_{\mathrm{bur}}(n).
\]
\end{definition}

\begin{proposition}[The crystal persistence correction datum admits a conservative two-parameter compression]
\label{prop:crystal-persistence-correction-datum-admits-two-parameter-compression}
Every local crystal persistence correction datum
\[
  (\eta_{\Theta}(n),\eta_T(n),\zeta_T(n),\lambda_T(n))
\]
admits a conservative two-parameter compression into
\[
  (A_{\mathrm{bur}}(n),A_{\mathrm{geo}}(n)).
\]
Moreover, the compressed pair dominates the original repair data in the sense that:
\begin{enumerate}[leftmargin=2em,label=(\alph*)]
  \item $A_{\mathrm{bur}}(n)$ retains the full burden-side strengthening requirement;
  \item $A_{\mathrm{geo}}(n)$ bounds every geometric persistence strengthening requirement from above.
\end{enumerate}
\end{proposition}

\begin{proof}
The first claim is immediate from Definition~\ref{def:local-compact-reinforcement-pair}.
The second follows because $A_{\mathrm{geo}}(n)$ is defined as the maximum of the three geometric persistence reinforcement terms.
Hence the compressed pair conservatively dominates the original local correction datum.
\end{proof}

\begin{proposition}[Current local reinforcement is not purely geometric]
\label{prop:current-local-reinforcement-is-not-purely-geometric}
Assume the first local repair-orientation theorem and the first local Teichmüller-persistence repair theorem.
Then the current local compact reinforcement pair is not geometrically dominant by default.
More precisely, the burden-side component
\[
  A_{\mathrm{bur}}(n)=\eta_{\Theta}(n)
\]
remains genuinely active in every admissible local repair reading.
\end{proposition}

\begin{proof}
By the local repair-orientation theorem, the first repair is Theta-dominant rather than purely geometric.
By the Teichmüller-persistence repair theorem, the geometric side is genuinely active but enters only after the burden-side drift has already been identified as the first required strengthening.
Therefore the current local compact reinforcement is not purely geometric by default.
\end{proof}

\begin{theorem}[First compact reinforcement theorem]
\label{thm:first-compact-reinforcement-theorem}
Assume:
\begin{enumerate}[leftmargin=2em,label=(\alph*)]
  \item the first local crystal-equation correction theorem applies;
  \item the first local repair-to-trigger calibration theorem applies;
  \item the first local Teichmüller-persistence repair theorem applies.
\end{enumerate}
Then the current post-v3.32.0 local crystal persistence problem admits a first compact reinforcement form.

More precisely:
\begin{enumerate}[leftmargin=2em,label=(\arabic*)]
  \item the four local repair parameters
  \[
    (\eta_{\Theta}(n),\eta_T(n),\zeta_T(n),\lambda_T(n))
  \]
  compress conservatively to the pair
  \[
    (A_{\mathrm{bur}}(n),A_{\mathrm{geo}}(n));
  \]
  \item the burden-side component $A_{\mathrm{bur}}(n)$ remains the first active repair scale;
  \item the geometric component $A_{\mathrm{geo}}(n)$ records the total persistence reinforcement required to keep the corrected crystal equation valid on the active two-step window;
  \item the scalar shadow $A_{\mathrm{tot}}(n)$ is only a coarse readout and should not replace the compact pair at the current stage.
\end{enumerate}
\end{theorem}

\begin{proof}
By Proposition~\ref{prop:crystal-persistence-correction-datum-admits-two-parameter-compression}, the four repair parameters admit a conservative compression into the compact pair.
By Proposition~\ref{prop:current-local-reinforcement-is-not-purely-geometric}, the burden-side component remains genuinely active and therefore cannot be discarded.
The geometric component dominates the three persistence-side repair contributions by definition, so it records the total persistence reinforcement requirement on the active two-step window.
The scalar shadow is only the maximum of the two compact parameters and therefore loses the burden/geometry distinction still needed at the current stage.
\end{proof}

\begin{corollary}[Current compact reinforcement working rule]
\label{cor:current-compact-reinforcement-working-rule}
The next local audit should not first track the full four-parameter repair datum.
It should first track the compact reinforcement pair
\[
  (A_{\mathrm{bur}}(n),A_{\mathrm{geo}}(n)),
\]
using the scalar shadow $A_{\mathrm{tot}}(n)$ only as a coarse size indicator.
\end{corollary}

\begin{proof}
Immediate from Theorem~\ref{thm:first-compact-reinforcement-theorem}.
\end{proof}


\begin{definition}[Local reinforcement ratio]
\label{def:local-reinforcement-ratio}
The \emph{local reinforcement ratio} is the quantity
\[
  \chi_{\mathrm{loc}}(n)
  :=
  \frac{A_{\mathrm{geo}}(n)}{A_{\mathrm{bur}}(n)}
\]
whenever $A_{\mathrm{bur}}(n)>0$.
\end{definition}

\begin{remark}[Meaning of the local reinforcement ratio]
\label{rem:meaning-of-the-local-reinforcement-ratio}
The ratio $\chi_{\mathrm{loc}}(n)$ measures how much geometric persistence reinforcement is required relative to the burden-side phase reinforcement on the active two-step window.
It is therefore the first compact indicator of whether the local repair is still burden-led or already close to a balanced coupled regime.
\end{remark}

\begin{definition}[Local reinforcement regime classes]
\label{def:local-reinforcement-regime-classes}
Assume $A_{\mathrm{bur}}(n)>0$.
The local compact reinforcement pair is said to be:
\begin{enumerate}[leftmargin=2em,label=(\alph*)]
  \item \emph{strictly burden-dominant} if
  \[
    \chi_{\mathrm{loc}}(n)<1;
  \]
  \item \emph{balanced-coupled} if
  \[
    \chi_{\mathrm{loc}}(n)=1;
  \]
  \item \emph{geometrically overcarried} if
  \[
    \chi_{\mathrm{loc}}(n)>1.
  \]
\end{enumerate}
\end{definition}

\begin{remark}[Why the ratio matters]
\label{rem:why-the-ratio-matters}
The compact pair already tells us that both burden-side and geometric reinforcement are active.
The ratio tells us which of the two is currently driving the local repair budget.
\end{remark}

\begin{proposition}[Current compact repair is not balanced by default]
\label{prop:current-compact-repair-is-not-balanced-by-default}
Assume the first compact reinforcement theorem and the first local repair-orientation theorem.
Then the current local compact repair is not balanced-coupled by default.
More precisely, the post-v3.32.0 line should be read first as strictly burden-dominant unless a later local audit shows that the geometric reinforcement has risen to the same scale.
\end{proposition}

\begin{proof}
By the first compact reinforcement theorem, the burden-side component $A_{\mathrm{bur}}(n)$ remains the first active repair scale, while the geometric component $A_{\mathrm{geo}}(n)$ records the total persistence reinforcement still required on the same window.
By the local repair-orientation theorem, the first repair is Theta-dominant rather than geometric-dominant.
Hence the default reading is not balanced-coupled, and the local repair is first to be read as burden-dominant unless later evidence forces a different regime class.
\end{proof}

\begin{definition}[Local reinforcement balance defect]
\label{def:local-reinforcement-balance-defect}
The \emph{local reinforcement balance defect} is
\[
  B_{\mathrm{loc}}(n)
  :=
  A_{\mathrm{geo}}(n)-A_{\mathrm{bur}}(n).
\]
\end{definition}

\begin{remark}[Interpretation of the balance defect]
\label{rem:interpretation-of-the-balance-defect}
Negative balance defect means burden-side dominance.
Zero balance defect means balanced coupled repair.
Positive balance defect means that the geometric persistence side has overtaken the burden-side repair scale.
\end{remark}

\begin{theorem}[First local reinforcement-ratio theorem]
\label{thm:first-local-reinforcement-ratio-theorem}
Assume:
\begin{enumerate}[leftmargin=2em,label=(\alph*)]
  \item the first compact reinforcement theorem applies;
  \item the first local Theta-drift witness theorem applies;
  \item the first local Teichmüller-persistence repair theorem applies.
\end{enumerate}
Then the current post-v3.32.0 local repair supports a first compact ratio reading.
More precisely:
\begin{enumerate}[leftmargin=2em,label=(\arabic*)]
  \item the compact pair admits the ratio $\chi_{\mathrm{loc}}(n)$ and the balance defect $B_{\mathrm{loc}}(n)$;
  \item the default local repair reading is strictly burden-dominant, hence
  \[
    B_{\mathrm{loc}}(n)<0
    \quad\text{and therefore}\quad
    \chi_{\mathrm{loc}}(n)<1
  \]
  in the default regime reading;
  \item the geometric reinforcement remains genuinely active even in this burden-dominant regime, so the ratio does not collapse to zero.
\end{enumerate}
\end{theorem}

\begin{proof}
Items (1) and the sign interpretation are immediate from the definitions of the ratio and balance defect.
By the first compact reinforcement theorem, the burden-side component remains the first active repair scale.
By the first local Teichmüller-persistence repair theorem, the geometric reinforcement is still genuinely active and therefore cannot be discarded.
Hence the default regime is burden-dominant but not purely one-sided, which yields $B_{\mathrm{loc}}(n)<0$ and therefore $\chi_{\mathrm{loc}}(n)<1$ in the default reading.
\end{proof}

\begin{corollary}[Current compact-ratio working rule]
\label{cor:current-compact-ratio-working-rule}
The next local audit should not only record the pair
\[
  (A_{\mathrm{bur}}(n),A_{\mathrm{geo}}(n)),
\]
but also track the ratio $\chi_{\mathrm{loc}}(n)$ and balance defect $B_{\mathrm{loc}}(n)$ in order to measure whether the current repair remains burden-led or is moving toward balanced coupled repair.
\end{corollary}

\begin{proof}
Immediate from Theorem~\ref{thm:first-local-reinforcement-ratio-theorem}.
\end{proof}


\begin{definition}[Local asymptotic repair orientation]
\label{def:local-asymptotic-repair-orientation}
The \emph{local asymptotic repair orientation} is the long-step tendency of the compact reinforcement ratio
\[
  \chi_{\mathrm{loc}}(n)
  =
  \frac{A_{\mathrm{geo}}(n)}{A_{\mathrm{bur}}(n)}
\]
along an admissible sequence of repaired local windows.
\end{definition}

\begin{definition}[Balance-seeking local repair family]
\label{def:balance-seeking-local-repair-family}
An admissible local repair family is called \emph{balance-seeking} if
\[
  \chi_{\mathrm{loc}}(n)\to 1
\]
along the family, equivalently if
\[
  B_{\mathrm{loc}}(n)\to 0.
\]
\end{definition}

\begin{definition}[Persistently burden-dominant local repair family]
\label{def:persistently-burden-dominant-local-repair-family}
An admissible local repair family is called \emph{persistently burden-dominant} if there exists some \(\delta_{\mathrm{bal}}>0\) such that
\[
  \chi_{\mathrm{loc}}(n)\le 1-\delta_{\mathrm{bal}}
\]
for all sufficiently large local repair steps.
\end{definition}

\begin{remark}[Why the asymptotic question matters]
\label{rem:why-the-asymptotic-question-matters}
The current compact repair is already burden-dominant at first order.
The real next question is whether this is only an initial orientation or whether the burden-side dominance remains structurally built into the repaired persistence problem.
\end{remark}

\begin{proposition}[The current line does not yet force balance-seeking behavior]
\label{prop:current-line-does-not-yet-force-balance-seeking-behavior}
Assume the current post-v3.32.0 local repair reading.
Then the existing local theorems do not yet force
\[
  \chi_{\mathrm{loc}}(n)\to 1.
\]
In particular, no current result requires the geometric reinforcement to rise all the way to burden-side balance.
\end{proposition}

\begin{proof}
The current local repair-orientation theorem yields burden-side dominance at first order, while the Teichmüller-side repair theorems only assert that the geometric reinforcement remains genuinely active. They do not assert that it asymptotically reaches the same scale as the burden-side reinforcement. Hence balance-seeking behavior is not yet forced.
\end{proof}

\begin{proposition}[The current line is naturally read as burden-dominant unless later balancing evidence appears]
\label{prop:current-line-naturally-read-as-burden-dominant-unless-later-balancing-evidence-appears}
Assume the first local reinforcement-ratio theorem and the first local repair-to-trigger calibration theorem.
Then the default asymptotic reading of the current line is persistently burden-dominant unless a later local audit exhibits a genuine balancing mechanism reducing \(B_{\mathrm{loc}}(n)\) toward zero.
\end{proposition}

\begin{proof}
The current ratio theorem places the local repair in the burden-dominant regime, and the first local repair-to-trigger calibration theorem still assigns the primary tightening to the Theta-side threshold. Thus the available repair logic continues to privilege the burden-side channel. Without an additional balancing mechanism, the natural default reading is that burden-side dominance persists.
\end{proof}

\begin{theorem}[First local asymptotic orientation theorem]
\label{thm:first-local-asymptotic-orientation-theorem}
Assume:
\begin{enumerate}[leftmargin=2em,label=(\alph*)]
  \item the first local reinforcement-ratio theorem applies;
  \item the first local repair-to-trigger calibration theorem applies;
  \item no later local theorem has yet exhibited a balancing mechanism forcing \(B_{\mathrm{loc}}(n)\to 0\).
\end{enumerate}
Then the current post-v3.32.0 line should be read asymptotically as burden-dominant rather than balance-seeking.
More precisely, the present local evidence supports persistent burden-side dominance as the default asymptotic orientation, while balanced-coupled repair remains only a later possibility to be justified by new local evidence.
\end{theorem}

\begin{proof}
By the first local reinforcement-ratio theorem, the current compact repair is burden-dominant at first order. By the first local repair-to-trigger calibration theorem, the primary tightening remains Theta-side, although the geometric repair remains genuinely active. In the absence of an additional balancing mechanism driving the balance defect toward zero, the natural asymptotic reading continues the burden-dominant orientation rather than forcing convergence to balance.
\end{proof}

\begin{corollary}[Current asymptotic working rule]
\label{cor:current-asymptotic-working-rule}
The next local audit should not presuppose
\[
  \chi_{\mathrm{loc}}(n)\to 1.
\]
Instead, it should first test whether any genuine balancing mechanism exists at all; until then, the current local repair should be treated as asymptotically burden-dominant with genuine but subordinate geometric carry.
\end{corollary}

\begin{proof}
Immediate from Theorem~\ref{thm:first-local-asymptotic-orientation-theorem}.
\end{proof}

\begin{definition}[Candidate local balancing mechanisms]
\label{def:candidate-local-balancing-mechanisms}
A \emph{candidate local balancing mechanism} for the compact reinforcement pair
\[
  \mathfrak A_{\mathrm{loc}}^{\sharp}(n)=(A_{\mathrm{bur}}(n),A_{\mathrm{geo}}(n))
\]
is any local structural effect capable of reducing the balance defect
\[
  B_{\mathrm{loc}}(n)=A_{\mathrm{geo}}(n)-A_{\mathrm{bur}}(n)
\]
toward zero. Three candidate types are distinguished:
\begin{enumerate}[leftmargin=2em,label=(\alph*)]
  \item a \emph{crystal-side balancing mechanism}, acting through the corrected crystal compatibility equation;
  \item a \emph{Teichm\"uller-side balancing mechanism}, acting through the geometric persistence layer alone;
  \item a \emph{hybrid crystal--Teichm\"uller balancing mechanism}, acting by coupled improvement of the crystal compatibility side and the Teichm\"uller persistence side.
\end{enumerate}
\end{definition}

\begin{remark}[Why a pure crystal correction is probably not enough]
\label{rem:why-pure-crystal-correction-probably-not-enough}
The current local line already required a first burden-side repair followed by an explicit Teichm\"uller-side repair. This suggests that crystal-side correction alone is unlikely to be the full balancing mechanism.
\end{remark}

\begin{remark}[Why a pure Teichm\"uller correction is also probably not enough]
\label{rem:why-pure-teichmueller-correction-probably-not-enough}
The local repair-orientation theorem still forces positive burden-side reinforcement
\[
  \eta_{\Theta}(n)>0.
\]
Hence the geometric side alone does not erase the burden-side asymmetry.
\end{remark}

\begin{definition}[Hybrid balancing witness]
\label{def:hybrid-balancing-witness}
A \emph{hybrid balancing witness} is a local witness package showing simultaneously:
\begin{enumerate}[leftmargin=2em,label=(\alph*)]
  \item improvement of the corrected crystal compatibility equation,
  \item improvement of the Teichm\"uller-side persistence inequalities,
  \item strict reduction of the balance defect
  \[
    B_{\mathrm{loc}}(n+1)>B_{\mathrm{loc}}(n)
  \]
  toward zero, equivalently
  \[
    |B_{\mathrm{loc}}(n+1)|<|B_{\mathrm{loc}}(n)|
  \]
  on the burden-dominant side.
\end{enumerate}
\end{definition}

\begin{proposition}[The current line does not yet support a purely crystal balancing mechanism]
\label{prop:current-line-does-not-yet-support-purely-crystal-balancing}
Assume the current post-v3.32.0 local repair line. Then the available local evidence does not yet support reading the first balancing mechanism as purely crystal-side.
\end{proposition}

\begin{proof}
The local crystal-equation correction theorem already required the prior burden-side and Teichm\"uller-side repair data. Hence the crystal correction enters as part of a repaired persistence package, not as a self-sufficient balancing source on its own.
\end{proof}

\begin{proposition}[The current line does not yet support a purely Teichm\"uller balancing mechanism]
\label{prop:current-line-does-not-yet-support-purely-teichmueller-balancing}
Assume the first local repair-orientation theorem. Then the available local evidence does not yet support reading the first balancing mechanism as purely Teichm\"uller-side.
\end{proposition}

\begin{proof}
The first local repair-orientation theorem keeps the dominant local repair on the burden side while the geometric side remains genuinely active but subordinate. Hence the currently visible repair is not purely Teichm\"uller-driven.
\end{proof}

\begin{theorem}[First local balancing-mechanism orientation theorem]
\label{thm:first-local-balancing-mechanism-orientation-theorem}
Assume:
\begin{enumerate}[leftmargin=2em,label=(\alph*)]
  \item the first local asymptotic orientation theorem applies;
  \item the first local crystal-equation correction theorem applies;
  \item the first local Teichm\"uller-persistence repair theorem applies.
\end{enumerate}
Then the most plausible candidate for a future local balancing mechanism is not purely crystal-side and not purely Teichm\"uller-side, but hybrid crystal--Teichm\"uller.
\end{theorem}

\begin{proof}
By the first local asymptotic orientation theorem, the current line remains burden-dominant unless a later balancing mechanism appears. By the local crystal-equation correction theorem, crystal-side correction is relevant but not self-sufficient. By the local Teichm\"uller-persistence repair theorem, geometric repair is genuinely active but likewise not self-sufficient. Therefore the first plausible balancing mechanism must couple both corrections if it is to reduce the balance defect toward zero.
\end{proof}

\begin{corollary}[Current balancing-mechanism working rule]
\label{cor:current-balancing-mechanism-working-rule}
The next local audit should not first search for a purely crystal balancing mechanism or a purely Teichm\"uller balancing mechanism. Instead, it should search for the first hybrid witness that lowers the balance defect by coordinated crystal-side and geometric-side repair.
\end{corollary}

\begin{proof}
Immediate from Theorem~\ref{thm:first-local-balancing-mechanism-orientation-theorem}.
\end{proof}


\begin{definition}[First hybrid balancing step]
\label{def:first-hybrid-balancing-step}
A \emph{first hybrid balancing step} from the local state at step \(n\) to the next step \(n+1\) is a local repaired transition satisfying all of the following:
\begin{enumerate}[leftmargin=2em,label=(\alph*)]
  \item the corrected crystal compatibility equation improves from
  \[
    \mathcal C_{\mathrm{cell}}^{\sharp}(X_n,X_{n+1})
  \]
  to a closer-to-compatible value on the updated two-step window;
  \item the Teichm\"uller-side persistence repair data improve in the sense that the repaired geometric persistence test is strictly better satisfied on the updated window;
  \item the compact reinforcement pair
  \[
    \mathfrak A_{\mathrm{loc}}^{\sharp}(n)=(A_{\mathrm{bur}}(n),A_{\mathrm{geo}}(n))
  \]
  is updated to
  \[
    \mathfrak A_{\mathrm{loc}}^{\sharp}(n+1)=(A_{\mathrm{bur}}(n+1),A_{\mathrm{geo}}(n+1))
  \]
  with strictly smaller balance defect magnitude
  \[
    |B_{\mathrm{loc}}(n+1)|<|B_{\mathrm{loc}}(n)|;
  \]
  \item the update remains inside the same local admissibility regime.
\end{enumerate}
\end{definition}

\begin{definition}[First hybrid balancing witness package]
\label{def:first-hybrid-balancing-witness-package}
A \emph{first hybrid balancing witness package} at step \(n\) is a finite witness datum
\[
  \mathfrak W_{\mathrm{hyb}}(n)
  =
  \bigl(
    \delta_{\mathrm{cell}}(n),
    \delta_{T}(n),
    \Delta B_{\mathrm{loc}}(n)
  \bigr)
\]
such that:
\begin{enumerate}[leftmargin=2em,label=(\alph*)]
  \item \(\delta_{\mathrm{cell}}(n)>0\) measures the crystal-side compatibility gain of the first hybrid balancing step;
  \item \(\delta_{T}(n)>0\) measures the Teichm\"uller-side persistence gain of the same step;
  \item
  \[
    \Delta B_{\mathrm{loc}}(n)
    :=
    |B_{\mathrm{loc}}(n)|-|B_{\mathrm{loc}}(n+1)|
    >0
  \]
  measures the strict gain toward local balance.
\end{enumerate}
\end{definition}

\begin{remark}[Meaning of the hybrid witness package]
\label{rem:meaning-of-the-hybrid-witness-package}
The first hybrid balancing witness package is the first local datum that does more than merely diagnose burden-dominance.
It certifies that crystal-side correction and Teichm\"uller-side repair improve together and that the resulting coupled improvement genuinely lowers the balance defect.
\end{remark}

\begin{definition}[Hybrid balancing efficiency ratio]
\label{def:hybrid-balancing-efficiency-ratio}
The \emph{hybrid balancing efficiency ratio} of a first hybrid balancing witness package is
\[
  \Xi_{\mathrm{hyb}}(n)
  :=
  \frac{\Delta B_{\mathrm{loc}}(n)}{\delta_{\mathrm{cell}}(n)+\delta_{T}(n)},
\]
provided the denominator is positive.
\end{definition}

\begin{remark}[Why the efficiency ratio is useful]
\label{rem:why-the-efficiency-ratio-is-useful}
This ratio records how much balance-defect reduction is obtained per unit of coupled crystal-plus-Teichm\"uller repair.
It is therefore the first local scalar witness that the hybrid mechanism is not merely present, but actually effective.
\end{remark}

\begin{proposition}[A first hybrid balancing step is stronger than separate side repairs]
\label{prop:first-hybrid-balancing-step-stronger-than-separate-side-repairs}
Every first hybrid balancing step is stronger than a purely crystal-side repair and stronger than a purely Teichm\"uller-side repair considered separately.
\end{proposition}

\begin{proof}
By definition, a first hybrid balancing step requires simultaneous crystal-side improvement, simultaneous Teichm\"uller-side persistence improvement, and strict reduction of the balance defect magnitude.
A purely one-sided repair can satisfy at most one of the two side-improvement clauses, but not both together with the coupled balance reduction requirement.
Hence the hybrid step is strictly stronger.
\end{proof}

\begin{lemma}[First hybrid balancing witness lemma]
\label{lem:first-hybrid-balancing-witness-lemma}
Assume:
\begin{enumerate}[leftmargin=2em,label=(\alph*)]
  \item the first local balancing-mechanism orientation theorem applies;
  \item a first hybrid balancing step exists on the current local window.
\end{enumerate}
Then a first hybrid balancing witness package exists.
\end{lemma}

\begin{proof}
If a first hybrid balancing step exists, then by definition it produces strictly positive crystal-side compatibility gain, strictly positive Teichm\"uller-side persistence gain, and strict decrease of the balance defect magnitude.
These three gains form the witness package
\[
  \mathfrak W_{\mathrm{hyb}}(n)
  =
  (\delta_{\mathrm{cell}}(n),\delta_{T}(n),\Delta B_{\mathrm{loc}}(n)).
\]
\end{proof}

\begin{theorem}[First hybrid balancing witness theorem]
\label{thm:first-hybrid-balancing-witness-theorem}
Assume:
\begin{enumerate}[leftmargin=2em,label=(\alph*)]
  \item the first local balancing-mechanism orientation theorem applies;
  \item the current line is still asymptotically burden-dominant;
  \item a local repaired two-step window exists on which crystal-side and Teichm\"uller-side improvements occur simultaneously.
\end{enumerate}
Then the current line supports a first hybrid balancing witness exactly when the simultaneous side improvements produce strict reduction of
\[
  |B_{\mathrm{loc}}(n)|.
\]
Equivalently, the first honest evidence for a genuine balancing mechanism is not the mere coexistence of crystal and geometric repair, but the existence of a positive witness package
\[
  \mathfrak W_{\mathrm{hyb}}(n)
  =
  (\delta_{\mathrm{cell}}(n),\delta_{T}(n),\Delta B_{\mathrm{loc}}(n))
\]
with
\[
  \Delta B_{\mathrm{loc}}(n)>0.
\]
\end{theorem}

\begin{proof}
By the balancing-mechanism orientation theorem, any plausible future balancing mechanism must be hybrid rather than purely one-sided.
If simultaneous crystal-side and Teichm\"uller-side improvements occur but do not reduce \(|B_{\mathrm{loc}}(n)|\), then no genuine balancing has yet been witnessed; the system remains merely burden-dominant with parallel repairs.
Conversely, if the simultaneous side improvements do reduce \(|B_{\mathrm{loc}}(n)|\), then the resulting gains define a positive hybrid witness package, and the first honest evidence for a balancing mechanism is present.
\end{proof}

\begin{corollary}[Current hybrid-balancing working rule]
\label{cor:current-hybrid-balancing-working-rule}
The next local audit should not merely ask whether crystal-side and Teichm\"uller-side repair are both visible.
It should ask whether they jointly produce
\[
  \Delta B_{\mathrm{loc}}(n)>0,
\]
and hence a first genuine hybrid balancing witness package.
\end{corollary}


\begin{definition}[Hybrid witness existence obstruction]
\label{def:hybrid-witness-existence-obstruction}
A \emph{hybrid witness existence obstruction} at step \(n\) is the failure of the current local line to produce any witness package
\[
  \mathfrak W_{\mathrm{hyb}}(n)
  =
  (\delta_{\mathrm{cell}}(n),\delta_T(n),\Delta B_{\mathrm{loc}}(n))
\]
with
\[
  \delta_{\mathrm{cell}}(n)>0,
  \qquad
  \delta_T(n)>0,
  \qquad
  \Delta B_{\mathrm{loc}}(n)>0.
\]
Equivalently, the line remains burden-dominant because no first honest hybrid balancing witness is yet produced.
\end{definition}

\begin{definition}[Weak hybrid efficiency obstruction]
\label{def:weak-hybrid-efficiency-obstruction}
A \emph{weak hybrid efficiency obstruction} at step \(n\) is the situation in which a first hybrid balancing witness package exists but its hybrid balancing efficiency ratio
\[
  \Xi_{\mathrm{hyb}}(n)
  =
  \frac{\Delta B_{\mathrm{loc}}(n)}{\delta_{\mathrm{cell}}(n)+\delta_T(n)}
\]
remains too small to alter the current asymptotic burden-dominant orientation.
\end{definition}

\begin{remark}[Existence versus efficiency]
\label{rem:existence-versus-efficiency}
There are therefore two distinct post-r30 obstructions:
first, the line may fail to produce any honest hybrid witness at all;
second, it may produce such a witness, but only with negligible balancing efficiency.
These two failure modes should not be confused.
\end{remark}

\begin{definition}[Existence-prior local regime]
\label{def:existence-prior-local-regime}
The current local line is said to be in an \emph{existence-prior local regime} if the first balancing obstruction is of witness-existence type rather than weak-efficiency type.
\end{definition}

\begin{definition}[Efficiency-prior local regime]
\label{def:efficiency-prior-local-regime}
The current local line is said to be in an \emph{efficiency-prior local regime} if a first hybrid balancing witness is already present, but its efficiency ratio remains too small to change the burden-dominant asymptotic orientation.
\end{definition}

\begin{proposition}[Asymptotic burden-dominance blocks efficiency-first readings]
\label{prop:asymptotic-burden-dominance-blocks-efficiency-first-readings}
Assume the first local asymptotic orientation theorem and the first hybrid balancing witness theorem.
If no positive witness package
\[
  \mathfrak W_{\mathrm{hyb}}(n)
\]
with
\[
  \Delta B_{\mathrm{loc}}(n)>0
\]
is yet available, then the current line cannot honestly be read as efficiency-prior.
\end{proposition}

\begin{proof}
An efficiency-prior reading presupposes that a first hybrid balancing witness already exists and that only its efficiency remains insufficient. If no positive witness package is available at all, then the obstruction occurs before efficiency becomes a meaningful variable. Hence the line cannot honestly be read as efficiency-prior.
\end{proof}

\begin{theorem}[First hybrid obstruction orientation theorem]
\label{thm:first-hybrid-obstruction-orientation-theorem}
Assume:
\begin{enumerate}[leftmargin=2em,label=(\alph*)]
  \item the first local asymptotic orientation theorem applies;
  \item the first local balancing-mechanism orientation theorem applies;
  \item the first hybrid balancing witness theorem identifies positive hybrid witness data as the first honest balancing evidence.
\end{enumerate}
Then the current post-r30 line is obstruction-prior on the witness-existence side.
More precisely, its first balancing obstruction is not that an existing hybrid witness is too weak, but that no first positive hybrid balancing witness package has yet been produced.
\end{theorem}

\begin{proof}
By the first local asymptotic orientation theorem, the line remains burden-dominant by default unless a genuine balancing mechanism is exhibited. By the balancing-mechanism orientation theorem, any such genuine mechanism must be hybrid. By the first hybrid balancing witness theorem, the first honest evidence for that mechanism is exactly a positive witness package with \(\Delta B_{\mathrm{loc}}(n)>0\). Therefore, as long as no such package is present, the primary obstruction lies on the witness-existence side rather than on the efficiency side.
\end{proof}

\begin{corollary}[Current hybrid-balancing obstruction rule]
\label{cor:current-hybrid-balancing-obstruction-rule}
The next local audit should first ask whether a positive hybrid witness package exists at all. Only after such a package is found should the secondary question of hybrid efficiency be treated as the leading obstruction.
\end{corollary}

\begin{proof}
Immediate from Theorem~\ref{thm:first-hybrid-obstruction-orientation-theorem}.
\end{proof}




\begin{definition}[Candidate positive hybrid construction step]
\label{def:candidate-positive-hybrid-construction-step}
A \emph{candidate positive hybrid construction step} at stage $n$ is a local update from step $n$ to $n+1$ such that:
\begin{enumerate}[leftmargin=2em,label=(\alph*)]
  \item the burden-side repair is active with
  \[
    \delta_{\mathrm{cell}}(n)>0;
  \]
  \item the Teichm\"uller-side persistence repair is active with
  \[
    \delta_T(n)>0;
  \]
  \item both repairs act on the same local closure-phase window;
  \item the updated local balance defect satisfies
  \[
    B_{\mathrm{loc}}(n+1)>B_{\mathrm{loc}}(n).
  \]
\end{enumerate}
\end{definition}

\begin{remark}[Why the construction is phrased this way]
\label{rem:why-construction-phrased-this-way}
The current line does not yet justify assuming that a positive hybrid witness appears automatically once both repair channels are nonzero.
The real construction requirement is stronger: the simultaneous crystal-side and Teichm\"uller-side repairs must improve the local balance defect itself.
\end{remark}

\begin{definition}[Positive hybrid construction datum]
\label{def:positive-hybrid-construction-datum}
A \emph{positive hybrid construction datum} at stage $n$ is a quadruple
\[
  \mathfrak C_{\mathrm{hyb}}^{+}(n)
  :=
  \bigl(
    \delta_{\mathrm{cell}}(n),
    \delta_T(n),
    B_{\mathrm{loc}}(n),
    B_{\mathrm{loc}}(n+1)
  \bigr)
\]
attached to a candidate positive hybrid construction step.
\end{definition}

\begin{definition}[Strictly positive hybrid construction]
\label{def:strictly-positive-hybrid-construction}
A candidate positive hybrid construction step is called \emph{strictly positive} if, in addition,
\[
  \Delta B_{\mathrm{loc}}(n)
  :=
  B_{\mathrm{loc}}(n+1)-B_{\mathrm{loc}}(n)
  >0.
\]
\end{definition}

\begin{definition}[Hybrid construction gate]
\label{def:hybrid-construction-gate}
The \emph{hybrid construction gate} at stage $n$ is the conjunction of the conditions
\[
  \delta_{\mathrm{cell}}(n)>0,
  \qquad
  \delta_T(n)>0,
  \qquad
  \Delta B_{\mathrm{loc}}(n)>0.
\]
\end{definition}

\begin{remark}[Existence obstruction as gate failure]
\label{rem:existence-obstruction-as-gate-failure}
Under the current obstruction-oriented reading, the missing positive hybrid witness is equivalently the failure of the hybrid construction gate.
This makes the existence problem testable in the same local language as the later efficiency problem.
\end{remark}

\begin{proposition}[Positive hybrid witness construction implies witness existence]
\label{prop:positive-hybrid-witness-construction-implies-witness-existence}
Every strictly positive hybrid construction step produces a positive hybrid witness package
\[
  \mathfrak W_{\mathrm{hyb}}(n)
  =
  \bigl(
    \delta_{\mathrm{cell}}(n),
    \delta_T(n),
    \Delta B_{\mathrm{loc}}(n)
  \bigr)
\]
with
\[
  \delta_{\mathrm{cell}}(n)>0,
  \qquad
  \delta_T(n)>0,
  \qquad
  \Delta B_{\mathrm{loc}}(n)>0.
\]
\end{proposition}

\begin{proof}
By definition, a strictly positive hybrid construction step already supplies positive crystal-side repair, positive Teichm\"uller-side repair, and a strictly positive balance-defect gain.
These are exactly the components required for a positive hybrid witness package.
\end{proof}

\begin{proposition}[The current existence obstruction is equivalent to absence of a strictly positive hybrid construction]
\label{prop:current-existence-obstruction-equivalent-to-absence-of-construction}
Assume the first hybrid obstruction orientation theorem.
Then the current line lies in the existence-prior regime if and only if no strictly positive hybrid construction step is yet available on the active local window.
\end{proposition}

\begin{proof}
By the first hybrid obstruction orientation theorem, the current primary obstruction is the absence of a positive hybrid witness package.
By Proposition~\ref{prop:positive-hybrid-witness-construction-implies-witness-existence}, every strictly positive hybrid construction step would already produce such a package.
Hence the existence-prior obstruction holds exactly when no such construction step is yet available.
\end{proof}

\begin{theorem}[First positive hybrid witness construction theorem]
\label{thm:first-positive-hybrid-witness-construction-theorem}
Assume:
\begin{enumerate}[leftmargin=2em,label=(\alph*)]
  \item the first hybrid obstruction orientation theorem applies;
  \item the local crystal-side correction datum is available;
  \item the local Teichm\"uller-side repair datum is available;
  \item both data act on the same active closure-phase window.
\end{enumerate}
Then the first honest transition out of the existence-prior regime occurs exactly when a strictly positive hybrid construction step is produced.
Equivalently, the first positive hybrid balancing witness arises exactly when the hybrid construction gate is opened on the active local window.
\end{theorem}

\begin{proof}
By the first hybrid obstruction orientation theorem, the current line remains stuck before the efficiency question precisely because no positive hybrid witness yet exists.
By Proposition~\ref{prop:positive-hybrid-witness-construction-implies-witness-existence}, every strictly positive hybrid construction step produces such a witness.
Conversely, on the active local window a first positive witness must come from simultaneous positive crystal-side and Teichm\"uller-side repair that genuinely improves the local balance defect, which is exactly the content of a strictly positive hybrid construction step.
Hence exiting the existence-prior regime is equivalent to producing such a step, or equivalently to opening the hybrid construction gate.
\end{proof}

\begin{corollary}[Current hybrid-construction working rule]
\label{cor:current-hybrid-construction-working-rule}
The next local audit should not yet optimize hybrid efficiency.
It should first test whether the active local crystal-side and Teichm\"uller-side repairs can be synchronized into a strictly positive hybrid construction step with
\[
  \Delta B_{\mathrm{loc}}(n)>0.
\]
\end{corollary}

\begin{proof}
Immediate from Theorem~\ref{thm:first-positive-hybrid-witness-construction-theorem}.
\end{proof}


\begin{definition}[Crystal-side gate obstruction]
\label{def:crystal-side-gate-obstruction}
A \emph{crystal-side gate obstruction} at stage $n$ is the failure of the hybrid construction gate caused primarily by the crystal-side component, i.e. by one of the following:
\begin{enumerate}[leftmargin=2em,label=(\alph*)]
  \item
  \[
    \delta_{\mathrm{cell}}(n)\le 0;
  \]
  \item the crystal-side correction datum is not available on the active local window;
  \item the crystal-side correction fails to synchronize with the active Teichm\"uller-side repair on the same closure-phase window.
\end{enumerate}
\end{definition}

\begin{definition}[Teichm\"uller-side gate obstruction]
\label{def:teichmueller-side-gate-obstruction}
A \emph{Teichm\"uller-side gate obstruction} at stage $n$ is the failure of the hybrid construction gate caused primarily by the geometric persistence side, i.e. by one of the following:
\begin{enumerate}[leftmargin=2em,label=(\alph*)]
  \item
  \[
    \delta_T(n)\le 0;
  \]
  \item the Teichm\"uller-side repair datum is not available on the active local window;
  \item the geometric repair remains too weak to support
  \[
    \Delta B_{\mathrm{loc}}(n)>0
  \]
  even when the crystal-side correction is already positive.
\end{enumerate}
\end{definition}

\begin{definition}[Gate-side desynchronization defect]
\label{def:gate-side-desynchronization-defect}
The \emph{gate-side desynchronization defect} at stage $n$ is the obstruction datum
\[
  \mathfrak D_{\mathrm{gate}}(n)
  :=
  \bigl(
    \delta_{\mathrm{cell}}(n),
    \delta_T(n),
    \Delta B_{\mathrm{loc}}(n)
  \bigr)
\]
read under the question which side first fails to keep the hybrid construction gate open.
\end{definition}

\begin{remark}[What ``first out of sync'' means]
\label{rem:what-first-out-of-sync-means}
The present question is not yet one of hybrid efficiency.
It asks which side fails first at the level of gate opening:
the crystal-side correction, the Teichm\"uller-side persistence repair, or the synchronization of both on the same local closure-phase window.
\end{remark}

\begin{proposition}[The current line is not crystal-obstruction-prior]
\label{prop:current-line-is-not-crystal-obstruction-prior}
Under the current post-v3.32.0 reading, the active local line is not primarily blocked by absence of crystal-side correction.
\end{proposition}

\begin{proof}
The local crystal-equation correction theorem and the compact reinforcement line already support a positive crystal-side correction datum on the active local window.
The remaining obstruction therefore does not first arise from the total absence of the crystal-side contribution.
\end{proof}

\begin{proposition}[The current line is Teichm\"uller-side-prior at the gate]
\label{prop:current-line-is-teichmueller-side-prior-at-the-gate}
Under the current post-v3.32.0 reading, the first gate-side desynchronization is read primarily on the Teichm\"uller-side persistence channel.
\end{proposition}

\begin{proof}
The tightened-package plausibility reading and the subsequent Teichm\"uller-persistence repair line already isolate the residual insufficiency on the geometric persistence side after the first burden-side repair has been introduced.
The hybrid-construction obstruction then remains because a positive crystal-side correction does not yet suffice to produce
\[
  \Delta B_{\mathrm{loc}}(n)>0.
\]
Hence the first gate-side desynchronization is read primarily on the Teichm\"uller-side channel rather than on the crystal side.
\end{proof}

\begin{theorem}[First gate-side desynchronization orientation theorem]
\label{thm:first-gate-side-desynchronization-orientation-theorem}
Assume:
\begin{enumerate}[leftmargin=2em,label=(\alph*)]
  \item the first positive hybrid witness construction theorem applies;
  \item the local crystal-equation correction theorem applies;
  \item the first local Teichm\"uller-persistence repair theorem applies;
  \item the current line still lies in the existence-prior hybrid obstruction regime.
\end{enumerate}
Then the current gate failure is Teichm\"uller-side-prior.

More precisely:
\begin{enumerate}[leftmargin=2em,label=(\roman*)]
  \item the crystal-side correction is not the first component to fall out of sync;
  \item the primary remaining defect is the insufficient geometric persistence contribution needed to turn the positive local corrections into
  \[
    \Delta B_{\mathrm{loc}}(n)>0;
  \]
  \item the next local construction step should therefore strengthen the Teichm\"uller-side synchronization before attempting any efficiency optimization.
\end{enumerate}
\end{theorem}

\begin{proof}
By Proposition~\ref{prop:current-line-is-not-crystal-obstruction-prior}, the crystal-side correction is already locally present and is not the primary missing ingredient.
By Proposition~\ref{prop:current-line-is-teichmueller-side-prior-at-the-gate}, the residual desynchronization sits on the Teichm\"uller-side persistence channel.
Since the hybrid witness still does not yet exist, the active local task is not efficiency optimization but restoration of sufficient geometric persistence to open the gate.
\end{proof}

\begin{corollary}[Current gate-side working rule]
\label{cor:current-gate-side-working-rule}
The next local audit should first ask how the Teichm\"uller-side repair must be strengthened or synchronized so that the already present crystal-side correction becomes hybrid-gate-effective.
Only after that should the line pass to hybrid efficiency questions.
\end{corollary}

\begin{proof}
Immediate from Theorem~\ref{thm:first-gate-side-desynchronization-orientation-theorem}.
\end{proof}




\begin{definition}[Phase-lock synchronization defect]
\label{def:phase-lock-synchronization-defect}
A \emph{phase-lock synchronization defect} at stage $n$ is the failure of the second-step local repair to keep the phase-side burden correction aligned with the same closure-phase window on which the first-step near-trigger state was read.
\end{definition}

\begin{definition}[Teichm\"uller-carry synchronization defect]
\label{def:teichmueller-carry-synchronization-defect}
A \emph{Teichm\"uller-carry synchronization defect} at stage $n$ is the failure of the geometric persistence correction to carry the repaired second step inside the same local Teichm\"uller-safe band as the first step.
\end{definition}

\begin{definition}[Local Teichm\"uller-side synchronization datum]
\label{def:local-teichmueller-side-synchronization-datum}
The \emph{local Teichm\"uller-side synchronization datum} at stage $n$ is
\[
  \mathfrak S_T^{\mathrm{loc}}(n)
  :=
  \bigl(
    \sigma_{\Theta}^{\mathrm{lock}}(n),
    \sigma_{T}^{\mathrm{carry}}(n)
  \bigr),
\]
where:
\begin{enumerate}[leftmargin=2em,label=(\alph*)]
  \item $\sigma_{\Theta}^{\mathrm{lock}}(n)$ measures the remaining phase-lock misalignment of the repaired second step;
  \item $\sigma_T^{\mathrm{carry}}(n)$ measures the remaining geometric carry defect on the same two-step local window.
\end{enumerate}
\end{definition}

\begin{remark}[What synchronization means locally]
\label{rem:what-synchronization-means-locally}
The present question is no longer whether geometry matters at all. It asks more precisely which kind of Teichm\"uller-side synchronization is still missing: the phase-side lock itself, the geometric carry of that lock, or both together.
\end{remark}

\begin{definition}[Locally synchronized repaired second step]
\label{def:locally-synchronized-repaired-second-step}
A repaired second step at stage $n$ is called \emph{locally synchronized} if
\[
  \sigma_{\Theta}^{\mathrm{lock}}(n)=0
  \qquad\text{and}\qquad
  \sigma_T^{\mathrm{carry}}(n)=0.
\]
\end{definition}

\begin{proposition}[The current line is not raw-geometry-prior]
\label{prop:current-line-is-not-raw-geometry-prior}
Under the current post-v3.32.0 reading, the active local line is not primarily blocked by raw geometric carry alone.
\end{proposition}

\begin{proof}
The current line already carries a positive Teichm\"uller-side repair contribution and a positive geometric witness contribution. The residual failure therefore cannot be read first as total absence of geometry, but only as insufficient synchronization of that geometric contribution with the repaired phase-side step.
\end{proof}

\begin{proposition}[The current line is phase-lock-prior on the Teichm\"uller side]
\label{prop:current-line-is-phase-lock-prior-on-the-teichmueller-side}
Under the current post-v3.32.0 reading, the first remaining Teichm\"uller-side synchronization defect is phase-lock-prior.
\end{proposition}

\begin{proof}
The gate-side desynchronization orientation theorem already identifies the first remaining gate defect as Teichm\"uller-side-prior rather than crystal-side-prior. The local repair-orientation and tightened-package plausibility lines then show that the geometric channel is active but still fails to hold the repaired second step in the same closure-phase window. Hence the first remaining defect is best read as insufficient phase-lock synchronization, with geometric carry acting as the witness and support channel.
\end{proof}

\begin{theorem}[First Teichm\"uller-side synchronization orientation theorem]
\label{thm:first-teichmueller-side-synchronization-orientation-theorem}
Assume:
\begin{enumerate}[leftmargin=2em,label=(\alph*)]
  \item the first local Theta-drift witness theorem applies;
  \item the first local phase-side repair theorem applies;
  \item the first local Teichm\"uller-persistence repair theorem applies;
  \item the first gate-side desynchronization orientation theorem applies.
\end{enumerate}
Then the first remaining Teichm\"uller-side synchronization problem is phase-lock-prior rather than raw-carry-prior.

More precisely:
\begin{enumerate}[leftmargin=2em,label=(\roman*)]
  \item the missing synchronization is not first the total absence of geometric carry;
  \item the missing synchronization is first the failure to keep the repaired second step phase-locked inside the same closure-phase window;
  \item the geometric carry defect remains genuinely active, but as the supporting witness of the phase-lock failure rather than as the primary independent obstruction.
\end{enumerate}
\end{theorem}

\begin{proof}
By Proposition~\ref{prop:current-line-is-not-raw-geometry-prior}, the active line is not blocked first by total absence of geometry. By Proposition~\ref{prop:current-line-is-phase-lock-prior-on-the-teichmueller-side}, the first remaining synchronization failure is read on the phase-lock component of the repaired second step. Therefore the current Teichm\"uller-side synchronization problem is phase-lock-prior, with the geometric carry defect remaining real but subordinate.
\end{proof}

\begin{corollary}[Current Teichm\"uller-side synchronization working rule]
\label{cor:current-teichmueller-side-synchronization-working-rule}
The next local audit should not ask first for more geometry in the abstract. It should ask first how the repaired second step can be phase-locked back into the same closure-phase window, while using the geometric channel to certify and carry that repaired lock.
\end{corollary}

\begin{proof}
Immediate from Theorem~\ref{thm:first-teichmueller-side-synchronization-orientation-theorem}.
\end{proof}



\begin{definition}[Local phase-lock repair datum]
\label{def:local-phase-lock-repair-datum}
The \emph{local phase-lock repair datum} at stage $n$ is
\[
  \mathfrak R_{\mathrm{lock}}^{\mathrm{loc}}(n)
  :=
  \bigl(
    \ell_{\Theta}^{\mathrm{lock}}(n),
    \ell_{T}^{\mathrm{lock}}(n)
  \bigr),
\]
where $\ell_{\Theta}^{\mathrm{lock}}(n)$ denotes the extra phase-side lock correction required to remove the residual phase-lock defect and $\ell_{T}^{\mathrm{lock}}(n)$ denotes the extra Teichm\"uller-side carry needed to stabilize that repaired lock on the same two-step window.
\end{definition}

\begin{definition}[Phase-lock repaired local window]
\label{def:phase-lock-repaired-local-window}
A repaired two-step local window at stage $n$ is called \emph{phase-lock repaired} if, after applying the local phase-lock repair datum, one has
\[
  \sigma_{\Theta}^{\mathrm{lock}}(n)=0
  \qquad\text{and}\qquad
  \sigma_T^{\mathrm{carry}}(n)\le \ell_T^{\mathrm{lock}}(n).
\]
Equivalently, the repaired second step is brought back into the same active closure-phase window while the remaining geometric carry defect is absorbed by the stated lock-carry allowance.
\end{definition}

\begin{definition}[Minimal local phase-lock repair]
\label{def:minimal-local-phase-lock-repair}
A local phase-lock repair datum
\[
  \mathfrak R_{\mathrm{lock}}^{\mathrm{loc}}(n)
  =
  \bigl(
    \ell_{\Theta}^{\mathrm{lock}}(n),
    \ell_{T}^{\mathrm{lock}}(n)
  \bigr)
\]
is called \emph{minimal} if no strictly smaller pair of nonnegative corrections removes the phase-lock synchronization defect on the same two-step local window.
\end{definition}

\begin{remark}[What the phase-lock repair really does]
\label{rem:what-the-phase-lock-repair-really-does}
The phase-lock repair is not a fresh bulk law. It is the smallest local correction that forces the repaired second step back onto the same closure-phase slice as the first near-trigger step. Geometrically, one may think of it as a local re-aiming plus carry-stabilization of the second point so that both points lie again in the same admissible narrow corridor rather than in neighboring but misaligned corridors.
\end{remark}

\begin{lemma}[First local phase-lock repair lemma]
\label{lem:first-local-phase-lock-repair-lemma}
Assume:
\begin{enumerate}[leftmargin=2em,label=(\alph*)]
  \item the first Teichm\"uller-side synchronization orientation theorem applies;
  \item the active local obstruction is phase-lock-prior;
  \item a local phase-lock repair datum exists.
\end{enumerate}
Then the vanishing condition
\[
  \sigma_{\Theta}^{\mathrm{lock}}(n)=0
\]
is the first honest local repair target for reopening the hybrid construction gate.
\end{lemma}

\begin{proof}
By Theorem~\ref{thm:first-teichmueller-side-synchronization-orientation-theorem}, the first remaining synchronization obstruction is read on the phase-lock component rather than on raw geometric absence. Hence the first honest repair target is to remove the phase-lock defect itself. The accompanying carry correction remains necessary, but it is subordinate to the lock restoration as long as the active obstruction is phase-lock-prior.
\end{proof}

\begin{proposition}[Current line supports phase-lock-prior local repair]
\label{prop:current-line-supports-phase-lock-prior-local-repair}
Under the present post-v3.32.0 reading, the next local strengthening step should first reduce $\sigma_{\Theta}^{\mathrm{lock}}(n)$ and only then optimize the associated geometric carry.
\end{proposition}

\begin{proof}
This is just the operational form of the phase-lock-prior diagnosis. Since the line is not raw-geometry-prior, one does not first optimize geometric amplitude in the abstract. One first restores alignment of the repaired second step with the active closure-phase window, and then uses the Teichm\"uller-side correction to support that restored lock.
\end{proof}

\begin{theorem}[First phase-lock repair theorem]
\label{thm:first-phase-lock-repair-theorem}
Assume:
\begin{enumerate}[leftmargin=2em,label=(\alph*)]
  \item the first local Theta-drift witness theorem applies;
  \item the first local phase-side repair theorem applies;
  \item the first local Teichm\"uller-persistence repair theorem applies;
  \item the first Teichm\"uller-side synchronization orientation theorem applies;
  \item a minimal local phase-lock repair datum exists.
\end{enumerate}
Then the current post-v3.32.0 line supports a first concrete phase-lock repair candidate.

More precisely:
\begin{enumerate}[leftmargin=2em,label=(\roman*)]
  \item the primary repair target is
  \[
    \sigma_{\Theta}^{\mathrm{lock}}(n) \to 0;
  \]
  \item the geometric carry repair remains genuinely active through
  \[
    \ell_T^{\mathrm{lock}}(n)>0
  \]
  whenever pure phase-side re-alignment does not by itself stabilize the repaired second step;
  \item a phase-lock repaired local window is the first local candidate for reopening the hybrid construction gate.
\end{enumerate}
\end{theorem}

\begin{proof}
Items (a)--(d) identify the present obstruction chain: the first failure is Theta-drift-prior, the first repair is burden-side with geometric carry, the next residual obstruction is Teichm\"uller-side-prior, and that obstruction is phase-lock-prior rather than raw-geometry-prior. Therefore the first concrete local repair must target vanishing of $\sigma_{\Theta}^{\mathrm{lock}}(n)$. If pure re-alignment is not enough to keep the repaired second step in the same local Teichm\"uller-safe band, a positive carry correction $\ell_T^{\mathrm{lock}}(n)$ remains necessary. This yields the first concrete phase-lock repair candidate.
\end{proof}

\begin{corollary}[Current phase-lock repair working rule]
\label{cor:current-phase-lock-repair-working-rule}
The next local audit should no longer ask merely whether the repaired second step is better than before. It should ask whether the applied local correction actually forces
\[
  \sigma_{\Theta}^{\mathrm{lock}}(n)=0
\]
and, if not, which remaining carry correction still prevents a fully phase-locked repaired window.
\end{corollary}

\begin{proof}
Immediate from Theorem~\ref{thm:first-phase-lock-repair-theorem}.
\end{proof}



\begin{definition}[Local Heisenberg compensator]
\label{def:local-heisenberg-compensator}
A \emph{local Heisenberg compensator} at stage $n$ is a conservative local synchronization operator
\[
  \mathrm{HC}_{\mathrm{loc}}(n):(X_n,X_{n+1})\longmapsto (X_n,\widetilde X_{n+1})
\]
which acts only on the repaired second step and is intended to reduce the remaining phase-lock and Teichm\"uller-carry defects without changing the already accepted first-step near-trigger reading.
\end{definition}

\begin{definition}[HC-admissible local compensation]
\label{def:hc-admissible-local-compensation}
A local Heisenberg compensator is called \emph{HC-admissible} if it satisfies:
\begin{enumerate}[leftmargin=2em,label=(\alph*)]
  \item the closure burden remains controlled;
  \item the corrected second step stays inside the same admissibility band;
  \item the phase-lock defect decreases:
  \[
    \sigma_{\Theta}^{\mathrm{lock}}(n;\widetilde X_{n+1})
    <
    \sigma_{\Theta}^{\mathrm{lock}}(n;X_{n+1});
  \]
  \item the Teichm\"uller-carry defect does not increase, and preferably also decreases;
  \item the compensator does not manufacture a false entry certificate by destroying the route-faithful burden reading.
\end{enumerate}
\end{definition}

\begin{remark}[How the HC is read here]
\label{rem:how-the-hc-is-read-here}
The Heisenberg compensator is not used here as a new bulk law. It is read as the minimal local re-phasing operator that tries to snap the repaired second step back into the same narrow closure-phase corridor as the first near-trigger step.
\end{remark}

\begin{definition}[Bombe-guided compensator menu]
\label{def:bombe-guided-compensator-menu}
A \emph{Bombe-guided compensator menu} at stage $n$ is a finite family
\[
  \mathcal M_{\mathrm{HC}}(n)=\{H_1,\dots,H_r\}
\]
of candidate local Heisenberg compensators together with a contradiction-sensitive elimination rule that removes every candidate violating closure control, admissibility, or route-faithful burden preservation.
\end{definition}

\begin{remark}[Why the Bombe idea fits here]
\label{rem:why-the-bombe-idea-fits-here}
The present problem is not to guess one magical compensator. It is to eliminate locally incompatible synchronizers from a finite candidate menu until only the route-faithful phase-lock-capable ones remain. This is exactly the right local role for a Bombe-style elimination discipline.
\end{remark}

\begin{definition}[Compensator filter tower]
\label{def:compensator-filter-tower}
A \emph{compensator filter tower} is a finite ordered stack
\[
  \mathfrak F_{\mathrm{HC}}(n)=
  (F_1,\dots,F_s)
\]
of local HC-compatible filters through which the repaired second step is passed before the final trigger test. The intended reading is:
\begin{enumerate}[leftmargin=2em,label=(\roman*)]
  \item burden-side pre-alignment;
  \item phase-lock re-alignment;
  \item Teichm\"uller-carry stabilization;
  \item final closure-phase re-entry test.
\end{enumerate}
\end{definition}

\begin{remark}[Filter tower and the electromagnetic analogy]
\label{rem:filter-tower-and-electromagnetic-analogy}
The filter-tower language should be read structurally rather than as a literal new electrodynamics. It expresses that several local synchronization layers may act in series, each removing one kind of phase/carry mismatch before the final trigger decision is taken.
\end{remark}

\begin{definition}[Readout-relative informational gravitation principle]
\label{def:readout-relative-informational-gravitation-principle}
A \emph{readout-relative informational gravitation principle} is the principle that the apparent curvature/attraction acting on a local state is not necessarily a primitive force attached only to rest mass, but may be read as the effective curvature law induced by the chosen admissible readout on the underlying information-carrying structure.
\end{definition}

\begin{remark}[How far this principle should be claimed]
\label{rem:how-far-this-principle-should-be-claimed}
In the present document this principle is only an internal architecture hypothesis. It does not claim multiple literal physical gravities. It claims instead that one deeper ART/QFT-compatible carrier may have different gravitational-like readout shadows depending on the active admissible information/readout channel.
\end{remark}

\begin{proposition}[HC can serve as a local phase-lock operator]
\label{prop:hc-can-serve-as-local-phase-lock-operator}
Assume an HC-admissible local compensation exists at stage $n$. Then the Heisenberg compensator can serve as a local phase-lock operator for the repaired second step.
\end{proposition}

\begin{proof}
By HC-admissibility, the compensator reduces the phase-lock defect while preserving closure control, admissibility, and the route-faithful burden reading. Therefore it acts exactly as the required local phase-lock operator.
\end{proof}

\begin{lemma}[Bombe-guided HC elimination lemma]
\label{lem:bombe-guided-hc-elimination-lemma}
Assume the compensator candidate menu is finite. Then Bombe-guided elimination reduces the local HC search to the subfamily of candidates that are simultaneously closure-compatible, admissibility-safe, and phase-lock-improving.
\end{lemma}

\begin{proof}
By definition, Bombe-guided elimination removes every candidate violating one of the required local compatibility conditions. Hence the surviving menu consists exactly of those candidates satisfying the required conditions.
\end{proof}

\begin{theorem}[First local Heisenberg compensator theorem]
\label{thm:first-local-heisenberg-compensator-theorem}
Assume:
\begin{enumerate}[leftmargin=2em,label=(\alph*)]
  \item the first phase-lock repair theorem applies;
  \item a finite Bombe-guided compensator menu is available;
  \item at least one surviving candidate is HC-admissible.
\end{enumerate}
Then the current post-v3.32.0 line supports a first local Heisenberg-compensated phase-lock repair.

More precisely:
\begin{enumerate}[leftmargin=2em,label=(\roman*)]
  \item the repaired second step can be re-phased by a surviving HC candidate;
  \item the corrected step is brought closer to the same closure-phase corridor as the first near-trigger step;
  \item the compensator acts as a synchronization layer between the local crystal correction and the Teichm\"uller-side persistence repair.
\end{enumerate}
\end{theorem}

\begin{proof}
By the first phase-lock repair theorem, the current line already isolates the local phase-lock defect as the next concrete obstruction. By the Bombe-guided HC elimination lemma, the finite menu reduces to the candidates that are compatible with closure, admissibility, and phase-lock improvement. If one surviving candidate is HC-admissible, then by Proposition~\ref{prop:hc-can-serve-as-local-phase-lock-operator} it provides the required local re-phasing. Therefore the Heisenberg compensator acts as the missing synchronization layer between the crystal-side correction and the Teichm\"uller-side persistence repair.
\end{proof}

\begin{corollary}[Current HC working rule]
\label{cor:current-hc-working-rule}
The next local audit should not ask first for a new bulk equation. It should ask first for the smallest HC-admissible local compensator that phase-locks the repaired second step back into the same closure-phase corridor while preserving closure control and admissibility.
\end{corollary}

\begin{proof}
Immediate from Theorem~\ref{thm:first-local-heisenberg-compensator-theorem}.
\end{proof}



\begin{definition}[Phase-side dominant HC candidate]
\label{def:phase-side-dominant-hc-candidate}
A local Heisenberg compensator candidate at stage $n$ is called \emph{phase-side dominant} if its primary improvement acts on the phase-lock defect, in the sense that it strictly lowers
\[
  \sigma_{\Theta}^{\mathrm{lock}}(n)
\]
while any improvement of the Teichm\"uller-carry defect remains subordinate to that phase-lock gain.
\end{definition}

\begin{definition}[Carry-dominant HC candidate]
\label{def:carry-dominant-hc-candidate}
A local Heisenberg compensator candidate at stage $n$ is called \emph{carry-dominant} if its primary improvement acts on the geometric carry defect, in the sense that it lowers
\[
  \sigma_T^{\mathrm{carry}}(n)
\]
without first removing the phase-lock misalignment.
\end{definition}

\begin{definition}[Two-component HC candidate]
\label{def:two-component-hc-candidate}
A local Heisenberg compensator candidate at stage $n$ is called \emph{two-component coupled} if it lowers both
\[
  \sigma_{\Theta}^{\mathrm{lock}}(n)
  \qquad\text{and}\qquad
  \sigma_T^{\mathrm{carry}}(n)
\]
with neither correction honestly subordinate to the other at first local order.
\end{definition}

\begin{remark}[What the classification is trying to decide]
\label{rem:what-the-hc-classification-is-trying-to-decide}
The present question is not whether a local HC helps at all. It asks which kind of local HC is the first plausible one for the current line: one that primarily restores phase-lock, one that primarily carries geometry, or one that is already fully two-component at first local order.
\end{remark}

\begin{proposition}[The current line is not carry-dominant on the HC side]
\label{prop:current-line-is-not-carry-dominant-on-the-hc-side}
Under the present post-v3.32.0 reading, the first plausible local HC candidate is not carry-dominant.
\end{proposition}

\begin{proof}
By the first Teichm\"uller-side synchronization orientation theorem and the first phase-lock repair theorem, the current residual obstruction is phase-lock-prior rather than raw-carry-prior. Therefore a compensator that tries first to optimize geometric carry without first restoring phase-lock cannot be the default first candidate.
\end{proof}

\begin{proposition}[The current line is not yet fully two-component on the HC side]
\label{prop:current-line-is-not-yet-fully-two-component-on-the-hc-side}
Under the present post-v3.32.0 reading, the first plausible local HC candidate is not yet fully two-component at first local order.
\end{proposition}

\begin{proof}
The current local repair line remains asymptotically burden-dominant and phase-lock-prior. Although the geometric carry correction is genuinely active, the present line has not yet produced a positive hybrid balancing witness with balanced first-order contributions. Hence the first plausible HC candidate cannot yet be classified as fully two-component at first local order.
\end{proof}

\begin{theorem}[First HC-candidate type orientation theorem]
\label{thm:first-hc-candidate-type-orientation-theorem}
Assume:
\begin{enumerate}[leftmargin=2em,label=(\alph*)]
  \item the first local Heisenberg compensator theorem applies;
  \item the first Teichm\"uller-side synchronization orientation theorem applies;
  \item the first phase-lock repair theorem applies;
  \item the current line has not yet produced a positive hybrid balancing witness.
\end{enumerate}
Then the first plausible local HC candidate is phase-side dominant rather than carry-dominant or fully two-component.

More precisely:
\begin{enumerate}[leftmargin=2em,label=(\roman*)]
  \item the primary local HC task is to reduce
  \[
    \sigma_{\Theta}^{\mathrm{lock}}(n);
  \]
  \item the Teichm\"uller-side carry correction remains genuinely active as a supporting synchronization channel;
  \item the current HC candidate should therefore be read as phase-side dominant but geometrically carried.
\end{enumerate}
\end{theorem}

\begin{proof}
By Proposition~\ref{prop:current-line-is-not-carry-dominant-on-the-hc-side}, the current line is not blocked first by raw geometric carry. By Proposition~\ref{prop:current-line-is-not-yet-fully-two-component-on-the-hc-side}, the current line has not yet reached a first-order balanced two-component compensator regime. Since the active residual obstruction is phase-lock-prior, the first plausible HC candidate must therefore act primarily on the phase-lock defect, with geometric carry serving as the supporting channel that stabilizes the repaired lock.
\end{proof}

\begin{corollary}[Current HC-candidate working rule]
\label{cor:current-hc-candidate-working-rule}
The next local audit should not search first for a carry-dominant or perfectly balanced HC. It should search first for the smallest phase-side dominant compensator whose geometric carry contribution is just strong enough to keep the repaired second step inside the same closure-phase corridor.
\end{corollary}

\begin{proof}
Immediate from Theorem~\ref{thm:first-hc-candidate-type-orientation-theorem}.
\end{proof}


\begin{definition}[Local lock corridor]
\label{def:local-lock-corridor}
A \emph{local lock corridor} at stage $n$ is a narrow admissible closure-phase band
\[
  \mathfrak C_{\mathrm{lock}}^{\mathrm{loc}}(n)
\]
containing the first near-trigger step and characterized by the following properties:
\begin{enumerate}[leftmargin=2em,label=(\alph*)]
  \item every point in $\mathfrak C_{\mathrm{lock}}^{\mathrm{loc}}(n)$ lies inside the same local closure-phase window;
  \item the local burden pair remains trigger-near throughout the corridor;
  \item the Teichm\"uller-side distortion stays inside the active safe band;
  \item the corridor is narrow enough that leaving it produces a detectable phase-lock or carry defect.
\end{enumerate}
\end{definition}

\begin{remark}[Geometric meaning of the lock corridor]
\label{rem:geometric-meaning-of-the-lock-corridor}
The lock corridor is the thin local channel in which both the first near-trigger step and the repaired second step must lie if one is to read them as belonging to the same closure-phase passage. The current obstruction does not come from total loss of that channel, but from the failure of the second step to remain locked inside it.
\end{remark}

\begin{definition}[HC-captured repaired second step]
\label{def:hc-captured-repaired-second-step}
A repaired second step at stage $n$ is called \emph{HC-captured} if there exists an HC-admissible local Heisenberg compensator such that the compensated second step lies in the same local lock corridor as the first near-trigger step.
\end{definition}

\begin{definition}[Lock-corridor leakage defect]
\label{def:lock-corridor-leakage-defect}
The \emph{lock-corridor leakage defect} at stage $n$ is the residual amount by which the repaired second step fails to remain inside
\[
  \mathfrak C_{\mathrm{lock}}^{\mathrm{loc}}(n).
\]
It is denoted by
\[
  \lambda_{\mathrm{lock}}(n)\ge 0,
\]
and vanishes exactly when the repaired second step is HC-captured.
\end{definition}

\begin{definition}[Minimal local lock-corridor compensation]
\label{def:minimal-local-lock-corridor-compensation}
A local HC candidate is called \emph{minimally lock-corridor compensating} if it is HC-admissible and reduces the lock-corridor leakage defect to zero with no strictly smaller phase-side dominant correction achieving the same result on the active two-step local window.
\end{definition}

\begin{proposition}[Current line still exhibits lock-corridor leakage]
\label{prop:current-line-still-exhibits-lock-corridor-leakage}
Under the present post-v3.32.0 reading, the repaired second step still exhibits a nonzero lock-corridor leakage defect.
\end{proposition}

\begin{proof}
The current line is phase-lock-prior on the Teichm\"uller side and has not yet produced a positive hybrid balancing witness. Hence the repaired second step is not yet fully locked into the same local closure-phase passage as the first near-trigger step, so the local lock-corridor leakage defect remains nonzero.
\end{proof}

\begin{lemma}[Phase-side dominant HC yields lock-corridor capture]
\label{lem:phase-side-dominant-hc-yields-lock-corridor-capture}
Assume:
\begin{enumerate}[leftmargin=2em,label=(\alph*)]
  \item a local lock corridor exists;
  \item the current HC candidate is phase-side dominant but geometrically carried;
  \item the compensator reduces the phase-lock defect to zero while keeping admissibility and the geometric carry inside the safe band.
\end{enumerate}
Then the compensated second step is HC-captured.
\end{lemma}

\begin{proof}
By assumption, the compensator restores the phase-lock condition and preserves the geometric safe-band control. Since the local lock corridor is defined exactly as the admissible narrow closure-phase band carrying the first near-trigger step, the compensated second step now lies in the same corridor. Hence it is HC-captured.
\end{proof}

\begin{theorem}[First explicit lock-corridor theorem]
\label{thm:first-explicit-lock-corridor-theorem}
Assume:
\begin{enumerate}[leftmargin=2em,label=(\alph*)]
  \item the first HC-candidate type orientation theorem applies;
  \item the first phase-lock repair theorem applies;
  \item a local lock corridor exists around the first near-trigger step;
  \item there exists a minimally lock-corridor compensating HC candidate.
\end{enumerate}
Then the current post-v3.32.0 line admits a first explicit local lock-corridor reading.

More precisely:
\begin{enumerate}[leftmargin=2em,label=(\roman*)]
  \item the geometric problem is not to invent a new bulk path, but to return the repaired second step to the same narrow local closure-phase corridor as the first step;
  \item the current HC task is to drive
  \[
    \lambda_{\mathrm{lock}}(n)\to 0;
  \]
  \item once the leakage defect vanishes under an HC-admissible phase-side dominant compensator, the repaired second step is locally synchronized in the same lock corridor as the first near-trigger step.
\end{enumerate}
\end{theorem}

\begin{proof}
By the HC-candidate orientation theorem, the first plausible compensator is phase-side dominant but geometrically carried. By the phase-lock repair theorem, the primary local repair target is vanishing of the phase-lock defect, with geometric carry remaining genuinely active. The local lock corridor isolates the thin admissible passage already occupied by the first near-trigger step. A minimally lock-corridor compensating HC candidate removes the leakage defect while preserving admissibility, so by Lemma~\ref{lem:phase-side-dominant-hc-yields-lock-corridor-capture} the repaired second step is returned to the same corridor.
\end{proof}

\begin{corollary}[Current lock-corridor working rule]
\label{cor:current-lock-corridor-working-rule}
The next local audit should no longer ask only whether the second step is improved. It should ask whether an HC-admissible compensator has actually reduced
\[
  \lambda_{\mathrm{lock}}(n)
\]
to zero, i.e. whether the repaired second step has truly re-entered the same local lock corridor as the first near-trigger step.
\end{corollary}

\begin{proof}
Immediate from Theorem~\ref{thm:first-explicit-lock-corridor-theorem}.
\end{proof}



\begin{definition}[Lock-side leakage defect]
\label{def:lock-side-leakage-defect}
A \emph{lock-side leakage defect} at stage $n$ is the portion of the local leakage defect attributable to failure of the repaired second step to satisfy the phase-lock condition under the current corridor calibration. It is denoted by
\[
  \lambda_{\mathrm{lock}}^{\Theta}(n)\ge 0.
\]
\end{definition}

\begin{definition}[Corridor-calibration leakage defect]
\label{def:corridor-calibration-leakage-defect}
A \emph{corridor-calibration leakage defect} at stage $n$ is the portion of the local leakage defect attributable to the active local lock corridor being too narrow or otherwise miscalibrated for the repaired second step, even after the primary phase-lock correction has been applied. It is denoted by
\[
  \lambda_{\mathrm{lock}}^{\mathrm{cal}}(n)\ge 0.
\]
\end{definition}

\begin{remark}[How the leakage splits]
\label{rem:how-the-leakage-splits}
The total local lock-corridor leakage is now read as having two possible sources: a remaining failure of lock restoration itself, or a failure of the current corridor calibration to reflect the true admissible narrow passage. In symbols, one may think heuristically of
\[
  \lambda_{\mathrm{lock}}(n)
  \sim
  \lambda_{\mathrm{lock}}^{\Theta}(n)
  +
  \lambda_{\mathrm{lock}}^{\mathrm{cal}}(n),
\]
without yet forcing a rigid additive identity.
\end{remark}

\begin{definition}[Lock-side-prior local leakage regime]
\label{def:lock-side-prior-local-leakage-regime}
The local line is said to be in a \emph{lock-side-prior local leakage regime} at stage $n$ if
\[
  \lambda_{\mathrm{lock}}^{\Theta}(n)
  >
  \lambda_{\mathrm{lock}}^{\mathrm{cal}}(n).
\]
\end{definition}

\begin{definition}[Corridor-calibration-prior local leakage regime]
\label{def:corridor-calibration-prior-local-leakage-regime}
The local line is said to be in a \emph{corridor-calibration-prior local leakage regime} at stage $n$ if
\[
  \lambda_{\mathrm{lock}}^{\mathrm{cal}}(n)
  \ge
  \lambda_{\mathrm{lock}}^{\Theta}(n).
\]
\end{definition}

\begin{proposition}[The current line is not corridor-calibration-prior by default]
\label{prop:current-line-is-not-corridor-calibration-prior-by-default}
Under the present post-v3.32.0 reading, the local line should not yet be read as corridor-calibration-prior by default.
\end{proposition}

\begin{proof}
The current diagnosis remains phase-lock-prior on the Teichm\"uller side, and the first local repair target is still the vanishing of the phase-lock defect rather than widening or re-centering the lock corridor itself. Hence the default reading is not that the corridor is already known to be primarily miscalibrated, but rather that the repaired second step still fails first to lock into the active corridor.
\end{proof}

\begin{proposition}[The current line is lock-side-prior on the leakage level]
\label{prop:current-line-is-lock-side-prior-on-the-leakage-level}
Under the present post-v3.32.0 reading, the first remaining local lock-corridor leakage is lock-side-prior.
\end{proposition}

\begin{proof}
By the first phase-lock repair theorem and the current phase-lock repair working rule, the next local task is still to force
\[
  \sigma_{\Theta}^{\mathrm{lock}}(n)=0.
\]
This means that the dominant unresolved leakage continues to arise from incomplete phase-lock restoration rather than from demonstrated over-narrowness or miscentering of the corridor itself. Therefore the first remaining local leakage is lock-side-prior.
\end{proof}

\begin{theorem}[First lock-corridor leakage orientation theorem]
\label{thm:first-lock-corridor-leakage-orientation-theorem}
Assume:
\begin{enumerate}[leftmargin=2em,label=(\alph*)]
  \item the first explicit lock-corridor theorem applies;
  \item the first phase-lock repair theorem applies;
  \item the current line still exhibits nonzero lock-corridor leakage.
\end{enumerate}
Then the current post-v3.32.0 line is lock-side-prior on the leakage level rather than corridor-calibration-prior.

More precisely:
\begin{enumerate}[leftmargin=2em,label=(\roman*)]
  \item the next honest reduction target is first
  \[
    \lambda_{\mathrm{lock}}^{\Theta}(n)\downarrow;
  \]
  \item local corridor recalibration should be treated only as a secondary follow-on question unless the lock-side leakage has already been made negligible;
  \item the present line should therefore first seek a stronger HC lock action before widening or re-centering the local corridor.
\end{enumerate}
\end{theorem}

\begin{proof}
By the first explicit lock-corridor theorem, the active local task is to return the repaired second step to the same narrow corridor as the first near-trigger step. By the first phase-lock repair theorem, the primary local repair target remains vanishing of the phase-lock defect. Therefore, as long as nonzero leakage remains, the default reading is that the dominant leakage sits on the lock-restoration side rather than already on the calibration side of the corridor. Hence the line is lock-side-prior on the leakage level.
\end{proof}

\begin{corollary}[Current leakage working rule]
\label{cor:current-leakage-working-rule}
The next local audit should not first widen, thicken, or re-center the active lock corridor. It should first ask whether the remaining leakage is still produced by incomplete lock restoration of the repaired second step under the current corridor geometry.
\end{corollary}

\begin{proof}
Immediate from Theorem~\ref{thm:first-lock-corridor-leakage-orientation-theorem}.
\end{proof}



\begin{definition}[Minimal local HC-lock step]
\label{def:minimal-local-hc-lock-step}
A \emph{minimal local HC-lock step} at stage $n$ is an HC-admissible local compensation
\[
  \mathcal H^{\min}_{\mathrm{lock}}(n)
\]
which acts on the repaired second step only through the smallest phase-side dominant correction required to reduce the lock-side leakage defect from
\[
  \lambda_{\mathrm{lock}}^{\Theta}(n)>0
\]
to
\[
  \lambda_{\mathrm{lock}}^{\Theta,+}(n)=0,
\]
while preserving the active local lock corridor geometry.
\end{definition}

\begin{definition}[Lock-neutral corridor preservation]
\label{def:lock-neutral-corridor-preservation}
A local HC-lock step is said to satisfy \emph{lock-neutral corridor preservation} if it does not widen, thicken, or re-center the current local lock corridor
\[
  \mathfrak C_{\mathrm{lock}}^{\mathrm{loc}}(n),
\]
but only changes the repaired second-step placement inside that already fixed corridor geometry.
\end{definition}

\begin{definition}[HC-lock effectiveness defect]
\label{def:hc-lock-effectiveness-defect}
The \emph{HC-lock effectiveness defect} of a local HC compensation at stage $n$ is
\[
  E_{\mathrm{HC}}^{\mathrm{lock}}(n)
  :=
  \lambda_{\mathrm{lock}}^{\Theta,+}(n),
\]
that is, the residual lock-side leakage after the local HC action has been applied.
\end{definition}

\begin{remark}[What the minimal HC step is allowed to do]
\label{rem:what-the-minimal-hc-step-is-allowed-to-do}
The minimal HC-lock step is not allowed to solve the problem by changing the corridor itself. It must work only by re-locking the repaired second step inside the already accepted corridor. Geometrically, it is the smallest admissible local twist that makes the second point snap back into the same narrow channel without moving the channel walls.
\end{remark}

\begin{definition}[HC-captured local repaired second step]
\label{def:hc-captured-local-repaired-second-step}
A repaired second step at stage $n$ is called \emph{HC-captured} if there exists an HC-admissible local compensation satisfying lock-neutral corridor preservation such that
\[
  \lambda_{\mathrm{lock}}^{\Theta,+}(n)=0.
\]
\end{definition}

\begin{proposition}[The current line calls first for a minimal HC-lock step]
\label{prop:the-current-line-calls-first-for-a-minimal-hc-lock-step}
Under the current post-v3.32.0 reading, the next local correction should first be sought as a minimal local HC-lock step rather than as a corridor recalibration.
\end{proposition}

\begin{proof}
By the first lock-corridor leakage orientation theorem, the active leakage is lock-side-prior rather than corridor-calibration-prior. Therefore the first honest correction should try to eliminate the residual lock-side leakage while preserving the existing corridor geometry. This is exactly the role of a minimal local HC-lock step.
\end{proof}

\begin{lemma}[Minimal HC-lock step implies HC capture]
\label{lem:minimal-hc-lock-step-implies-hc-capture}
Assume a minimal local HC-lock step exists at stage $n$. Then the repaired second step is HC-captured at that stage.
\end{lemma}

\begin{proof}
By definition, a minimal local HC-lock step is HC-admissible, preserves the corridor geometry, and reduces the residual lock-side leakage defect to zero. Hence the repaired second step satisfies the defining conditions of an HC-captured local repaired second step.
\end{proof}

\begin{theorem}[First minimal HC-lock step theorem]
\label{thm:first-minimal-hc-lock-step-theorem}
Assume:
\begin{enumerate}[leftmargin=2em,label=(\alph*)]
  \item the first local Theta-drift witness theorem applies;
  \item the first phase-lock repair theorem applies;
  \item the first explicit lock-corridor theorem applies;
  \item the first lock-corridor leakage orientation theorem applies;
  \item an HC-admissible local compensation menu is available.
\end{enumerate}
Then the current post-v3.32.0 line admits the following first local repair target:
\begin{enumerate}[leftmargin=2em,label=(\roman*)]
  \item search first for a minimal local HC-lock step;
  \item require lock-neutral corridor preservation;
  \item accept the first candidate only if it yields
  \[
    E_{\mathrm{HC}}^{\mathrm{lock}}(n)=0;
  \]
  \item treat any remaining failure first as lack of sufficient HC lock action rather than as corridor miscalibration.
\end{enumerate}
\end{theorem}

\begin{proof}
Items (a)--(d) identify the current obstruction chain: the line is phase-lock-prior, the local corridor is already the right object to preserve, and the leakage is lock-side-prior rather than corridor-calibration-prior. Hence the first honest local HC task is not to redesign the corridor but to search within the available HC menu for the smallest admissible lock action that drives the residual lock-side leakage to zero while leaving the corridor fixed. This is exactly the asserted minimal HC-lock step search.
\end{proof}

\begin{corollary}[Current minimal HC working rule]
\label{cor:current-minimal-hc-working-rule}
The next local audit should not first modify the local lock corridor. It should first scan the local HC menu for the smallest phase-side dominant compensator whose geometric carry is just sufficient to force
\[
  \lambda_{\mathrm{lock}}^{\Theta,+}(n)=0
\]
under lock-neutral corridor preservation.
\end{corollary}

\begin{proof}
Immediate from Theorem~\ref{thm:first-minimal-hc-lock-step-theorem}.
\end{proof}



\begin{definition}[Local HC candidate family]
\label{def:local-hc-candidate-family}
The \emph{local HC candidate family} at stage $n$ is a finite admissible menu
\[
  \mathfrak M_{\mathrm{HC}}^{\mathrm{loc}}(n)
  =
  \{H_1^{(n)},\dots,H_{N(n)}^{(n)}\}
\]
of local Heisenberg-compensator actions acting on the repaired second step under the fixed local lock corridor.
\end{definition}

\begin{definition}[HC menu admissibility filter]
\label{def:hc-menu-admissibility-filter}
A local HC candidate
\[
  H_j^{(n)}\in \mathfrak M_{\mathrm{HC}}^{\mathrm{loc}}(n)
\]
passes the \emph{HC menu admissibility filter} if:
\begin{enumerate}[leftmargin=2em,label=(\alph*)]
  \item it is phase-side dominant with genuinely active geometric carry;
  \item it preserves the active local lock corridor;
  \item it preserves local admissibility on the same two-step window;
  \item it does not enlarge the corridor merely to hide the residual leakage.
\end{enumerate}
\end{definition}

\begin{definition}[HC lock score]
\label{def:hc-lock-score}
The \emph{HC lock score} of an admissible local HC candidate is the pair
\[
  \mathfrak L_{\mathrm{HC}}(H_j^{(n)})
  :=
  \bigl(
    \lambda_{\mathrm{lock}}^{\Theta,+}(n;H_j^{(n)}),
    C_{\mathrm{HC}}(H_j^{(n)})
  \bigr),
\]
where the first component is the residual lock-side leakage after applying the candidate and the second component is its local compensator cost.
\end{definition}

\begin{definition}[Bombe-guided HC elimination]
\label{def:bombe-guided-hc-elimination}
The \emph{bombe-guided HC elimination} is the finite contradiction-sensitive elimination of all candidates in
\[
  \mathfrak M_{\mathrm{HC}}^{\mathrm{loc}}(n)
\]
that fail the HC menu admissibility filter or leave a residual lock score incompatible with minimal HC capture.
\end{definition}

\begin{remark}[Why the HC scan is finite]
\label{rem:why-the-hc-scan-is-finite}
The present line does not search over an arbitrary continuum of compensators. It searches only over the finite local HC menu already licensed by the current route-faithful readout state. This is why the bombe-guided elimination picture is appropriate here.
\end{remark}

\begin{definition}[Minimal surviving HC candidate]
\label{def:minimal-surviving-hc-candidate}
A candidate
\[
  H_{\min}^{(n)}\in \mathfrak M_{\mathrm{HC}}^{\mathrm{loc}}(n)
\]
is called the \emph{minimal surviving HC candidate} if:
\begin{enumerate}[leftmargin=2em,label=(\alph*)]
  \item it survives bombe-guided HC elimination;
  \item it minimizes the residual lock-side leakage
  \[
    \lambda_{\mathrm{lock}}^{\Theta,+}(n;H_j^{(n)});
  \]
  \item among all candidates with the same minimal leakage, it minimizes the compensator cost
  \[
    C_{\mathrm{HC}}(H_j^{(n)}).
  \]
\end{enumerate}
\end{definition}

\begin{proposition}[The current line reduces the HC problem to a finite menu scan]
\label{prop:current-line-reduces-hc-to-finite-menu-scan}
Under the current post-v3.32.0 reading, the first local HC problem is reduced to scanning the finite local HC candidate family under admissibility and lock-corridor preservation.
\end{proposition}

\begin{proof}
By the current minimal HC working rule, the line must first search for a minimal local HC-lock step while preserving the active lock corridor. By the local Heisenberg-compensator theorem and the bombe-guided compensator menu reading, the admissible compensator actions are already organized as a finite local menu. Hence the first local HC problem is reduced to a finite menu scan with admissibility and corridor-preservation constraints.
\end{proof}

\begin{lemma}[Minimal surviving HC candidate yields the first executable HC-lock step]
\label{lem:minimal-surviving-hc-candidate-yields-first-executable-hc-lock-step}
Assume a minimal surviving HC candidate exists. Then it yields the first executable local HC-lock step under the current lock-side-prior regime.
\end{lemma}

\begin{proof}
A minimal surviving HC candidate has already passed the admissibility and corridor-preservation filter. By minimality, it leaves no smaller admissible compensator with strictly lower residual lock-side leakage. Hence it is the first executable local HC-lock step compatible with the current regime.
\end{proof}

\begin{theorem}[First local HC menu-scan theorem]
\label{thm:first-local-hc-menu-scan-theorem}
Assume:
\begin{enumerate}[leftmargin=2em,label=(\alph*)]
  \item the first minimal HC-lock step theorem applies;
  \item the current line is lock-side-prior on the leakage level;
  \item the local HC candidate family is finite;
  \item bombe-guided HC elimination is available.
\end{enumerate}
Then the next local HC task is a finite menu-scan problem.

More precisely:
\begin{enumerate}[leftmargin=2em,label=(\roman*)]
  \item all inadmissible or corridor-breaking HC candidates may be eliminated first;
  \item the surviving candidates are ranked by residual lock-side leakage and compensator cost;
  \item if a minimal surviving HC candidate exists, it determines the first executable HC-lock step;
  \item if no such candidate exists, the current line records a local HC-menu exhaustion obstruction rather than immediately blaming corridor calibration.
\end{enumerate}
\end{theorem}

\begin{proof}
By Proposition~\ref{prop:current-line-reduces-hc-to-finite-menu-scan}, the present HC problem is already reduced to a finite candidate family with admissibility and corridor-preservation constraints. Bombe-guided HC elimination removes candidates that fail these constraints. The remaining candidates are compared by lock score. If a minimal surviving candidate exists, Lemma~\ref{lem:minimal-surviving-hc-candidate-yields-first-executable-hc-lock-step} yields the first executable HC-lock step. If not, the failure lies at the menu level itself, hence the correct reading is local HC-menu exhaustion rather than premature corridor redesign.
\end{proof}

\begin{corollary}[Current HC menu-scan working rule]
\label{cor:current-hc-menu-scan-working-rule}
The next local audit should first enumerate the finite HC menu, eliminate inadmissible and corridor-breaking candidates, and only then ask whether the surviving menu still contains a minimal lock-effective compensator.
\end{corollary}

\begin{proof}
Immediate from Theorem~\ref{thm:first-local-hc-menu-scan-theorem}.
\end{proof}



\begin{definition}[Phase-side micro-lock HC candidate]
\label{def:phase-side-micro-lock-hc-candidate}
A local HC candidate
\[
  H_{\mathrm{lock}}^{(1)}(n)\in \mathfrak M_{\mathrm{HC}}^{\mathrm{loc}}(n)
\]
is called a \emph{phase-side micro-lock HC candidate} if it satisfies:
\begin{enumerate}[leftmargin=2em,label=(\alph*)]
  \item phase-side dominance in the sense that its primary action falls on the residual lock component;
  \item lock-neutral corridor preservation;
  \item nonzero but minimal geometric carry sufficient to support the repaired second step inside the same local lock corridor;
  \item no extra crystal-side widening or recentering action.
\end{enumerate}
\end{definition}

\begin{remark}[Why this is the first plausible HC menu entry]
\label{rem:why-this-is-the-first-plausible-hc-menu-entry}
The present line is phase-lock-prior, lock-side-prior on the leakage level, and not carry-dominant. Hence the first plausible HC candidate is not a broad geometric carrier and not a corridor redesign. It is the smallest local phase-side micro-lock action whose geometric carry is only as large as needed to prevent renewed leakage.
\end{remark}

\begin{definition}[HC micro-lock score]
\label{def:hc-micro-lock-score}
The \emph{HC micro-lock score} of a phase-side micro-lock candidate
\[
  H_{\mathrm{lock}}^{(1)}(n)
\]
is the tuple
\[
  \mathfrak L_{\mathrm{HC}}^{\mathrm{micro}}(H_{\mathrm{lock}}^{(1)}(n))
  :=
  \bigl(
    \lambda_{\mathrm{lock}}^{\Theta,+}(n),
    c_{\mathrm{HC}}(H_{\mathrm{lock}}^{(1)}(n))
  \bigr),
\]
ordered lexicographically by residual lock-side leakage first and compensator cost second.
\end{definition}

\begin{definition}[Plausible surviving HC candidate]
\label{def:plausible-surviving-hc-candidate}
A phase-side micro-lock HC candidate is called \emph{plausibly surviving} if it passes the HC menu admissibility filter, preserves the local lock corridor, and yields no worse lock score than any other currently known phase-side dominant candidate.
\end{definition}

\begin{proposition}[The current line favors a phase-side micro-lock candidate]
\label{prop:current-line-favors-a-phase-side-micro-lock-candidate}
Under the present post-v3.32.0 reading, the first plausible HC menu entry is phase-side micro-lock rather than carry-dominant or fully two-component balanced.
\end{proposition}

\begin{proof}
By the first HC-candidate type orientation theorem, the current line is phase-side dominant on the HC level and is neither carry-dominant nor already genuinely two-component balanced. By the first lock-corridor leakage orientation theorem, the residual leakage is lock-side-prior rather than corridor-calibration-prior. Therefore the first plausible HC entry is the smallest phase-side dominant candidate that closes the lock-side leakage while preserving the existing corridor.
\end{proof}

\begin{lemma}[Plausible survival of the phase-side micro-lock candidate]
\label{lem:plausible-survival-of-the-phase-side-micro-lock-candidate}
Assume:
\begin{enumerate}[leftmargin=2em,label=(\alph*)]
  \item the local HC candidate family is finite;
  \item the current line is phase-side dominant on the HC level;
  \item lock-neutral corridor preservation is enforced.
\end{enumerate}
Then the first plausible surviving HC candidate, if it exists at all, may be taken from the phase-side micro-lock class.
\end{lemma}

\begin{proof}
Because corridor redesign is excluded and carry-dominant candidates are not favored by the current orientation, the admissible search is restricted first to phase-side dominant candidates. Within that class, the lexicographic HC micro-lock score selects the candidates with smallest residual leakage and then smallest cost. Hence any first plausible survivor, if one exists, may be taken from the phase-side micro-lock class.
\end{proof}

\begin{theorem}[First plausible HC menu-entry theorem]
\label{thm:first-plausible-hc-menu-entry-theorem}
Assume:
\begin{enumerate}[leftmargin=2em,label=(\alph*)]
  \item the first HC-candidate type orientation theorem applies;
  \item the first lock-corridor leakage orientation theorem applies;
  \item the first local HC menu-scan theorem applies.
\end{enumerate}
Then the first plausible HC menu entry is a phase-side micro-lock HC candidate.

More precisely:
\begin{enumerate}[leftmargin=2em,label=(\roman*)]
  \item the candidate is sought first in the phase-side dominant subclass;
  \item its geometric carry is only auxiliary but genuinely active;
  \item its admissibility is tested under lock-neutral corridor preservation;
  \item if a first plausible survivor exists, it determines the leading local HC hypothesis before any broader compensator class is tried.
\end{enumerate}
\end{theorem}

\begin{proof}
Items (a) and (b) imply that the line is not looking first for a carry-dominant or fully balanced HC candidate. Item (c) reduces the search to the finite local HC menu with corridor-preserving admissibility. By Lemma~\ref{lem:plausible-survival-of-the-phase-side-micro-lock-candidate}, any first plausible survivor may be taken from the phase-side micro-lock class. This yields the asserted first plausible HC menu entry.
\end{proof}

\begin{corollary}[Current HC entry working rule]
\label{cor:current-hc-entry-working-rule}
The next local audit should first test the phase-side micro-lock HC candidates in the finite menu, ranked by residual lock-side leakage and cost, before considering broader compensator types.
\end{corollary}

\begin{proof}
Immediate from Theorem~\ref{thm:first-plausible-hc-menu-entry-theorem}.
\end{proof}



\begin{definition}[Micro-lock update rule]
\label{def:micro-lock-update-rule}
Let
\[
  H_{\mathrm{lock}}^{(1)}(n)
\]
be a phase-side micro-lock HC candidate acting on the repaired second local step. The \emph{micro-lock update rule} is the local update
\[
  X_{n+1}^{\mathrm{rep}}
  \longmapsto
  X_{n+1}^{\mathrm{lock}}
\]
induced by $H_{\mathrm{lock}}^{(1)}(n)$ under the following constraints:
\begin{enumerate}[leftmargin=2em,label=(\alph*)]
  \item lock-neutral corridor preservation is maintained;
  \item the updated second step stays inside the same local admissibility band;
  \item the lock-side leakage does not increase;
  \item the phase-side component of the update is dominant, while the geometric carry remains auxiliary but genuinely active.
\end{enumerate}
\end{definition}

\begin{definition}[Successful micro-lock update]
\label{def:successful-micro-lock-update}
A micro-lock update is called \emph{successful} if it satisfies
\[
  \lambda_{\mathrm{lock}}^{\Theta,+}(n)=0
\]
while preserving the current local lock corridor and the active admissibility constraints.
\end{definition}

\begin{definition}[HC-captured local continuation]
\label{def:hc-captured-local-continuation}
A local continuation after stage $n$ is called \emph{HC-captured} if it is produced by a successful micro-lock update and the resulting repaired second step remains in the same local lock corridor as the first near-trigger step.
\end{definition}

\begin{remark}[Why capture is the right next target]
\label{rem:why-capture-is-the-right-next-target}
The present line no longer needs a broad compensator search. It needs the first concrete capture event in which the second repaired step is actually brought back into the same narrow local corridor. That is the smallest nontrivial success criterion beyond plausible HC menu survival.
\end{remark}

\begin{lemma}[Minimal HC-lock step yields local capture]
\label{lem:minimal-hc-lock-step-yields-local-capture}
Assume:
\begin{enumerate}[leftmargin=2em,label=(\alph*)]
  \item a minimal surviving HC candidate exists in the local menu;
  \item its associated micro-lock update is successful;
  \item lock-neutral corridor preservation is enforced.
\end{enumerate}
Then the updated repaired second step is HC-captured.
\end{lemma}

\begin{proof}
By success of the micro-lock update, the residual lock-side leakage is driven to zero. By lock-neutral corridor preservation, the active corridor is not widened or re-centered during the update. Hence the updated second step lies in the same local lock corridor as the first near-trigger step, which is exactly HC capture.
\end{proof}

\begin{definition}[First local HC-capture chain]
\label{def:first-local-hc-capture-chain}
A \emph{first local HC-capture chain} is a finite local sequence
\[
  X_n^{\mathrm{near}}
  \to
  X_{n+1}^{\mathrm{rep}}
  \to
  X_{n+1}^{\mathrm{lock}}
\]
consisting of:
\begin{enumerate}[leftmargin=2em,label=(\roman*)]
  \item a near-trigger first step;
  \item a repaired but not yet fully locked second step;
  \item a successful micro-lock update producing an HC-captured continuation.
\end{enumerate}
\end{definition}

\begin{proposition}[Current local line reduces to HC-capture chain search]
\label{prop:current-local-line-reduces-to-hc-capture-chain-search}
Under the present post-v3.32.0 reading, once the local corridor, the phase-lock repair target, and the HC menu scan are fixed, the next local task reduces to the search for a first local HC-capture chain.
\end{proposition}

\begin{proof}
The preceding line already fixes: the right corridor object, the lock-side-prior leakage orientation, the minimal HC-lock target, and the first plausible HC candidate family. No broader redesign is currently preferred. Therefore the next local task is exactly to determine whether one of the plausible HC candidates yields an actual capture of the repaired second step.
\end{proof}

\begin{theorem}[First local HC-capture chain theorem]
\label{thm:first-local-hc-capture-chain-theorem}
Assume:
\begin{enumerate}[leftmargin=2em,label=(\alph*)]
  \item the first explicit lock-corridor theorem applies;
  \item the first minimal HC-lock step theorem applies;
  \item the first plausible HC menu-entry theorem applies;
  \item a successful micro-lock update exists for the minimal surviving HC candidate.
\end{enumerate}
Then the current post-v3.32.0 line supports a first local HC-capture chain.

More precisely:
\begin{enumerate}[leftmargin=2em,label=(\roman*)]
  \item the repaired second step is first updated by a phase-side dominant, geometrically carried micro-lock compensator;
  \item the lock-side leakage is removed without redesigning the corridor;
  \item the resulting continuation is HC-captured in the same local lock corridor as the first near-trigger step.
\end{enumerate}
\end{theorem}

\begin{proof}
By the first plausible HC menu-entry theorem, the first plausible surviving HC candidate is sought in the phase-side micro-lock class. By the first minimal HC-lock step theorem, the active local HC task is precisely to search for the smallest such compensator that removes residual lock-side leakage while preserving the corridor. If such a micro-lock update is successful, Lemma~\ref{lem:minimal-hc-lock-step-yields-local-capture} yields HC capture. This is exactly the stated HC-capture chain.
\end{proof}

\begin{corollary}[Current HC-capture working rule]
\label{cor:current-hc-capture-working-rule}
The next local audit should no longer stop at plausible HC menu survival. It should ask whether the best phase-side micro-lock candidate actually yields a successful micro-lock update and hence the first local HC-capture chain.
\end{corollary}

\begin{proof}
Immediate from Theorem~\ref{thm:first-local-hc-capture-chain-theorem}.
\end{proof}

\begin{corollary}[Freeze threshold after v3.32.0]
\label{cor:freeze-threshold-after-v3320}
If the local line is stably organized by
\[
  \text{lock corridor}
  \to
  \text{minimal HC-lock step}
  \to
  \text{HC menu scan}
  \to
  \text{first local HC-capture chain},
\]
then the post-v3.32.0 line has reached a natural freeze threshold for a new minor release candidate.
\end{corollary}

\begin{proof}
At that point the local HC line is no longer merely diagnostic. It carries a coherent operational chain from the residual lock defect to the first concrete capture event. This is exactly the kind of self-contained new kernel that supports a new freeze candidate before audit.
\end{proof}



\subsection{Transport-stable HC capture and OP-entry handoff}
\label{subsec:transport-stable-hc-capture-op-entry-handoff}

\begin{definition}[Transport-stable HC-capture chain]
\label{def:transport-stable-hc-capture-chain}
Let
\[
  \mathcal C_{n}^{\mathrm{HC}}
  :=
  \bigl(
    X_n^{\mathrm{near}}
    \to
    X_{n+1}^{\mathrm{rep}}
    \to
    X_{n+1}^{\mathrm{lock}}
  \bigr)
\]
be a first local HC-capture chain in the sense of Definition~\ref{def:first-local-hc-capture-chain}.  It is called
\emph{transport-stable on an admissible window} $[n,n+k]$ if every route-faithful admissible transport step
\[
  \mathsf T_j:X_j\longmapsto X_{j+1},
  \qquad n\leq j<n+k,
\]
including the symmetrized forward/backward transport readout when that readout is active, preserves:
\begin{enumerate}[leftmargin=2em,label=(\alph*)]
  \item membership of the locked continuation in the transported lock corridor;
  \item non-increase of the lock-side leakage under the transported HC readout;
  \item the inherited admissibility and boundary-stop constraints of the active route;
  \item the distinction between horizon-side and stop-preserving boundary-side packets.
\end{enumerate}
Thus transport-stability is not a new closure theorem.  It is the statement that the local capture event survives the already-admissible route motion needed before OP-facing use.
\end{definition}

\begin{definition}[HC transport defect]
\label{def:hc-transport-defect}
For a captured local chain $\mathcal C_n^{\mathrm{HC}}$ and an admissible transport step $\mathsf T_j$, define the associated HC transport defect schematically by
\[
  \Delta_{\mathrm{tr}}^{\mathrm{HC}}(j)
  :=
  \max\Bigl\{
    d_{\mathfrak C}\bigl(\mathsf T_j X_{j}^{\mathrm{lock}},\mathfrak C_{\mathrm{lock}}^{\mathrm{loc}}(j+1)\bigr),
    \bigl(\lambda_{\mathrm{lock}}^{\Theta,+}(j+1)-\lambda_{\mathrm{lock}}^{\Theta,+}(j)\bigr)_{+},
    \chi_{\mathrm{stop}}(j+1)
  \Bigr\}.
\]
Here $d_{\mathfrak C}$ denotes corridor distance in the current local readout, $(\cdot)_+$ denotes the positive part, and $\chi_{\mathrm{stop}}$ is the indicator of an inadmissible boundary-stop violation.  The defect vanishes exactly when the transported step remains corridor-inside, does not re-open lock leakage, and does not violate the inherited stop discipline.
\end{definition}

\begin{lemma}[Zero HC transport defect preserves capture]
\label{lem:zero-hc-transport-defect-preserves-capture}
Assume that $\mathcal C_n^{\mathrm{HC}}$ is a first local HC-capture chain and that
\[
  \Delta_{\mathrm{tr}}^{\mathrm{HC}}(j)=0
\]
for an admissible route-faithful transport step $\mathsf T_j$.  Then the transported locked state is again HC-captured in the transported local lock corridor.
\end{lemma}

\begin{proof}
The equality $\Delta_{\mathrm{tr}}^{\mathrm{HC}}(j)=0$ simultaneously says that the transported locked state has zero corridor distance from the next transported lock corridor, that lock-side leakage has not been positively re-opened, and that no boundary-stop violation has occurred.  These are exactly the transported versions of corridor membership, lock-side capture, and route admissibility in Definition~\ref{def:transport-stable-hc-capture-chain}.  Hence the transported locked state remains HC-captured.
\end{proof}

\begin{definition}[HC-to-OP handoff packet]
\label{def:hc-to-op-handoff-packet}
An \emph{HC-to-OP handoff packet} on $[n,n+k]$ is a tuple
\[
  \mathcal P_{\mathrm{HC}\to\mathrm{OP}}(n,k)
  :=
  \bigl(
    \mathcal C_n^{\mathrm{HC}},
    \{\mathsf T_j\}_{j=n}^{n+k-1},
    \{\Delta_{\mathrm{tr}}^{\mathrm{HC}}(j)\}_{j=n}^{n+k-1},
    \widetilde{\mathbf R}_{\mathrm{live}}(n+k),
    \tau_{\mathrm{route}}
  \bigr),
\]
where:
\begin{enumerate}[leftmargin=2em,label=(\alph*)]
  \item $\mathcal C_n^{\mathrm{HC}}$ is a first local HC-capture chain;
  \item the transports $\mathsf T_j$ are route-faithful and admissible;
  \item all displayed HC transport defects vanish on the window;
  \item the terminal normalized live residual $\widetilde{\mathbf R}_{\mathrm{live}}(n+k)$ lies in the admissible OP-entry normalization regime;
  \item $\tau_{\mathrm{route}}$ records whether the packet is horizon-side or stop-preserving boundary-side.
\end{enumerate}
The handoff packet is therefore the smallest object that can be consumed by the OP-facing machine without forgetting how the local HC capture was transported.
\end{definition}

\begin{proposition}[Handoff packets are admissible pre-OP frontier objects]
\label{prop:handoff-packets-are-admissible-pre-op-frontier-objects}
Every HC-to-OP handoff packet is an admissible pre-OP frontier object in the sense of the current formal OP-entry framework.
\end{proposition}

\begin{proof}
By definition, the packet contains a valid local HC-capture chain, an admissible route-faithful transport family, vanishing HC transport defects, an admissibly normalized terminal live residual, and an explicit route tag.  These are precisely the data required to prevent three forbidden moves: erasing the route distinction, re-opening a boundary stop, or treating local capture as OP closure without transport.  Hence the packet is admissible as a pre-OP frontier object.
\end{proof}

\begin{theorem}[First transport-stable HC-capture to OP-entry criterion]
\label{thm:first-transport-stable-hc-capture-to-op-entry-criterion}
Assume:
\begin{enumerate}[leftmargin=2em,label=(\alph*)]
  \item the first local HC-capture chain theorem applies;
  \item the captured chain is transport-stable on a finite admissible window $[n,n+k]$;
  \item the terminal normalized live residual $\widetilde{\mathbf R}_{\mathrm{live}}(n+k)$ lies in the admissible OP-entry normalization regime;
  \item the inherited route tag is preserved throughout the handoff.
\end{enumerate}
Then the current post-v3.33.0 line supports a controlled OP-entry handoff packet
\[
  \mathcal P_{\mathrm{HC}\to\mathrm{OP}}(n,k).
\]
If, in addition, the already-defined closure-phase trigger persistence condition holds at the terminal state, the packet yields an immediate OP-entry certificate.  If that persistence condition is not yet met, the same packet yields a narrowed frontier-certified pre-entry state rather than a false OP closure claim.
\end{theorem}

\begin{proof}
By the first local HC-capture chain theorem, the local repaired second step can be locked into the same local corridor as the first near-trigger step.  By transport-stability, Lemma~\ref{lem:zero-hc-transport-defect-preserves-capture} applies at each admissible transport step in the finite window.  Therefore the captured state remains captured after the finite route-faithful transport sequence, with no positive re-opening of lock-side leakage and no boundary-stop violation.  The normalized terminal residual and route tag then supply exactly the remaining fields in Definition~\ref{def:hc-to-op-handoff-packet}.  The final alternative follows from the existing OP-entry framework: trigger persistence upgrades the handoff packet to an entry certificate, while absence of trigger persistence leaves a legitimate, narrowed pre-entry frontier state.
\end{proof}

\begin{corollary}[v3.34.0-r01 working rule]
\label{cor:v3340-r01-working-rule}
The next audit should no longer ask merely whether a local HC capture chain exists.  It should test whether the captured chain is transport-stable and whether the resulting handoff packet terminates in an immediate OP-entry certificate or in a smaller frontier-certified pre-entry state.
\end{corollary}

\begin{proof}
Immediate from Theorem~\ref{thm:first-transport-stable-hc-capture-to-op-entry-criterion}.
\end{proof}



\subsection{First OP-entry normal form extracted from HC handoff}
\label{subsec:first-op-entry-normal-form-extracted-from-hc-handoff}

\begin{definition}[HC-extracted OP-entry normal form]
\label{def:hc-extracted-op-entry-normal-form}
Let $\mathcal P_{\mathrm{HC}\to\mathrm{OP}}(n,k)$ be an HC-to-OP handoff packet in the sense of Definition~\ref{def:hc-to-op-handoff-packet}.  Its \emph{HC-extracted OP-entry normal form} is the tuple
\[
  \mathcal N_{\mathrm{OP}}^{\mathrm{HC}}(n,k)
  :=
  \bigl(
    \widetilde{\mathbf R}_{\mathrm{live}}^{\mathrm{OP}}(n+k),
    \tau_{\mathrm{route}},
    \sigma_{\mathrm{tr}},
    \mathcal M_{\mathrm{OP}}^{\mathrm{adm}}(n+k),
    \eta_{\mathrm{entry}}
  \bigr),
\]
where:
\begin{enumerate}[leftmargin=2em,label=(\alph*)]
  \item $\widetilde{\mathbf R}_{\mathrm{live}}^{\mathrm{OP}}(n+k)$ is the terminal normalized live residual read as an OP-facing residual triple
  \[
    \widetilde{\mathbf R}_{\mathrm{live}}^{\mathrm{OP}}
    =
    \bigl(\rho_{\mathrm{CL}},\rho_{\Theta R},\rho_A\bigr);
  \]
  \item $\tau_{\mathrm{route}}$ is the inherited horizon-side or stop-preserving boundary-side route tag;
  \item $\sigma_{\mathrm{tr}}$ records that the HC transport defects vanish on the handoff window;
  \item $\mathcal M_{\mathrm{OP}}^{\mathrm{adm}}(n+k)$ is the finite admissible OP-facing request menu available at the terminal state;
  \item $\eta_{\mathrm{entry}}\in\{\mathrm{entry},\mathrm{frontier}\}$ records whether the terminal closure-phase trigger persistence condition has already upgraded the packet to immediate OP-entry or only to a smaller frontier-certified pre-entry state.
\end{enumerate}
The normal form is therefore not a new theorem layer.  It is the smallest downstream-readable form of a transported HC capture event.
\end{definition}

\begin{definition}[Residual routing out of the OP-entry normal form]
\label{def:residual-routing-out-of-op-entry-normal-form}
The residual routing map associated with $\mathcal N_{\mathrm{OP}}^{\mathrm{HC}}$ is the finite assignment
\[
  \mathsf{Route}_{\mathrm{OP}}
  \bigl(\rho_{\mathrm{CL}},\rho_{\Theta R},\rho_A;\tau_{\mathrm{route}}\bigr)
  =
  \bigl(
    \mathcal Q_{1/4},
    \mathcal Q_{5B},
    \mathcal Q_A
  \bigr),
\]
where:
\begin{enumerate}[leftmargin=2em,label=(\alph*)]
  \item $\mathcal Q_{1/4}$ contains only continuum-limit / Hopf-correction requests controlled by $\rho_{\mathrm{CL}}$ and is the admissible intake for the later OP~1/OP~4 correction line;
  \item $\mathcal Q_{5B}$ contains only routed $\Theta$--$R$ balance requests controlled by $\rho_{\Theta R}$ and is the admissible intake for the later OP~5(B) line;
  \item $\mathcal Q_A$ contains only anchor, admissibility, and route-discipline requests controlled by $\rho_A$;
  \item no request may erase $\tau_{\mathrm{route}}$, deny the earliest-boundary stop, or promote a diagnostic residual to a theorematic closure certificate by naming alone.
\end{enumerate}
\end{definition}

\begin{lemma}[Normal-form extraction preserves route discipline]
\label{lem:normal-form-extraction-preserves-route-discipline}
If $\mathcal P_{\mathrm{HC}\to\mathrm{OP}}(n,k)$ is an HC-to-OP handoff packet, then its HC-extracted OP-entry normal form preserves the inherited route distinction and boundary-stop discipline.
\end{lemma}

\begin{proof}
The handoff packet already contains the route tag, the transported HC-capture chain, the vanishing HC transport defects, and the terminal normalized live residual.  Definition~\ref{def:hc-extracted-op-entry-normal-form} only repackages these data into the fields consumed by the OP-facing request layer.  Since neither the route tag nor the stop indicator is discarded during that repackaging, the inherited route distinction and boundary-stop discipline are preserved.
\end{proof}

\begin{proposition}[Finite OP-facing request menu from HC handoff]
\label{prop:finite-op-facing-request-menu-from-hc-handoff}
Every HC-extracted OP-entry normal form determines a finite admissible OP-facing request menu
\[
  \mathcal M_{\mathrm{OP}}^{\mathrm{adm}}
  =
  \mathcal Q_{1/4}\cup \mathcal Q_{5B}\cup \mathcal Q_A
\]
subordinate to the inherited route tag.
\end{proposition}

\begin{proof}
The terminal live residual triple has only three displayed coordinates: continuum-limit, routed $\Theta$--$R$, and anchor/admissibility.  Definition~\ref{def:residual-routing-out-of-op-entry-normal-form} assigns each coordinate to one finite request family and forbids requests that break route discipline or manufacture a closure certificate.  Hence the union is finite, admissible, and subordinate to the inherited route tag.
\end{proof}

\begin{definition}[Lead OP-entry residual selector]
\label{def:lead-op-entry-residual-selector}
A \emph{lead OP-entry residual selector} is any admissible rule
\[
  \mathsf L_{\mathrm{OP}}
  \bigl(\mathcal N_{\mathrm{OP}}^{\mathrm{HC}}\bigr)
  \in
  \{\rho_{\mathrm{CL}},\rho_{\Theta R},\rho_A,\rho_{\mathrm{CL}}\oplus\rho_{\Theta R}\}
\]
that chooses the next finite OP-facing work branch according to the current residual ranking, subject to the following restrictions:
\begin{enumerate}[leftmargin=2em,label=(\alph*)]
  \item if $\rho_{\mathrm{CL}}$ is dominant, the next branch is the OP~1/OP~4 continuum-limit correction line;
  \item if $\rho_{\Theta R}$ is dominant, the next branch is the OP~5(B) routed balance line;
  \item if $\rho_A$ is dominant, the next branch is an admissibility/anchor refinement rather than a physical closure claim;
  \item if $\rho_{\mathrm{CL}}$ and $\rho_{\Theta R}$ are inseparable under the current ranking, the next branch is the paired OP~1/4--OP~5(B) comparison branch.
\end{enumerate}
\end{definition}

\begin{theorem}[First OP-entry normal-form theorem]
\label{thm:first-op-entry-normal-form-theorem}
Assume the first transport-stable HC-capture to OP-entry criterion of Theorem~\ref{thm:first-transport-stable-hc-capture-to-op-entry-criterion}.  Then the current v3.34.0 working line admits a unique downstream-readable OP-entry normal form
\[
  \mathcal N_{\mathrm{OP}}^{\mathrm{HC}}(n,k)
\]
for the transported HC capture event.  Moreover:
\begin{enumerate}[leftmargin=2em,label=(\roman*)]
  \item if $\eta_{\mathrm{entry}}=\mathrm{entry}$, the normal form is an immediate OP-entry surface;
  \item if $\eta_{\mathrm{entry}}=\mathrm{frontier}$, the normal form is a smaller frontier-certified pre-entry surface;
  \item in both cases the admissible next work is restricted to the residual-routed menu $\mathcal Q_{1/4}$, $\mathcal Q_{5B}$, and $\mathcal Q_A$.
\end{enumerate}
\end{theorem}

\begin{proof}
Theorem~\ref{thm:first-transport-stable-hc-capture-to-op-entry-criterion} supplies the controlled handoff packet or, if trigger persistence is absent, the narrowed frontier-certified pre-entry state carried by the same packet.  Definition~\ref{def:hc-extracted-op-entry-normal-form} reads exactly those fields into one terminal residual triple, one route tag, one transport-stability record, one admissible request menu, and one entry/frontier flag.  No additional choice is made during this extraction, so the downstream-readable normal form is unique relative to the given handoff packet.  The final restriction follows from Proposition~\ref{prop:finite-op-facing-request-menu-from-hc-handoff}.
\end{proof}

\begin{corollary}[v3.34.0-r02 working rule]
\label{cor:v3340-r02-working-rule}
The next mathematical audit should select the lead OP-entry residual from $\mathcal N_{\mathrm{OP}}^{\mathrm{HC}}$.  If the continuum-limit residual dominates, start with the OP~1/OP~4 correction branch.  If the routed $\Theta$--$R$ residual dominates, start with OP~5(B).  If neither is dominant, use the paired comparison branch rather than forcing a premature single-problem split.
\end{corollary}

\begin{proof}
Immediate from Definition~\ref{def:lead-op-entry-residual-selector} and Theorem~\ref{thm:first-op-entry-normal-form-theorem}.
\end{proof}

\subsection{First OP-entry task routing extracted from the normal form}
\label{subsec:first-op-entry-task-routing-extracted-from-normal-form}

\begin{definition}[OP-entry task-routing object]
\label{def:op-entry-task-routing-object}
Let
\[
  \mathcal N_{\mathrm{OP}}^{\mathrm{HC}}(n,k)
\]
be an HC-extracted OP-entry normal form in the sense of Definition~\ref{def:hc-extracted-op-entry-normal-form}, and let
\[
  \mathsf L_{\mathrm{OP}}
  \bigl(\mathcal N_{\mathrm{OP}}^{\mathrm{HC}}\bigr)
\]
be a lead OP-entry residual selector in the sense of Definition~\ref{def:lead-op-entry-residual-selector}.  The associated \emph{OP-entry task-routing object} is the tuple
\[
  \mathcal T_{\mathrm{OP}}^{\mathrm{HC}}(n,k)
  :=
  \bigl(
    \mathcal N_{\mathrm{OP}}^{\mathrm{HC}}(n,k),
    \mathsf L_{\mathrm{OP}},
    q_{\mathrm{lead}},
    \mathcal W_{\mathrm{OP}}^{\mathrm{lead}},
    \tau_{\mathrm{route}},
    \eta_{\mathrm{entry}}
  \bigr),
\]
where:
\begin{enumerate}[leftmargin=2em,label=(\alph*)]
  \item $q_{\mathrm{lead}}$ is the selected residual branch in
  \[
    \{\rho_{\mathrm{CL}},\rho_{\Theta R},\rho_A,
      \rho_{\mathrm{CL}}\oplus\rho_{\Theta R}\};
  \]
  \item $\mathcal W_{\mathrm{OP}}^{\mathrm{lead}}$ is the finite route-faithful work family assigned to that branch;
  \item $\tau_{\mathrm{route}}$ is inherited unchanged from the normal form;
  \item $\eta_{\mathrm{entry}}$ records whether the object is consumed as immediate OP-entry or only as a frontier-certified sharpening task.
\end{enumerate}
Thus the task-routing object does not decide an OP claim.  It decides only which finite admissible work family is allowed to be executed next.
\end{definition}

\begin{definition}[Residual-to-task routing map]
\label{def:residual-to-task-routing-map}
The \emph{residual-to-task routing map} associated with an OP-entry normal form is the admissible assignment
\[
  \mathfrak R_{\mathrm{task}}
  :
  \{\rho_{\mathrm{CL}},\rho_{\Theta R},\rho_A,
      \rho_{\mathrm{CL}}\oplus\rho_{\Theta R}\}
  \longrightarrow
  \{\mathcal W_{1/4}^{\mathrm{CL}},
    \mathcal W_{5B}^{\Theta R},
    \mathcal W_A^{\mathrm{adm}},
    \mathcal W_{1/4\mid 5B}^{\mathrm{cmp}}
  \}
\]
defined by
\[
\begin{array}{rcl}
  \rho_{\mathrm{CL}} &\longmapsto&
  \mathcal W_{1/4}^{\mathrm{CL}},\\[0.25em]
  \rho_{\Theta R} &\longmapsto&
  \mathcal W_{5B}^{\Theta R},\\[0.25em]
  \rho_A &\longmapsto&
  \mathcal W_A^{\mathrm{adm}},\\[0.25em]
  \rho_{\mathrm{CL}}\oplus\rho_{\Theta R} &\longmapsto&
  \mathcal W_{1/4\mid 5B}^{\mathrm{cmp}}.
\end{array}
\]
Here $\mathcal W_{1/4}^{\mathrm{CL}}$ contains only continuum-limit correction tasks for OP~1/OP~4, $\mathcal W_{5B}^{\Theta R}$ contains only routed $\Theta$--$R$ balance tasks for OP~5(B), $\mathcal W_A^{\mathrm{adm}}$ contains only anchor/admissibility refinement tasks, and $\mathcal W_{1/4\mid 5B}^{\mathrm{cmp}}$ contains paired comparison tasks that keep the common HC-capture origin visible when the two leading residuals cannot yet be separated.
\end{definition}

\begin{lemma}[Task routing preserves OP-entry normal-form discipline]
\label{lem:task-routing-preserves-op-entry-normal-form-discipline}
Let $\mathcal T_{\mathrm{OP}}^{\mathrm{HC}}(n,k)$ be an OP-entry task-routing object.  Then executing its selected finite work family cannot erase the inherited route tag, cannot remove the boundary-stop discipline, and cannot manufacture a new OP closure certificate.
\end{lemma}

\begin{proof}
By Definition~\ref{def:op-entry-task-routing-object}, the task-routing object carries the normal form, the selected residual branch, the finite work family, the inherited route tag, and the entry/frontier flag as explicit fields.  By Definition~\ref{def:residual-to-task-routing-map}, each work family is subordinate to one already displayed residual coordinate or to the paired comparison of the two leading displayed coordinates.  None of these assignments changes the route tag or the entry/frontier flag.  Hence task execution may refine, compare, or package the selected branch, but it cannot promote itself to a closure certificate unless the pre-existing OP-entry framework supplies the missing trigger persistence.
\end{proof}

\begin{proposition}[Paired comparison branch is the safe non-dominant branch]
\label{prop:paired-comparison-branch-safe-non-dominant-branch}
If the current residual ranking cannot separate $\rho_{\mathrm{CL}}$ from $\rho_{\Theta R}$, the admissible next branch is the paired comparison family
\[
  \mathcal W_{1/4\mid 5B}^{\mathrm{cmp}}
\]
rather than a forced single OP~1/OP~4 or OP~5(B) branch.
\end{proposition}

\begin{proof}
Definition~\ref{def:lead-op-entry-residual-selector} already allows the paired symbol $\rho_{\mathrm{CL}}\oplus\rho_{\Theta R}$ exactly when the two residuals are inseparable under the current ranking.  Definition~\ref{def:residual-to-task-routing-map} sends that symbol to the paired comparison family.  Choosing a single branch in that case would add information not present in the normal form and could hide the common HC-capture origin of both residuals.  The paired branch is therefore the route-faithful and conservative choice.
\end{proof}

\begin{theorem}[First OP-entry task-routing theorem]
\label{thm:first-op-entry-task-routing-theorem}
Assume that an HC-to-OP handoff packet has been converted into the HC-extracted OP-entry normal form
\[
  \mathcal N_{\mathrm{OP}}^{\mathrm{HC}}(n,k)
\]
of Theorem~\ref{thm:first-op-entry-normal-form-theorem}.  Then every admissible lead selector $\mathsf L_{\mathrm{OP}}$ determines a finite route-faithful OP-entry task-routing object
\[
  \mathcal T_{\mathrm{OP}}^{\mathrm{HC}}(n,k).
\]
Moreover:
\begin{enumerate}[leftmargin=2em,label=(\roman*)]
  \item if $\eta_{\mathrm{entry}}=\mathrm{entry}$, the selected family is a first admissible OP-facing work packet;
  \item if $\eta_{\mathrm{entry}}=\mathrm{frontier}$, the selected family is only a finite sharpening/comparison packet for the smaller frontier-certified pre-entry state;
  \item if the selector returns $\rho_{\mathrm{CL}}\oplus\rho_{\Theta R}$, the next work packet must keep OP~1/OP~4 and OP~5(B) in paired comparison until a later ranking separates them.
\end{enumerate}
\end{theorem}

\begin{proof}
The normal form contains a finite residual triple, a finite admissible request menu, a route tag, a transport-stability record, and an entry/frontier flag.  The lead selector chooses one of the explicitly allowed residual branches.  Applying Definition~\ref{def:residual-to-task-routing-map} to that branch yields one finite work family, and adjoining the inherited normal-form fields gives the task-routing object of Definition~\ref{def:op-entry-task-routing-object}.  Lemma~\ref{lem:task-routing-preserves-op-entry-normal-form-discipline} gives the route and certificate restrictions.  Proposition~\ref{prop:paired-comparison-branch-safe-non-dominant-branch} gives the paired-branch rule when the two leading residuals are not yet separated.
\end{proof}

\begin{corollary}[v3.34.0-r03 working rule]
\label{cor:v3340-r03-working-rule}
The next OP-facing audit should not yet try to close OP~1/OP~4 or OP~5(B) directly.  It should first instantiate the task-routing object $\mathcal T_{\mathrm{OP}}^{\mathrm{HC}}(n,k)$, rank the displayed residuals, and execute only the selected finite work family: $\mathcal W_{1/4}^{\mathrm{CL}}$, $\mathcal W_{5B}^{\Theta R}$, $\mathcal W_A^{\mathrm{adm}}$, or the paired comparison family $\mathcal W_{1/4\mid 5B}^{\mathrm{cmp}}$.
\end{corollary}

\begin{proof}
Immediate from Theorem~\ref{thm:first-op-entry-task-routing-theorem}.  The task-routing object is the conservative bridge between OP-entry normal form and explicit OP-facing execution.
\end{proof}


\subsection{First paired OP-entry comparison execution layer}
\label{subsec:first-paired-op-entry-comparison-execution-layer}

\begin{definition}[Paired OP-entry comparison execution packet]
\label{def:paired-op-entry-comparison-execution-packet}
Assume that the OP-entry task-routing object of Definition~\ref{def:op-entry-task-routing-object} selects the paired branch
\[
  q_{\mathrm{lead}}=\rho_{\mathrm{CL}}\oplus\rho_{\Theta R},
  \qquad
  \mathcal W_{\mathrm{OP}}^{\mathrm{lead}}
  =\mathcal W_{1/4\mid 5B}^{\mathrm{cmp}}.
\]
A \emph{paired OP-entry comparison execution packet} is the tuple
\[
  \mathcal X_{1/4\mid 5B}^{\mathrm{cmp}}(n,k)
  :=
  \bigl(
    \mathcal T_{\mathrm{OP}}^{\mathrm{HC}}(n,k),
    \mathcal U_{\mathrm{HC}}^{\mathrm{adm}},
    D_{\mathrm{joint}}^{\mathrm{HC}},
    D_{\mathrm{sep}}^{\mathrm{HC}},
    \Delta_{\mathrm{sep}}^{\mathrm{HC}},
    \omega_{\mathrm{cmp}}
  \bigr),
\]
where:
\begin{enumerate}[leftmargin=2em,label=(\alph*)]
  \item $\mathcal T_{\mathrm{OP}}^{\mathrm{HC}}(n,k)$ is the inherited task-routing object;
  \item $\mathcal U_{\mathrm{HC}}^{\mathrm{adm}}$ is the finite admissible menu of HC-compatible micro-updates available at the terminal normal form;
  \item $D_{\mathrm{joint}}^{\mathrm{HC}}$ is the best joint readout error obtained when one and the same admissible HC-compatible update is asked to account for both $\rho_{\mathrm{CL}}$ and $\rho_{\Theta R}$;
  \item $D_{\mathrm{sep}}^{\mathrm{HC}}$ is the best separated readout error obtained when the two coordinates are allowed to choose their admissible update witnesses independently;
  \item
  \[
    \Delta_{\mathrm{sep}}^{\mathrm{HC}}
    :=D_{\mathrm{joint}}^{\mathrm{HC}}-D_{\mathrm{sep}}^{\mathrm{HC}}
    \geq 0
  \]
  is the HC-origin separation defect;
  \item $\omega_{\mathrm{cmp}}\in\{\mathrm{bind},\mathrm{split},\mathrm{frontier}\}$ records the comparison outcome.
\end{enumerate}
The packet executes only the comparison branch.  It does not yet assert an OP~1/4 correction theorem and does not yet assert the routed OP~5(B) balance theorem.
\end{definition}

\begin{definition}[HC-origin comparison decision rule]
\label{def:hc-origin-comparison-decision-rule}
Fix an admissible comparison tolerance $\varepsilon_{\mathrm{cmp}}>0$ and a strict descent requirement on the active frontier measure.  The outcome field of Definition~\ref{def:paired-op-entry-comparison-execution-packet} is read by
\[
  \omega_{\mathrm{cmp}}
  =
  \begin{cases}
    \mathrm{bind},
      & \Delta_{\mathrm{sep}}^{\mathrm{HC}}\leq \varepsilon_{\mathrm{cmp}},\\[0.25em]
    \mathrm{split},
      & \Delta_{\mathrm{sep}}^{\mathrm{HC}}>\varepsilon_{\mathrm{cmp}}
        \text{ and the separated ranking is stable},\\[0.25em]
    \mathrm{frontier},
      & \text{otherwise.}
  \end{cases}
\]
The bind case means that the two residual coordinates are still best treated as one common HC-origin comparison problem.  The split case means that the comparison has supplied enough information to return to a single OP~1/4 or OP~5(B) branch under the residual ranking.  The frontier case means that the comparison improved the state, but not enough to justify either binding or splitting as a theorematic route decision.
\end{definition}

\begin{lemma}[Paired comparison execution preserves the inherited handoff discipline]
\label{lem:paired-comparison-execution-preserves-handoff-discipline}
Executing a paired OP-entry comparison packet cannot erase the inherited route tag, cannot reopen the boundary-stop condition, and cannot promote the current frontier state to an OP closure certificate.
\end{lemma}

\begin{proof}
The execution packet contains the task-routing object as an explicit first component.  By Lemma~\ref{lem:task-routing-preserves-op-entry-normal-form-discipline}, that object already carries the inherited route tag, entry/frontier flag, and admissible work-family restriction.  The additional fields in Definition~\ref{def:paired-op-entry-comparison-execution-packet} compare only joint versus separated HC-compatible readout errors inside the finite admissible micro-update menu.  They do not change the route tag, do not delete the boundary-stop flag, and do not introduce a trigger-persistence certificate.  Hence the comparison can only bind, split, or sharpen the frontier; it cannot close an OP by itself.
\end{proof}

\begin{proposition}[Bind branch retains the common HC-capture origin]
\label{prop:bind-branch-retains-common-hc-capture-origin}
If $\omega_{\mathrm{cmp}}=\mathrm{bind}$, then OP~1/4 and OP~5(B) must remain in the coupled work family
\[
  \mathcal W_{1/4\mid 5B}^{\mathrm{cmp}}
\]
for the next step, with the common HC-capture origin kept as an explicit field.
\end{proposition}

\begin{proof}
In the bind case, Definition~\ref{def:hc-origin-comparison-decision-rule} says that the best joint HC-compatible readout is not worse than the separated readout by more than the admissible comparison tolerance.  Splitting the branch at that point would artificially discard the information that one HC-compatible source still accounts for both leading residuals.  The route-faithful next object therefore remains the paired comparison family, now with the common HC-origin field recorded rather than hidden.
\end{proof}

\begin{proposition}[Split branch returns to ranked single-branch work]
\label{prop:split-branch-returns-to-ranked-single-branch-work}
If $\omega_{\mathrm{cmp}}=\mathrm{split}$, then the comparison packet returns a ranked single-branch task-routing object selecting either
\[
  \mathcal W_{1/4}^{\mathrm{CL}}
  \qquad\text{or}\qquad
  \mathcal W_{5B}^{\Theta R},
\]
according to the stable separated residual ranking.
\end{proposition}

\begin{proof}
The split case is allowed only when the HC-origin separation defect exceeds the admissible comparison tolerance and the separated ranking is stable.  Thus the normal form now contains information that was absent in the inseparable r03 branch: one of the two residual coordinates has become a legitimate lead branch.  Applying the residual-to-task routing map of Definition~\ref{def:residual-to-task-routing-map} to that newly ranked coordinate gives either the OP~1/4 continuum-limit work family or the OP~5(B) routed $\Theta$--$R$ work family.
\end{proof}

\begin{theorem}[First paired OP-entry comparison execution theorem]
\label{thm:first-paired-op-entry-comparison-execution-theorem}
Assume that the r03 task-routing layer selects the paired comparison branch
\[
  \rho_{\mathrm{CL}}\oplus\rho_{\Theta R}
  \longmapsto
  \mathcal W_{1/4\mid 5B}^{\mathrm{cmp}}.
\]
Then every admissible finite HC-compatible comparison menu produces one paired OP-entry comparison execution packet
\[
  \mathcal X_{1/4\mid 5B}^{\mathrm{cmp}}(n,k)
\]
and exactly one of the following route-faithful next states:
\begin{enumerate}[leftmargin=2em,label=(\roman*)]
  \item a bound common-origin comparison state, still subordinate to $\mathcal W_{1/4\mid 5B}^{\mathrm{cmp}}$;
  \item a split ranked single-branch state, subordinate to either $\mathcal W_{1/4}^{\mathrm{CL}}$ or $\mathcal W_{5B}^{\Theta R}$;
  \item a smaller frontier-certified comparison state, if neither binding nor splitting is yet justified.
\end{enumerate}
In all three cases, no OP closure certificate is manufactured by the comparison execution itself.
\end{theorem}

\begin{proof}
The finite admissible menu $\mathcal U_{\mathrm{HC}}^{\mathrm{adm}}$ supplies the joint and separated minimization tests by construction.  Their difference defines $\Delta_{\mathrm{sep}}^{\mathrm{HC}}\geq0$, since the separated readout is at least as flexible as the joint readout.  Definition~\ref{def:hc-origin-comparison-decision-rule} assigns exactly one of the three comparison outcomes.  Proposition~\ref{prop:bind-branch-retains-common-hc-capture-origin} gives the bound next state; Proposition~\ref{prop:split-branch-returns-to-ranked-single-branch-work} gives the split next state; and the remaining case is, by definition, the smaller frontier-certified comparison state.  Lemma~\ref{lem:paired-comparison-execution-preserves-handoff-discipline} supplies the route, boundary, and no-false-closure restrictions.
\end{proof}

\begin{corollary}[v3.34.0-r04 working rule]
\label{cor:v3340-r04-working-rule}
The next OP-facing audit may now execute the paired comparison branch before choosing between OP~1/4 and OP~5(B).  The correct next question is not yet \emph{which OP closes first}, but whether $\rho_{\mathrm{CL}}$ and $\rho_{\Theta R}$ are still one common HC-origin residual or have separated into stable ranked branches.
\end{corollary}

\begin{proof}
Immediate from Theorem~\ref{thm:first-paired-op-entry-comparison-execution-theorem}.
\end{proof}



\subsection{Minimal MML/MMA interface for HC-routed OP execution}
\label{subsec:minimal-mml-mma-interface-hc-routed-op-execution}

\begin{definition}[r07 local language-version registry]
\label{def:r07-local-language-version-registry}
The current HC-routed OP-entry execution layer uses the following local interface tags; their interface stamps are recorded in prose, not as mathematical exponents:
\[
  \mathrm{QAM}_{\mathrm{HCOP}}^{\mathrm{loc}},
  \qquad
  \mathrm{MML}_{\mathrm{HCOP}}^{\mathrm{loc}},
  \qquad
  \mathrm{MMA}_{\mathrm{HCOP}}^{\mathrm{loc}}.
\]
Their status is intentionally limited.
\begin{enumerate}[leftmargin=2em,label=(\alph*)]
  \item \(\mathrm{QAM}_{\mathrm{HCOP}}^{\mathrm{loc}}\) is the finite tagged QAM bridge used to code the HC/OP objects of the present subsection inside the inherited QAM-over-ZF layer; its interface stamp is QAM-HC/OP 0.1.
  \item \(\mathrm{MML}_{\mathrm{HCOP}}^{\mathrm{loc}}\) is the finite instruction language for the r01--r04 HC-routed OP-entry chain; its interface stamp is MML-HC/OP 0.2.
  \item \(\mathrm{MMA}_{\mathrm{HCOP}}^{\mathrm{loc}}\) is the matrix/readout execution semantics for those instructions on finite route-faithful machine states; its interface stamp is MMA-HC/OP 0.2.
\end{enumerate}
This registry upgrades only the local execution interface. It does not replace the inherited public QAM layer, does not activate the deferred binary tag system, and does not turn the historical MML--MMA--MCP reserve into an independent proof engine.
\end{definition}

\begin{definition}[HC/OP MMA object types]
\label{def:hcop-mma-object-types}
The minimal MMA object layer for HC-routed OP execution consists of the finite type family
\[
\begin{aligned}
  \mathsf{HCState},\quad
  \mathsf{HandoffPacket},\quad
  \mathsf{OPNormalForm},\quad
  \mathsf{TaskRouting},\quad
  \mathsf{ComparisonPacket},\\
  \mathsf{Decision},\quad
  \mathsf{FrontierState},\quad
  \mathsf{RouteTag},\quad
  \mathsf{CertificateCandidate}.
\end{aligned}
\]
The intended instances are respectively: local HC-capture states, the handoff packets of Definition~\ref{def:hc-to-op-handoff-packet}, the normal forms of Definition~\ref{def:hc-extracted-op-entry-normal-form}, the routing objects of Definition~\ref{def:op-entry-task-routing-object}, the comparison packets of Definition~\ref{def:paired-op-entry-comparison-execution-packet}, the decision values \(\mathrm{bind},\mathrm{split},\mathrm{frontier}\), smaller pre-entry frontier states, inherited route tags, and possible certificate data. A \(\mathsf{CertificateCandidate}\) becomes an actual theorematic certificate only when the previously stated mathematical trigger conditions are discharged outside the machine interface.
\end{definition}

\begin{definition}[QAM coding bridge for HC/OP objects]
\label{def:qam-coding-bridge-for-hcop-objects}
A \emph{QAM-HC/OP bridge code} is a partial coding map
\[
  \ulcorner\cdot\urcorner_{\mathrm{QAM}}^{\mathrm{HCOP}}
  :
  \mathsf{Obj}_{\mathrm{MMA}}^{\mathrm{HCOP}}
  \dashrightarrow
  \mathsf{QAM}_{\mathrm{fin}},
\]
where \(\mathsf{Obj}_{\mathrm{MMA}}^{\mathrm{HCOP}}\) is the finite object family of Definition~\ref{def:hcop-mma-object-types} and \(\mathsf{QAM}_{\mathrm{fin}}\) denotes finite tagged QAM objects in the inherited QAM-over-ZF sense. The map is required to preserve object kind, route tag, entry/frontier flag, and no-false-certificate status.
\end{definition}

\begin{definition}[Minimal MML-HC/OP instruction alphabet]
\label{def:minimal-mml-hcop-instruction-alphabet}
The local instruction alphabet \(\mathcal I_{\mathrm{MML}}^{\mathrm{HCOP}}\) consists of the following finite instructions:
\[
\begin{gathered}
  \mathrm{CAPTURE}_{\mathrm{HC}},\qquad
  \mathrm{HANDOFF}_{\mathrm{HC}\to\mathrm{OP}},\qquad
  \mathrm{NORMALIZE}_{\mathrm{OP}}^{\mathrm{HC}},\qquad
  \mathrm{ROUTE}_{\mathrm{OP}}^{\mathrm{HC}},\\
  \mathrm{COMPARE}_{\mathrm{HC}}^{1/4\mid 5B},\qquad
  \mathrm{RESOLVE}_{\mathrm{bind/split/frontier}},\qquad
  \mathrm{RETURN}_{\mathrm{OP}}.
\end{gathered}
\]
The instructions are typed as follows:
\begin{enumerate}[leftmargin=2em,label=(\alph*)]
  \item \(\mathrm{CAPTURE}_{\mathrm{HC}}\) sends a local admissible HC-state to a candidate transport-stable HC state.
  \item \(\mathrm{HANDOFF}_{\mathrm{HC}\to\mathrm{OP}}\) sends such a state to an HC-to-OP handoff packet.
  \item \(\mathrm{NORMALIZE}_{\mathrm{OP}}^{\mathrm{HC}}\) sends the handoff packet to an OP-entry normal form.
  \item \(\mathrm{ROUTE}_{\mathrm{OP}}^{\mathrm{HC}}\) sends the normal form to a finite OP-entry task-routing object.
  \item \(\mathrm{COMPARE}_{\mathrm{HC}}^{1/4\mid 5B}\) executes the paired OP~1/4--OP~5(B) comparison when the routing object selects \(\rho_{\mathrm{CL}}\oplus\rho_{\Theta R}\).
  \item \(\mathrm{RESOLVE}_{\mathrm{bind/split/frontier}}\) records exactly one comparison decision.
  \item \(\mathrm{RETURN}_{\mathrm{OP}}\) returns either a bound common-origin state, a split ranked single-branch state, or a smaller frontier-certified comparison state.
\end{enumerate}
\end{definition}

\begin{definition}[MMA execution state for the HC/OP program]
\label{def:mma-execution-state-for-hcop-program}
An \emph{MMA-HC/OP execution state} is a finite tuple
\[
  \Sigma_{\mathrm{HCOP}}
  =
  \bigl(
    S,
    \tau,
    \eta,
    r,
    c
  \bigr),
\]
where \(S\) is the finite QAM-coded store, \(\tau\) is the inherited route tag, \(\eta\in\{\mathrm{entry},\mathrm{frontier}\}\) is the entry/frontier flag, \(r\) is the active residual or comparison register, and \(c\) is a certificate register. The invariant of the state is:
\[
  c=\varnothing
  \quad\text{unless a previously proved mathematical trigger supplies a certificate.}
\]
\end{definition}

\begin{definition}[Finite HC-routed OP-entry MML program]
\label{def:finite-hc-routed-op-entry-mml-program}
The \emph{minimal HC-routed OP-entry program} is the finite MML word
\[
\begin{aligned}
  P_{\mathrm{HCOP}}
  :=
  (&\mathrm{CAPTURE}_{\mathrm{HC}},
    \mathrm{HANDOFF}_{\mathrm{HC}\to\mathrm{OP}},
    \mathrm{NORMALIZE}_{\mathrm{OP}}^{\mathrm{HC}},
    \mathrm{ROUTE}_{\mathrm{OP}}^{\mathrm{HC}},\\
   &\mathrm{COMPARE}_{\mathrm{HC}}^{1/4\mid 5B},
    \mathrm{RESOLVE}_{\mathrm{bind/split/frontier}},
    \mathrm{RETURN}_{\mathrm{OP}}).
\end{aligned}
\]
If the task-routing object does not select the paired branch, the comparison instruction is skipped and the return instruction hands control to the selected single work family instead. Thus the program is finite in either the paired or the single-branch case.
\end{definition}

\begin{lemma}[Typing of the r01--r04 HC-routed OP chain]
\label{lem:typing-of-r01-r04-hc-routed-op-chain}
Each step from the transport-stable HC-capture handoff through the first paired OP-entry comparison execution has a matching type in Definitions~\ref{def:hcop-mma-object-types}--\ref{def:minimal-mml-hcop-instruction-alphabet}.
\end{lemma}

\begin{proof}
The r01 handoff packet is an instance of \(\mathsf{HandoffPacket}\). The r02 normal form is an instance of \(\mathsf{OPNormalForm}\). The r03 task-routing object is an instance of \(\mathsf{TaskRouting}\). The r04 paired comparison packet is an instance of \(\mathsf{ComparisonPacket}\), and its outcome belongs to \(\mathsf{Decision}\). The possible return states are exactly the bound, split, or frontier states already listed in Theorem~\ref{thm:first-paired-op-entry-comparison-execution-theorem}. These are precisely the input/output types of the finite instruction alphabet.
\end{proof}

\begin{lemma}[MML/MMA execution preserves route and frontier discipline]
\label{lem:mml-mma-execution-preserves-route-and-frontier-discipline}
Executing \(P_{\mathrm{HCOP}}\) in the MMA-HC/OP layer preserves the inherited route tag, preserves the entry/frontier flag unless a previously proved trigger changes it, and cannot fill the certificate register by machine action alone.
\end{lemma}

\begin{proof}
Each instruction in Definition~\ref{def:minimal-mml-hcop-instruction-alphabet} is typed over the objects already constrained by the r01--r04 mathematical layer. The handoff and normalization instructions inherit the route tag from the HC-capture chain; the routing and comparison instructions act only on the finite residual and comparison registers; the resolution instruction chooses only among \(\mathrm{bind}\), \(\mathrm{split}\), and \(\mathrm{frontier}\); and the return instruction returns one of the route-faithful next states of Theorem~\ref{thm:first-paired-op-entry-comparison-execution-theorem}. None of these instructions contains a rule that writes a theorematic certificate into \(c\). Hence route, frontier, and no-false-certificate discipline are preserved.
\end{proof}

\begin{proposition}[r04 comparison as one admissible MML/MMA step]
\label{prop:r04-comparison-as-one-admissible-mml-mma-step}
The paired comparison execution of Theorem~\ref{thm:first-paired-op-entry-comparison-execution-theorem} is represented in \(\mathrm{MML}_{\mathrm{HCOP}}^{\mathrm{loc}}\) by the instruction
\[
  \mathrm{COMPARE}_{\mathrm{HC}}^{1/4\mid 5B}
\]
followed by \(\mathrm{RESOLVE}_{\mathrm{bind/split/frontier}}\). Its MMA semantics are exactly the construction of \(D_{\mathrm{joint}}^{\mathrm{HC}}\), \(D_{\mathrm{sep}}^{\mathrm{HC}}\), \(\Delta_{\mathrm{sep}}^{\mathrm{HC}}\), and the decision \(\omega_{\mathrm{cmp}}\).
\end{proposition}

\begin{proof}
Definition~\ref{def:paired-op-entry-comparison-execution-packet} constructs the joint and separated HC-compatible readout errors and their difference. Definition~\ref{def:hc-origin-comparison-decision-rule} resolves the resulting comparison defect into exactly one decision. These are exactly the operational contents assigned to the comparison and resolution instructions. The statement is therefore only a typed re-reading of the r04 comparison execution, not an additional theorematic claim.
\end{proof}

\begin{theorem}[HC-routed OP-entry as a finite MML/MMA program]
\label{thm:hc-routed-op-entry-as-finite-mml-mma-program}
Assume the transport-stable HC-capture to OP-entry criterion of Theorem~\ref{thm:first-transport-stable-hc-capture-to-op-entry-criterion}. Then the current r01--r04 HC-routed OP-entry chain is executable as the finite program \(P_{\mathrm{HCOP}}\) over the local interface stack
\[
  \mathrm{QAM}_{\mathrm{HCOP}}^{\mathrm{loc}}
  \longrightarrow
  \mathrm{MML}_{\mathrm{HCOP}}^{\mathrm{loc}}
  \longrightarrow
  \mathrm{MMA}_{\mathrm{HCOP}}^{\mathrm{loc}}.
\]
The execution returns exactly one admissible next state: immediate OP-entry when the inherited mathematical trigger supplies it, a bound common-origin comparison state, a split ranked single-branch state, or a smaller frontier-certified state. It does not manufacture an OP closure certificate.
\end{theorem}

\begin{proof}
By Lemma~\ref{lem:typing-of-r01-r04-hc-routed-op-chain}, every mathematical object used in the r01--r04 chain has a corresponding local MMA type and QAM bridge code. Definition~\ref{def:finite-hc-routed-op-entry-mml-program} lists the finite instruction word realizing the chain. Lemma~\ref{lem:mml-mma-execution-preserves-route-and-frontier-discipline} supplies route preservation and the no-false-certificate condition. The possible next states are exactly the ones already obtained in the handoff, normal-form, routing, and paired-comparison theorems. Therefore the MML/MMA layer is an executable notation for the existing mathematical chain, not a new source of closure.
\end{proof}

\begin{corollary}[v3.34.0-r07 working rule]
\label{cor:v3340-r07-working-rule}
The QAM/MML/MMA layer may now be used as a notation and audit interface for the HC-routed OP-entry program. Any later language extension must preserve the local registry of Definition~\ref{def:r07-local-language-version-registry}: QAM remains a finite bridge over ZF-coded objects, MML remains a finite instruction language, MMA remains the execution semantics, and theorematic certificates must still come from the mathematical trigger conditions rather than from the machine layer itself.
\end{corollary}

\begin{proof}
Immediate from Theorem~\ref{thm:hc-routed-op-entry-as-finite-mml-mma-program} and the status restriction in Definition~\ref{def:r07-local-language-version-registry}.
\end{proof}


\subsection{MML/MMA compression map for inherited HC/OP text}
\label{subsec:mml-mma-compression-map-for-inherited-hcop-text}

\begin{definition}[Inherited HC/OP compression target]
\label{def:inherited-hcop-compression-target}
An \emph{inherited HC/OP compression target} is a paragraph, lemma, proposition, diagnostic block, or tool paragraph in the inherited active or historical Companion text whose proof-use is one of the following finite roles:
\[
\begin{gathered}
  \mathrm{capture},\quad
  \mathrm{handoff},\quad
  \mathrm{normalization},\quad
  \mathrm{routing},\quad
  \mathrm{comparison},\quad
  \mathrm{resolution},\quad
  \mathrm{return},\\
  \mathrm{scan},\quad
  \mathrm{rank},\quad
  \mathrm{certificate\ test},\quad
  \mathrm{route\ preservation},\quad
  \mathrm{frontier\ descent}.
\end{gathered}
\]
The target is \emph{compressible} only in the expository sense: it may be reread as a typed MML/MMA operation or short finite operation word, but its original mathematical content, hypotheses, and proof restrictions remain part of the document unless a later explicit editorial release removes or relocates them.
\end{definition}

\begin{definition}[Compression map from inherited text to MML/MMA operations]
\label{def:compression-map-from-inherited-text-to-mml-mma-operations}
The local \emph{HC/OP compression map}
\[
  \mathfrak C_{\mathrm{HCOP}}
\]
sends inherited compression targets to typed MML/MMA readings according to the following table.

{\small
\renewcommand{\arraystretch}{1.18}
\begin{longtable}{|p{0.28\textwidth}|p{0.24\textwidth}|p{0.36\textwidth}|}
\hline
\textbf{Inherited text pattern} & \textbf{MML/MMA reading} & \textbf{Compression use} \\
\hline
HC lock corridor, local HC candidate, minimal HC-lock step &
\(\mathrm{LOCK}/\mathrm{CAPTURE}_{\mathrm{HC}}\) &
Marks the block as a candidate-production and lock-stabilization step rather than as an OP certificate. \\
\hline
Transport-stable HC-to-OP packet &
\(\mathrm{HANDOFF}_{\mathrm{HC}\to\mathrm{OP}}\) &
Turns repeated handoff prose into one typed transfer from HC state to OP-facing packet. \\
\hline
OP-entry normal form and residual triple &
\(\mathrm{NORMALIZE}_{\mathrm{OP}}^{\mathrm{HC}}\) &
Collects \(\rho_{\mathrm{CL}}\), \(\rho_{\Theta R}\), \(\rho_A\), route tags, and entry/frontier flags into one typed normal form. \\
\hline
Residual-to-task selector &
\(\mathrm{ROUTE}_{\mathrm{OP}}^{\mathrm{HC}}\) &
Replaces repeated branch-selection explanations by one admissible routing step. \\
\hline
Paired OP~1/4--OP~5(B) branch &
\(\mathrm{COMPARE}_{\mathrm{HC}}^{1/4\mid 5B}\) &
Reads the joint/separated defect comparison as one admissible comparison instruction. \\
\hline
Bind/split/frontier decision &
\(\mathrm{RESOLVE}_{\mathrm{bind/split/frontier}}\) &
Records that the comparison returns only a bound state, a split state, or a smaller frontier state. \\
\hline
Next-state handback to OP-facing work &
\(\mathrm{RETURN}_{\mathrm{OP}}\) &
Compresses the return package without changing theorematic status. \\
\hline
Probe, Monte Carlo, decoder, or finite search block &
\(\mathrm{SCAN}\to\mathrm{RANK}\to\mathrm{CANDIDATE}\) &
Makes clear that computational probes rank admissible candidates but do not prove closure. \\
\hline
Admissibility / boundary-stop / route-faithfulness paragraph &
\(\mathrm{PRESERVE\_ROUTE}\) &
Marks the paragraph as an invariant-preservation rule for route tag and boundary-stop data. \\
\hline
Certificate discussion or trigger-persistence paragraph &
\(\mathrm{CERTIFY?}\) &
Keeps certification separate from execution: the operation asks whether a mathematical trigger exists; it never creates one. \\
\hline
Frontier descent or smaller pre-entry state &
\(\mathrm{DESCEND}_{\mathrm{frontier}}\) &
Records strict controlled progress without falsely upgrading the state to OP closure. \\
\hline
Historical tool or early proof-attempt block &
\(\mathrm{ARCHIVE}\to\mathrm{TOOL}\) &
Keeps the historical witness intact while extracting the reusable tool-role for the active line. \\
\hline
\end{longtable}
}
\end{definition}

\begin{lemma}[Compression map preserves proof content]
\label{lem:compression-map-preserves-proof-content}
Applying \(\mathfrak C_{\mathrm{HCOP}}\) to an inherited compression target does not change the truth status of that target, does not delete hypotheses, and does not promote a diagnostic or work packet to a theorematic certificate.
\end{lemma}

\begin{proof}
By Definition~\ref{def:inherited-hcop-compression-target}, a compression target is compressed only in its expository role. Definition~\ref{def:compression-map-from-inherited-text-to-mml-mma-operations} assigns a typed MML/MMA reading that records what the block does in the finite HC/OP program. The assignment neither alters the original statement nor supplies a new mathematical trigger. The no-false-certificate restriction of Lemma~\ref{lem:mml-mma-execution-preserves-route-and-frontier-discipline} therefore remains in force.
\end{proof}

\begin{proposition}[Main inherited text zones helped by the compression map]
\label{prop:main-inherited-text-zones-helped-by-compression-map}
The compression map gives immediate expository leverage in the following inherited zones:
\begin{enumerate}[leftmargin=2em,label=(\roman*)]
  \item HC-capture and lock-corridor paragraphs, by replacing repeated candidate/lock explanations with \(\mathrm{LOCK}/\mathrm{CAPTURE}_{\mathrm{HC}}\);
  \item HC-to-OP handoff and OP-entry paragraphs, by replacing repeated transfer and normal-form prose with \(\mathrm{HANDOFF}\), \(\mathrm{NORMALIZE}\), and \(\mathrm{ROUTE}\);
  \item OP~1/4--OP~5(B) comparison paragraphs, by writing the inseparable branch as \(\mathrm{COMPARE}\to\mathrm{RESOLVE}\to\mathrm{RETURN}\);
  \item probe and Monte-Carlo paragraphs, by reading them as \(\mathrm{SCAN}\to\mathrm{RANK}\to\mathrm{CANDIDATE}\) and not as proof generators;
  \item historical witness paragraphs, by marking reusable tools without deleting the preserved genealogy.
\end{enumerate}
\end{proposition}

\begin{proof}
Each listed zone contains one of the patterns displayed in Definition~\ref{def:compression-map-from-inherited-text-to-mml-mma-operations}. The assigned operation records the proof-use of the zone and can therefore replace later repetitive explanation. Lemma~\ref{lem:compression-map-preserves-proof-content} prevents the compression from changing mathematical status.
\end{proof}

\begin{theorem}[MML/MMA compression theorem for inherited HC/OP text]
\label{thm:mml-mma-compression-theorem-for-inherited-hcop-text}
Every inherited HC/OP compression target that matches one of the patterns in Definition~\ref{def:compression-map-from-inherited-text-to-mml-mma-operations} admits a finite MML/MMA reading under
\[
  \mathrm{QAM}_{\mathrm{HCOP}}^{\mathrm{loc}}
  \longrightarrow
  \mathrm{MML}_{\mathrm{HCOP}}^{\mathrm{loc}}
  \longrightarrow
  \mathrm{MMA}_{\mathrm{HCOP}}^{\mathrm{loc}}.
\]
This reading may be used for later compression, indexing, and audit of the inherited text. It may not be used to erase the historical witness or to manufacture an OP closure certificate.
\end{theorem}

\begin{proof}
The r07 interface theorem, Theorem~\ref{thm:hc-routed-op-entry-as-finite-mml-mma-program}, already supplies the finite typed execution stack for the current HC-routed OP-entry chain. Definition~\ref{def:compression-map-from-inherited-text-to-mml-mma-operations} extends this stack as an expository map over inherited text patterns that perform the same roles. Since the map is role-preserving and certificate-neutral by Lemma~\ref{lem:compression-map-preserves-proof-content}, it yields a finite MML/MMA reading for each matching target while leaving the original mathematical content intact.
\end{proof}

\begin{corollary}[v3.34.0-r08 working rule]
\label{cor:v3340-r08-working-rule}
The next editorial pass may shorten repeated HC/OP explanations by citing \(\mathfrak C_{\mathrm{HCOP}}\), but it should not delete active proof material yet. The safe immediate use of the compression map is to add labels, cross-references, tool tables, and short operation names to inherited blocks; stronger textual shortening should wait until the corresponding OP branch has either closed or returned to a stable frontier state.
\end{corollary}

\begin{proof}
Immediate from Theorem~\ref{thm:mml-mma-compression-theorem-for-inherited-hcop-text}. The compression map is currently an audit and refactoring layer, not a release-level deletion license.
\end{proof}



\subsection{First compression-guided HC/OP tool normal forms}
\label{subsec:first-compression-guided-hcop-tool-normal-forms}

\begin{definition}[Reusable HC/OP tool normal form]
\label{def:first-reusable-hcop-tool-normal-form}
A \emph{reusable HC/OP tool normal form} is a finite typed contract
\[
  \mathfrak T
  =
  \langle
    \mathrm{name},
    \mathrm{input},
    \mathrm{guard},
    \mathrm{operation},
    \mathrm{output},
    \mathrm{status}
  \rangle
\]
obtained from the compression map \(\mathfrak C_{\mathrm{HCOP}}\). Its entries have the following roles.
\begin{enumerate}[leftmargin=2em,label=(\roman*)]
  \item \(\mathrm{name}\) is one of the finite MML/MMA operation names already admitted by Definitions~\ref{def:minimal-mml-hcop-instruction-alphabet} and~\ref{def:compression-map-from-inherited-text-to-mml-mma-operations};
  \item \(\mathrm{input}\) is the typed object family consumed by the tool;
  \item \(\mathrm{guard}\) is the admissibility, route, frontier, or certificate-neutrality condition under which the tool may be invoked;
  \item \(\mathrm{operation}\) is the finite MML instruction word together with its MMA readout semantics;
  \item \(\mathrm{output}\) is the typed object family returned by the tool;
  \item \(\mathrm{status}\) records whether the tool is diagnostic, normalizing, routing, comparison-level, frontier-level, or certificate-testing.
\end{enumerate}
A tool normal form is an invocation interface only. It may shorten later exposition by replacing repeated safety prose with a typed reference, but it does not erase the original proof material and does not create a theorematic certificate.
\end{definition}

\begin{definition}[First extracted HC/OP tool library]
\label{def:first-extracted-hcop-tool-library}
The first compression-guided HC/OP tool library is the finite family
\[
  \mathbb T_{\mathrm{HCOP}}^{(1)}
  :=
  \{
    \mathfrak T_{\mathrm{lock}},
    \mathfrak T_{\mathrm{handoff}},
    \mathfrak T_{\mathrm{norm}},
    \mathfrak T_{\mathrm{route}},
    \mathfrak T_{\mathrm{cmp}},
    \mathfrak T_{\mathrm{resolve}},
    \mathfrak T_{\mathrm{probe}},
    \mathfrak T_{\mathrm{cert}},
    \mathfrak T_{\mathrm{frontier}},
    \mathfrak T_{\mathrm{archive}}
  \}.
\]
Its intended contracts are summarized as follows.

{\small
\renewcommand{\arraystretch}{1.18}
\begin{longtable}{|p{0.18\textwidth}|p{0.22\textwidth}|p{0.22\textwidth}|p{0.26\textwidth}|}
\hline
\textbf{Tool} & \textbf{Input} & \textbf{Output} & \textbf{Guard / status} \\
\hline
\endfirsthead
\hline
\textbf{Tool} & \textbf{Input} & \textbf{Output} & \textbf{Guard / status} \\
\hline
\endhead
\(\mathfrak T_{\mathrm{lock}}\) & local HC candidate, corridor data & locked/captured HC candidate & guard: inherited route tag and local admissibility; status: candidate/capture \\
\hline
\(\mathfrak T_{\mathrm{handoff}}\) & transport-stable HC-capture chain & OP-facing handoff packet & guard: zero or controlled HC transport defect; status: handoff \\
\hline
\(\mathfrak T_{\mathrm{norm}}\) & handoff packet & \(\mathcal N_{\mathrm{OP}}^{\mathrm{HC}}\) with live residual triple & guard: no change of theorematic status; status: normalization \\
\hline
\(\mathfrak T_{\mathrm{route}}\) & \(\mathcal N_{\mathrm{OP}}^{\mathrm{HC}}\), selected residual branch & finite OP-facing task family & guard: residual type and route tag match; status: routing \\
\hline
\(\mathfrak T_{\mathrm{cmp}}\) & paired \(\rho_{\mathrm{CL}}\oplus\rho_{\Theta R}\) branch & \(\Delta_{\mathrm{sep}}^{\mathrm{HC}}\) and comparison record & guard: inseparable branch has not yet split; status: comparison \\
\hline
\(\mathfrak T_{\mathrm{resolve}}\) & comparison record & bind, split, or frontier next state & guard: decision rule returns exactly one branch; status: resolution \\
\hline
\(\mathfrak T_{\mathrm{probe}}\) & probe/Monte-Carlo/decoder scan surface & ranked admissible candidate family & guard: diagnostic only; status: scan/rank/candidate \\
\hline
\(\mathfrak T_{\mathrm{cert}}\) & trigger-persistence or certificate question & yes/no/pending certificate-test record & guard: asks only whether a mathematical trigger exists; status: certificate-test \\
\hline
\(\mathfrak T_{\mathrm{frontier}}\) & non-closing but improving state & smaller frontier-certified pre-entry state & guard: strict controlled frontier descent; status: frontier \\
\hline
\(\mathfrak T_{\mathrm{archive}}\) & historical witness block & reusable tool role plus preserved genealogy & guard: no deletion of historical proof memory; status: archive-to-tool \\
\hline
\end{longtable}
}
\end{definition}

\begin{lemma}[Extracted tool contracts are typed MML/MMA readings]
\label{lem:extracted-tool-contracts-are-typed-mml-mma-readings}
Every member of \(\mathbb T_{\mathrm{HCOP}}^{(1)}\) is a typed MML/MMA reading of one or more rows of the compression map \(\mathfrak C_{\mathrm{HCOP}}\).
\end{lemma}

\begin{proof}
Definition~\ref{def:first-extracted-hcop-tool-library} is obtained by grouping the rows of Definition~\ref{def:compression-map-from-inherited-text-to-mml-mma-operations} by proof-use. The lock, handoff, normalization, routing, comparison, resolution, return, scan/rank, certificate-test, frontier-descent, and archive-to-tool roles are exactly the roles already admitted by the compression map and by the r07 instruction alphabet. Hence each extracted contract is a typed reading over the existing QAM/MML/MMA stack, not a new primitive operation.
\end{proof}

\begin{proposition}[Immediate simplification rule for inherited HC/OP prose]
\label{prop:immediate-simplification-rule-for-inherited-hcop-prose}
Let \(B\) be an inherited HC/OP text block that matches one of the patterns of Definition~\ref{def:compression-map-from-inherited-text-to-mml-mma-operations}. If \(B\) has already stated its hypotheses and proof restrictions once, then later repetitions of the same proof-use may be replaced by a citation to the corresponding tool contract in \(\mathbb T_{\mathrm{HCOP}}^{(1)}\), provided the replacement also cites the original local statement or keeps it in place.
\end{proposition}

\begin{proof}
By Lemma~\ref{lem:compression-map-preserves-proof-content}, the compression map changes only the expository role of a block and preserves the original mathematical content. By Lemma~\ref{lem:extracted-tool-contracts-are-typed-mml-mma-readings}, the corresponding tool contract records the same role in a finite typed form. Therefore a later repetition can be shortened to a tool invocation as long as the first occurrence and its hypotheses remain available. The proviso prevents the normal form from becoming a deletion rule.
\end{proof}

\begin{proposition}[Where the first tool layer helps the existing text]
\label{prop:where-first-tool-layer-helps-existing-text}
The library \(\mathbb T_{\mathrm{HCOP}}^{(1)}\) gives immediate refactoring leverage in the following existing zones of the Companion.
\begin{enumerate}[leftmargin=2em,label=(\roman*)]
  \item The HC-lock and capture paragraphs may cite \(\mathfrak T_{\mathrm{lock}}\) instead of restating the full candidate/corridor discipline.
  \item The r01--r03 OP-entry paragraphs may cite \(\mathfrak T_{\mathrm{handoff}}\), \(\mathfrak T_{\mathrm{norm}}\), and \(\mathfrak T_{\mathrm{route}}\) instead of repeating the full handoff-normalize-route chain.
  \item The r04 OP~1/4--OP~5(B) branch may cite \(\mathfrak T_{\mathrm{cmp}}\) and \(\mathfrak T_{\mathrm{resolve}}\) to express the joint-versus-separated comparison.
  \item Probe, Monte-Carlo, decoder, and finite-search paragraphs may cite \(\mathfrak T_{\mathrm{probe}}\) to emphasize that they rank candidates rather than certify closure.
  \item Trigger-persistence paragraphs may cite \(\mathfrak T_{\mathrm{cert}}\) to keep certification separate from execution.
  \item Historical witness paragraphs may cite \(\mathfrak T_{\mathrm{archive}}\) to extract reusable tools without shortening the preserved genealogy.
\end{enumerate}
\end{proposition}

\begin{proof}
Each zone listed is one of the main inherited zones already identified in Proposition~\ref{prop:main-inherited-text-zones-helped-by-compression-map}. Definition~\ref{def:first-extracted-hcop-tool-library} assigns the corresponding reusable contract. The proposition is therefore the first explicit refactoring consequence of the r08 compression map.
\end{proof}

\begin{theorem}[First compression-guided HC/OP reusable-tool layer]
\label{thm:first-compression-guided-hcop-reusable-tool-layer}
The compression map \(\mathfrak C_{\mathrm{HCOP}}\) induces the finite reusable-tool layer \(\mathbb T_{\mathrm{HCOP}}^{(1)}\). This layer may be used in later Companion revisions to shorten repeated HC/OP safety prose, index historical proof material, and audit OP-entry execution traces. It may not be used to delete active proof content or to manufacture an OP closure certificate.
\end{theorem}

\begin{proof}
Existence of \(\mathbb T_{\mathrm{HCOP}}^{(1)}\) is Definition~\ref{def:first-extracted-hcop-tool-library}. Its relation to the compression map is Lemma~\ref{lem:extracted-tool-contracts-are-typed-mml-mma-readings}. Its safe simplification use is Proposition~\ref{prop:immediate-simplification-rule-for-inherited-hcop-prose}. Its no-false-certificate restriction follows from Lemma~\ref{lem:mml-mma-execution-preserves-route-and-frontier-discipline} and Lemma~\ref{lem:compression-map-preserves-proof-content}. Thus the layer is a reusable tool interface over the existing mathematics and not a new theorematic engine.
\end{proof}

\begin{corollary}[v3.34.0-r09 working rule]
\label{cor:v3340-r09-working-rule}
The next safe editorial use of MML/MMA is to annotate repeated inherited HC/OP paragraphs by the tool contracts in \(\mathbb T_{\mathrm{HCOP}}^{(1)}\). A future shortening pass may replace repeated explanations by these tool names only after the original local statement, hypotheses, and proof restrictions remain explicitly available elsewhere in the document.
\end{corollary}

\begin{proof}
Immediate from Theorem~\ref{thm:first-compression-guided-hcop-reusable-tool-layer}.
\end{proof}

\subsection{First tool-contract application to the active HC/OP proof spine}
\label{subsec:first-tool-contract-application-active-hcop-proof-spine}

\begin{definition}[Tool-indexed active HC/OP spine]
\label{def:tool-indexed-active-hcop-spine}
The \emph{tool-indexed active HC/OP spine} is the finite contract trace
\[
  \mathcal P_{\mathrm{HCOP}}^{\mathrm{tool}}
  :=
  \mathfrak T_{\mathrm{lock}}
  ;\mathfrak T_{\mathrm{handoff}}
  ;\mathfrak T_{\mathrm{norm}}
  ;\mathfrak T_{\mathrm{route}}
  ;\mathfrak T_{\mathrm{cmp}}
  ;\mathfrak T_{\mathrm{resolve}}
  ;\mathfrak T_{\mathrm{cert/frontier}} .
\]
It is a compact citation layer over the r01--r04 HC/OP proof spine. The trace is read as follows: lock/capture prepares the local HC candidate, handoff transports it toward OP-facing form, normalization produces \(\mathcal N_{\mathrm{OP}}^{\mathrm{HC}}\), routing selects a live residual work family, comparison executes the paired OP~1/4--OP~5(B) readout when separation is not yet justified, resolution returns bind, split, or frontier, and the final certificate/frontier test records whether the branch may promote or must remain a smaller frontier-certified state.
\end{definition}

\begin{definition}[Tool-contract application site]
\label{def:tool-contract-application-site}
A paragraph, lemma, proposition, or proof step is a \emph{tool-contract application site} for a contract \(\mathfrak T\in\mathbb T_{\mathrm{HCOP}}^{(1)}\) if all of the following are explicit either locally or by a nearby reference.
\begin{enumerate}[leftmargin=2em,label=(\roman*)]
  \item the typed input object of \(\mathfrak T\);
  \item the guard of \(\mathfrak T\), including route, frontier, and admissibility restrictions;
  \item the operation or comparison performed;
  \item the typed output object;
  \item the status restriction that the invocation is diagnostic, routing, comparison-level, frontier-level, or certificate-testing rather than certificate-creating.
\end{enumerate}
At such a site, later repetitions may cite \(\mathfrak T\) without restating the full safety prose, provided the original statement and its hypotheses remain visible in the document.
\end{definition}

\begin{lemma}[Active HC/OP blocks are contract-application sites]
\label{lem:active-hcop-blocks-are-contract-application-sites}
The r01--r04 active HC/OP blocks are application sites for the corresponding contracts in Definition~\ref{def:tool-indexed-active-hcop-spine}.
\end{lemma}

\begin{proof}
The r01 transport-stable HC-capture and handoff block states the local candidate, route-faithful transport guard, and OP-facing packet output, hence matches \(\mathfrak T_{\mathrm{lock}}\) and \(\mathfrak T_{\mathrm{handoff}}\). The r02 normal form block states the handoff input, the live residual triple, and the non-closure status restriction, hence matches \(\mathfrak T_{\mathrm{norm}}\). The r03 task-routing block states the residual branch, admissible work family, and route tag, hence matches \(\mathfrak T_{\mathrm{route}}\). The r04 paired-comparison block states the inseparable \(\rho_{\mathrm{CL}}\oplus\rho_{\Theta R}\) input, the comparison defect, and the bind/split/frontier return, hence matches \(\mathfrak T_{\mathrm{cmp}}\) and \(\mathfrak T_{\mathrm{resolve}}\). The inherited no-false-certificate rule supplies the final certificate/frontier guard. Therefore each block satisfies Definition~\ref{def:tool-contract-application-site}.
\end{proof}

\begin{proposition}[Tool-indexed rewriting of the active HC/OP chain]
\label{prop:tool-indexed-rewriting-active-hcop-chain}
The active r01--r04 HC/OP chain admits the compressed reading
\[
  \mathrm{HC}_{\mathrm{loc}}
  \xrightarrow{\mathfrak T_{\mathrm{lock}}}
  \mathrm{HC}_{\mathrm{cap}}
  \xrightarrow{\mathfrak T_{\mathrm{handoff}}}
  \mathcal P_{\mathrm{handoff}}
  \xrightarrow{\mathfrak T_{\mathrm{norm}}}
  \mathcal N_{\mathrm{OP}}^{\mathrm{HC}}
  \xrightarrow{\mathfrak T_{\mathrm{route}}}
  \mathcal T_{\mathrm{OP}}^{\mathrm{HC}}
  \xrightarrow{\mathfrak T_{\mathrm{cmp}},\,\mathfrak T_{\mathrm{resolve}}}
  \{\mathrm{bind},\mathrm{split},\mathrm{frontier}\}.
\]
This rewriting is admissible as a cross-reference and audit normal form, but it is not an independent proof of any OP closure.
\end{proposition}

\begin{proof}
By Lemma~\ref{lem:active-hcop-blocks-are-contract-application-sites}, each arrow is backed by an already present active HC/OP block and by a typed contract in \(\mathbb T_{\mathrm{HCOP}}^{(1)}\). The chain is therefore a finite MML/MMA execution trace in the sense of Theorem~\ref{thm:hc-routed-op-entry-as-finite-mml-mma-program}. Since Theorem~\ref{thm:first-compression-guided-hcop-reusable-tool-layer} forbids tool contracts from manufacturing certificates, the compressed trace is an audit normal form only.
\end{proof}

\begin{lemma}[Tool application preserves local hypotheses]
\label{lem:tool-application-preserves-local-hypotheses}
Replacing a repeated explanation by a citation to a tool contract preserves the local hypotheses of the original explanation exactly when the application site cites the original statement or restates the required guard.
\end{lemma}

\begin{proof}
A tool contract records input, guard, operation, output, and status; it does not by itself reproduce every local hypothesis. Definition~\ref{def:tool-contract-application-site} therefore requires the relevant guard to be explicit locally or by reference. Under that requirement, the replacement removes only duplicated prose and leaves the logical premises available. Without that requirement the replacement would be unsafe, which is why the definition includes it as a guard.
\end{proof}

\begin{corollary}[First safe shortening target]
\label{cor:first-safe-shortening-target}
The first safe shortening target is repeated no-false-certificate, route-preservation, and frontier-preservation prose following an already stated HC/OP operation. Such repetitions may be replaced by a sentence of the form:
\begin{quote}
\small
By the corresponding tool contract in \(\mathbb T_{\mathrm{HCOP}}^{(1)}\), this is an admissible MML/MMA execution step and not a closure certificate.
\end{quote}
The original definition, theorem, or proof restriction must remain present elsewhere in the same local chain.
\end{corollary}

\begin{proof}
The repeated phrases listed are status and guard reminders, not new mathematical premises. Lemma~\ref{lem:tool-application-preserves-local-hypotheses} permits the repetition to be replaced once the underlying guard has already been stated. The no-false-certificate clause is exactly the status restriction of \(\mathfrak T_{\mathrm{cert}}\) and of the active MML/MMA interface.
\end{proof}

\begin{theorem}[First tool-contract application theorem]
\label{thm:first-tool-contract-application-theorem}
The reusable tool library \(\mathbb T_{\mathrm{HCOP}}^{(1)}\) can now be applied to the active HC/OP proof spine as a conservative cross-reference layer. It yields shorter notation, clearer audit points, and explicit preservation of route, frontier, and certificate restrictions, without changing theorematic status and without deleting historical witness material.
\end{theorem}

\begin{proof}
Shorter notation follows from Proposition~\ref{prop:tool-indexed-rewriting-active-hcop-chain}. Auditability follows because Definition~\ref{def:tool-contract-application-site} lists the required input, guard, operation, output, and status fields for every invocation. Preservation of local hypotheses follows from Lemma~\ref{lem:tool-application-preserves-local-hypotheses}. The no-deletion and no-false-certificate restrictions are inherited from Theorem~\ref{thm:first-compression-guided-hcop-reusable-tool-layer}. Hence the layer is a conservative application of the r09 tool library rather than a new proof step.
\end{proof}

\begin{corollary}[v3.34.0-r10 working rule]
\label{cor:v3340-r10-working-rule}
Subsequent revisions may begin annotating the active HC/OP chain by \(\mathcal P_{\mathrm{HCOP}}^{\mathrm{tool}}\). They should first replace only duplicated safety prose and local route/frontier reminders. They should not delete first occurrences of definitions, hypotheses, proof restrictions, historical witness blocks, or any statement needed to verify the no-false-certificate rule.
\end{corollary}

\begin{proof}
Immediate from Theorem~\ref{thm:first-tool-contract-application-theorem} and Corollary~\ref{cor:first-safe-shortening-target}.
\end{proof}


\subsection{First tool-guarded reduction protocol for inherited HC/OP prose}
\label{subsec:first-tool-guarded-reduction-protocol-inherited-hcop-prose}

\begin{definition}[Guarded HC/OP reduction packet]
\label{def:guarded-hcop-reduction-packet}
A \emph{guarded HC/OP reduction packet} is a finite tuple
\[
  \mathcal R_{\mathrm{HCOP}}^{(1)}
  :=
  \langle B,\, \mathfrak T,\, G_B,\, S_B,\, W_B\rangle
\]
with the following data.
\begin{enumerate}[leftmargin=2em,label=(\roman*)]
  \item \(B\) is an inherited or active HC/OP paragraph, lemma step, proof reminder, or diagnostic status clause.
  \item \(\mathfrak T\in\mathbb T_{\mathrm{HCOP}}^{(1)}\) is the tool contract assigned to the repeated role of \(B\).
  \item \(G_B\) records the local guard: typed input, route tag, frontier condition, admissibility condition, and no-false-certificate status.
  \item \(S_B\) records the theorematic status of the occurrence: definition, hypothesis, proof step, diagnostic reminder, certificate question, frontier update, or historical witness.
  \item \(W_B\) records where the unreduced witness of \(B\) remains available in the document.
\end{enumerate}
The packet is a reduction-control object only. It does not change the mathematical claim carried by \(B\).
\end{definition}

\begin{definition}[Reducible and non-reducible HC/OP occurrences]
\label{def:reducible-and-nonreducible-hcop-occurrences}
Let \(B\) be an occurrence in the HC/OP proof spine or in an inherited HC/OP support zone. The occurrence is \emph{r11-reducible} if it satisfies all of the following conditions.
\begin{enumerate}[leftmargin=2em,label=(\roman*)]
  \item It repeats a safety, route-preservation, frontier-preservation, diagnostic-status, scan/rank/candidate, comparison-status, or no-false-certificate clause already stated in the same local chain.
  \item Its assigned tool contract \(\mathfrak T\) is one of the contracts of Definition~\ref{def:first-extracted-hcop-tool-library}.
  \item The original hypotheses, guards, proof restrictions, and local statement remain visible either in \(B\)'s first occurrence or in a nearby cited statement.
  \item The replacement sentence explicitly says that the citation is a tool-contract invocation and not a closure certificate.
\end{enumerate}
The occurrence is \emph{r11-non-reducible} if it introduces a first definition, a new object, a new hypothesis, a new guard, a new route tag, a certificate test, a theorematic status change, a proof-critical inequality, or historical witness material.
\end{definition}

\begin{definition}[Guarded citation normal form]
\label{def:guarded-citation-normal-form}
The preferred normal form for replacing an r11-reducible occurrence is
\begin{quote}
\small
By the guarded tool contract \(\mathfrak T\) and the local guard recorded in \(G_B\), this is an admissible MML/MMA execution or diagnostic step; it preserves route/frontier discipline and does not create an OP closure certificate.
\end{quote}
When the local context is already clear, the sentence may be shortened to a parenthetical form:
\[
  \text{guarded by }\mathfrak T\text{; route/frontier preserved; no certificate manufactured.}
\]
\end{definition}

\begin{lemma}[Guarded reduction preserves proof content]
\label{lem:guarded-reduction-preserves-proof-content}
If \(B\) is r11-reducible and is replaced by the guarded citation normal form of Definition~\ref{def:guarded-citation-normal-form}, then the replacement preserves the proof content of the local chain.
\end{lemma}

\begin{proof}
By Definition~\ref{def:reducible-and-nonreducible-hcop-occurrences}, only repeated status and guard reminders are eligible for replacement. The original hypotheses, proof restrictions, and first local statement remain available. By Definition~\ref{def:guarded-citation-normal-form}, the replacement still records the assigned tool contract, preserves route/frontier discipline, and repeats the no-false-certificate status. Hence the replacement removes duplicated prose only; it does not remove the premise, operation, guard, output, or theorematic status of the original argument.
\end{proof}

\begin{lemma}[Non-reducible occurrences protect theorematic status]
\label{lem:nonreducible-occurrences-protect-theorematic-status}
Every r11-non-reducible occurrence must remain explicit in the document unless a later revision supplies a stronger local theorem that restates the same content with all hypotheses and guards.
\end{lemma}

\begin{proof}
The non-reducible occurrences listed in Definition~\ref{def:reducible-and-nonreducible-hcop-occurrences} are precisely the occurrences that create or alter mathematical content: they introduce objects, hypotheses, route tags, guards, certificate tests, status changes, proof-critical inequalities, or historical witness material. Replacing such an occurrence by a tool-name alone would hide a premise or change the audit surface. Therefore they are protected until a later theorem restates them explicitly.
\end{proof}

\begin{proposition}[First usable simplification map]
\label{prop:first-usable-hcop-simplification-map}
The r11 reduction protocol induces the partial simplification map
\[
  \mathsf{Red}_{\mathrm{HCOP}}^{(1)}:
  B
  \dashrightarrow
  \bigl(\mathfrak T(B),G_B,S_B,W_B\bigr),
\]
where the map is defined exactly on r11-reducible occurrences and undefined on r11-non-reducible occurrences.
\end{proposition}

\begin{proof}
For a reducible occurrence, Definitions~\ref{def:guarded-hcop-reduction-packet} and~\ref{def:guarded-citation-normal-form} assign the corresponding tool contract, guard, status, and witness pointer. Lemma~\ref{lem:guarded-reduction-preserves-proof-content} proves that this replacement is content-preserving. For a non-reducible occurrence, Lemma~\ref{lem:nonreducible-occurrences-protect-theorematic-status} forbids replacement by a mere tool invocation. Hence the simplification map is partial and conservative.
\end{proof}

\begin{theorem}[First tool-guarded HC/OP reduction theorem]
\label{thm:first-tool-guarded-hcop-reduction-theorem}
The active HC/OP tool layer now supports a first conservative reduction protocol. Repeated safety prose, route/frontier reminders, diagnostic-status clauses, and no-false-certificate clauses may be replaced by guarded tool citations according to \(\mathsf{Red}_{\mathrm{HCOP}}^{(1)}\). First definitions, first hypotheses, local guards, certificate tests, proof-critical inequalities, theorematic status changes, and historical witness material remain non-reducible. Therefore the protocol simplifies later exposition without altering the OP status of the document.
\end{theorem}

\begin{proof}
The permitted replacements are exactly the defined values of the partial map in Proposition~\ref{prop:first-usable-hcop-simplification-map}. Content preservation follows from Lemma~\ref{lem:guarded-reduction-preserves-proof-content}. Protection of theorematic status follows from Lemma~\ref{lem:nonreducible-occurrences-protect-theorematic-status}. The no-false-certificate restriction is inherited from Theorem~\ref{thm:first-tool-contract-application-theorem}. Thus the reduction protocol is a simplifier for repeated exposition and an audit aid, not a new OP closure mechanism.
\end{proof}

\begin{corollary}[v3.34.0-r11 working rule]
\label{cor:v3340-r11-working-rule}
The next safe use of the tool layer is a local annotation pass: mark repeated HC/OP safety paragraphs as reducible or non-reducible before shortening them. Actual deletion or replacement should occur only after the first occurrence, hypotheses, guards, and historical witness pointer have been verified.
\end{corollary}

\begin{proof}
Immediate from Theorem~\ref{thm:first-tool-guarded-hcop-reduction-theorem}.
\end{proof}



\subsection{First guarded-reduction application pass for the active HC/OP proof spine}
\label{subsec:first-guarded-reduction-application-pass-active-hcop-proof-spine}

\begin{definition}[First guarded-reduction application ledger]
\label{def:first-guarded-reduction-application-ledger}
The \emph{first guarded-reduction application ledger} is the finite list
\[
  \mathcal L_{\mathrm{HCOP}}^{(1)}
  :=
  \{\ell_{\mathrm{lock}},\ell_{\mathrm{handoff}},\ell_{\mathrm{norm}},
    \ell_{\mathrm{route}},\ell_{\mathrm{cmp}},\ell_{\mathrm{probe}},
    \ell_{\mathrm{archive}}\}.
\]
Each ledger entry has the form
\[
  \ell
  =
  \langle S_{\ell},\, W_{\ell},\, \mathfrak T_{\ell},\,
  G_{\ell},\, R_{\ell}\rangle,
\]
where \(S_{\ell}\) is the local site class, \(W_{\ell}\) is the protected unreduced witness, \(\mathfrak T_{\ell}\in\mathbb T_{\mathrm{HCOP}}^{(1)}\) is the tool contract used for later citation, \(G_{\ell}\) is the local guard data, and \(R_{\ell}\) is the repeated prose tail that may be replaced by a guarded citation.  The ledger is an application of the partial map \(\mathsf{Red}_{\mathrm{HCOP}}^{(1)}\) from Proposition~\ref{prop:first-usable-hcop-simplification-map}; it is not an instruction to remove the first occurrence of any definition, guard, certificate test, proof-critical inequality, or historical witness.
\end{definition}

\begin{definition}[r12 guarded citation sites]
\label{def:r12-guarded-citation-sites}
The entries of \(\mathcal L_{\mathrm{HCOP}}^{(1)}\) are interpreted as follows.
\begin{enumerate}[leftmargin=2em,label=(\roman*)]
  \item \(\ell_{\mathrm{lock}}\) covers later repeated lock/capture reminders after the first local HC-capture statement has been given.  Its citation contract is \(\mathfrak T_{\mathrm{lock}}\), and its protected witness is the unreduced local capture chain.
  \item \(\ell_{\mathrm{handoff}}\) covers later repeated claims that HC capture is only a handoff into OP-entry and not an OP closure.  Its citation contracts are \(\mathfrak T_{\mathrm{handoff}}\) and \(\mathfrak T_{\mathrm{cert}}\).
  \item \(\ell_{\mathrm{norm}}\) covers later repeated normalization reminders for the HC-routed OP normal form.  Its citation contract is \(\mathfrak T_{\mathrm{norm}}\), guarded by the already visible residual triple and route tag.
  \item \(\ell_{\mathrm{route}}\) covers later repeated route/frontier-preservation clauses in the finite task-routing object.  Its citation contracts are \(\mathfrak T_{\mathrm{route}}\) and \(\mathfrak T_{\mathrm{frontier}}\).
  \item \(\ell_{\mathrm{cmp}}\) covers later repeated status clauses for the paired OP~1/4--OP~5(B) comparison branch.  Its citation contracts are \(\mathfrak T_{\mathrm{cmp}}\) and \(\mathfrak T_{\mathrm{resolve}}\), and its protected witness includes the explicit comparison defect \(\Delta_{\mathrm{sep}}^{\mathrm{HC}}\).
  \item \(\ell_{\mathrm{probe}}\) covers later repeated probe-diagnostic clauses stating that scan/rank/candidate output is not a certificate.  Its citation contracts are \(\mathfrak T_{\mathrm{probe}}\) and \(\mathfrak T_{\mathrm{cert}}\).
  \item \(\ell_{\mathrm{archive}}\) covers later repeated archive-to-tool extraction reminders.  Its citation contract is \(\mathfrak T_{\mathrm{archive}}\), and its protected witness is the unreduced historical source block.
\end{enumerate}
\end{definition}

\begin{definition}[Ledger citation normal form]
\label{def:ledger-citation-normal-form}
For an occurrence assigned to an entry \(\ell\in\mathcal L_{\mathrm{HCOP}}^{(1)}\), the preferred r12 citation normal form is
\[
  \mathrm{GRCite}_{\mathrm{HCOP}}(\ell)
  :=
  \langle \mathfrak T_{\ell},G_{\ell},W_{\ell};\mathrm{noCert}\rangle.
\]
In prose this is read as:
\begin{quote}\small
Guarded by \(\mathfrak T_{\ell}\) with local guard \(G_{\ell}\) and protected witness \(W_{\ell}\); route/frontier discipline is preserved and no certificate is manufactured.
\end{quote}
The normal form may replace only a repeated tail \(R_{\ell}\), never the protected witness \(W_{\ell}\).
\end{definition}

\begin{lemma}[Ledger entries are r11-reducible where used]
\label{lem:ledger-entries-r11-reducible-where-used}
Every later occurrence covered by a ledger entry in Definition~\ref{def:r12-guarded-citation-sites} is r11-reducible exactly when its first local witness, guard, and theorematic status remain visible.
\end{lemma}

\begin{proof}
Each ledger entry records a repeated prose tail: safety status, route/frontier preservation, diagnostic non-certification, comparison status, or archive-to-tool extraction.  These are precisely the reducible classes of Definition~\ref{def:reducible-and-nonreducible-hcop-occurrences}.  The same definition excludes replacement whenever a first occurrence, guard, certificate test, proof-critical inequality, status change, or historical witness would be hidden.  Hence the ledger is admissible only at later occurrences for which the protected witness \(W_{\ell}\) remains visible.
\end{proof}

\begin{proposition}[First guarded-reduction application is conservative]
\label{prop:first-guarded-reduction-application-conservative}
Applying \(\mathrm{GRCite}_{\mathrm{HCOP}}(\ell)\) to the repeated tail of any admissible ledger entry preserves the mathematical content of the active HC/OP proof spine.
\end{proposition}

\begin{proof}
The replacement is defined only for repeated tails \(R_{\ell}\).  By Definition~\ref{def:first-guarded-reduction-application-ledger}, the protected witness \(W_{\ell}\), the tool contract \(\mathfrak T_{\ell}\), the local guard \(G_{\ell}\), and the no-certificate status remain part of the citation normal form.  Lemma~\ref{lem:ledger-entries-r11-reducible-where-used} reduces the claim to Lemma~\ref{lem:guarded-reduction-preserves-proof-content}.  Thus the operation shortens only repeated exposition and leaves the proof spine unchanged.
\end{proof}

\begin{theorem}[First guarded-reduction application theorem]
\label{thm:first-guarded-reduction-application-theorem}
The active HC/OP proof spine now carries a first practical guarded-reduction application layer.  Later repeated HC/OP prose may cite the finite ledger \(\mathcal L_{\mathrm{HCOP}}^{(1)}\) in the normal form of Definition~\ref{def:ledger-citation-normal-form}.  This permits local exposition to be shortened after first occurrence, but it does not delete proof material, alter route tags, alter frontier status, convert diagnostics into certificates, or change the OP status of the document.
\end{theorem}

\begin{proof}
Existence of the finite ledger follows from Definition~\ref{def:first-guarded-reduction-application-ledger}.  Admissibility of its local uses follows from Lemma~\ref{lem:ledger-entries-r11-reducible-where-used}.  Conservativity follows from Proposition~\ref{prop:first-guarded-reduction-application-conservative}.  The no-false-certificate restriction is explicitly retained in Definition~\ref{def:ledger-citation-normal-form} and is inherited from Theorem~\ref{thm:first-tool-guarded-hcop-reduction-theorem}.  Hence the layer is a guarded reduction application pass rather than a new theorematic closure step.
\end{proof}

\begin{corollary}[v3.34.0-r12 working rule]
\label{cor:v3340-r12-working-rule}
The next safe simplification pass may replace only later repeated HC/OP tails that match an entry of \(\mathcal L_{\mathrm{HCOP}}^{(1)}\).  The first occurrence of each definition, hypothesis, guard, certificate test, proof-critical inequality, and historical witness remains explicit.
\end{corollary}

\begin{proof}
Immediate from Theorem~\ref{thm:first-guarded-reduction-application-theorem}.
\end{proof}

\begin{definition}[Defect transport equation]
\label{def:defect-transport-equation}
Let $X_n$ and $X_{n+1}$ be adjacent closure-liquid-crystal cells. Their defect transport equation is written
\[
  \mathcal D_{\mathrm{cell}}(X_n,X_{n+1})\neq 0
\]
whenever the first local failure is no longer merely screening-visible, but appears as a witness-grade or shell-grade obstruction transported across the current channel. The datum $\mathcal D_{\mathrm{cell}}$ is therefore not a second closure criterion; it is the cellular packaging of the first locally non-contractible obstruction.
\end{definition}

\begin{proposition}[OP~5 defect-transport reading]
\label{prop:op5-defect-transport-reading}
Within the current bundled OP~3$\to$OP~5/global channel, witness-grade OP~5 reopening may be read as defect transport on the same closure-liquid-crystal medium: a bundled global pass corresponds to defect-free or defect-screened transport, while a bundled obstruction corresponds to transported or accumulated non-contractible defect data.
\end{proposition}

\begin{proof}
By Theorem~\ref{thm:bundled-op5-diagnostic-theorem}, Theorem~\ref{thm:first-op3-op5-bundled-transfer}, and Theorem~\ref{thm:first-global-compressed-proof-run-through-extracted-global-tools}, the current channel already distinguishes a bundled pass from a bundled obstruction while subordinating packet/phase/epsilon diagnostics to the carrier-level decision. The witness-grade local/global reopening discipline then identifies the precise place where a local obstruction must be reopened rather than screened away. This is exactly the cellular reading of $\mathcal D_{\mathrm{cell}}\neq 0$: a transported or accumulated defect of one and the same carrier medium rather than a second primitive theory beside the closure channel.
\end{proof}

\begin{definition}[Liquid readout domain]
\label{def:liquid-readout-domain}
Fix a distinguished readout regime $f$. A \emph{liquid readout domain} in regime $f$ is a family of closure-liquid-crystal cells whose carrier data lie in one closure-bearing class while their visible shell/phase data realize one coherent readout phase of that same bulk order. Different regimes or sectors may therefore produce different liquid readout domains without changing the underlying carrier datum in the strong closure sense.
\end{definition}

\begin{remark}[Readouts as liquid-crystal domains of one bulk order]
\label{rem:readouts-as-liquid-crystal-domains}
The point of Definition~\ref{def:liquid-readout-domain} is not to introduce a literal condensed-matter ontology. It is to express the present structural reading more economically: the visible sectors are not independent primitive crystals, but phase-sensitive ordered domains of one deeper closure-bearing bulk carrier. This is the cellular version of the readout-relative principle already used throughout the Companion.
\end{remark}

\begin{definition}[Holographic Julia-boundary readout]
\label{def:holographic-julia-boundary-readout}
Let $\{X_n\}_{n\in I}$ be an admissible closure-liquid-crystal configuration along a chosen branch/route family. A \emph{holographic Julia-boundary readout} is a boundary map
\[
  \mathcal H_{\mathrm{J}}(\{X_n\}_{n\in I})
\]
that records the visible boundary pattern produced by the phase-/shell-sensitive readout of the same bulk carrier configuration. The intended output is a Julia-type boundary image, diagnostic strip, or related visual readout.
\end{definition}

\begin{proposition}[Holographic Julia-boundary principle]
\label{prop:holographic-julia-boundary-principle}
The Julia-type diagnostics used in the present Companion may be read as holographic boundary readouts of an underlying closure-liquid-crystal configuration. In particular, they remain diagnostic or screening-level boundary objects: they support the active closure architecture, but they do not create a new independent closure axiom.
\end{proposition}

\begin{proof}
The current document already treats packet-, phase-lock-, epsilon-, and frame-flicker diagnostics as visible readout layers subordinate to the carrier-level and closure-level statements; see Proposition~\ref{prop:packet-phase-epsilon-screening} and Corollary~\ref{cor:flicker-helix-subordination}. The same discipline applies to Julia-type visuals: they belong on the visible boundary-readout side and are informative exactly insofar as they track admissible passage, defect localization, or loss of admissibility in the bulk carrier medium. Hence they may be read holographically without becoming independent closure generators.
\end{proof}

\begin{remark}[v3.20 pilot status]
\label{rem:v320-pilot-status}
The present section should be read as a first pilot layer for v3.20.0. It upgrades neither the current OP~3/QAM closure package nor the bundled OP~3$\to$OP~5/global chain by new quantitative claims. Its role is organizational: it packages OP~3 as local cell compatibility, OP~5 as defect transport, visible sectors as liquid-crystal domains, Julia-type visuals as holographic boundary readout of one common carrier order, and tests a first addressed role layer in which the 0-band is reserved for primitive ZF$\leftrightarrow$QAM transport strata beneath later derived role families.
\end{remark}


\begin{remark}[Status of the compressed addressed crystal reserve surface in the present release line]
\label{rem:status-of-the-compressed-addressed-crystal-reserve-surface}
The earlier addressed-role pilot, the 0-band transport reservation, the locally licensed mixed addressed tensors, and the packet/corridor/hologram/readout-class ladder have now served their main purpose: they show that one closure-stable carrier may support visible domain phase, local defect loading, and holographic boundary readout without introducing a second primitive bulk family. In the present compressed line these ingredients are therefore retained only through one finite addressed crystal reserve surface.
\end{remark}

\begin{definition}[Compressed addressed crystal reserve token]
\label{def:compressed-addressed-crystal-reserve-token}
A \emph{compressed addressed crystal reserve token} is a finite package
\[
  \mathfrak A^{\mathrm{cr}}_{\mathrm{res}}
  :=
  \langle \mathcal B_0,\, \mathfrak M_{\mathrm{pil}},\, \mathfrak C^{\mathrm{term}}_{\mathrm{cr}} \rangle
\]
with the following intended reading.
\begin{enumerate}[leftmargin=2em,label=(\roman*)]
  \item $\mathcal B_0$ is the reserved primitive transport substrate between ZF and QAM;
  \item $\mathfrak M_{\mathrm{pil}}$ denotes the finite licensed addressed mixed-tensor pilot family consisting of the carrier--domain, domain--defect, and carrier--boundary couplings above that substrate;
  \item $\mathfrak C^{\mathrm{term}}_{\mathrm{cr}}$ is the terminal contraction-stable addressed crystal readout object that retains only the carrier/domain/defect/boundary content needed for later citation.
\end{enumerate}
The token is used only as a compressed reserve object. It does not claim a global unrestricted crystal calculus.
\end{definition}

\begin{proposition}[Compressed addressed crystal reserve package]
\label{prop:compressed-addressed-crystal-reserve-package}
Whenever a later argument uses only the finite addressed-crystal content encoded by
\[
  \mathfrak A^{\mathrm{cr}}_{\mathrm{res}},
\]
it may cite this single reserve token instead of reopening the full addressed-role pilot, the three local mixed addressed tensor pilots, or the packet/corridor/hologram/class-normal-form ladder. In particular, the compressed reserve token is citation-equivalent to the removed detailed reserve block for every argument that depends only on:
\begin{enumerate}[leftmargin=2em,label=(\roman*)]
  \item the primitive 0-band transport support,
  \item the three locally licensed mixed addressed couplings,
  \item one closure-stable carrier line together with its visible domain-face data,
  \item the corresponding local defect-load data,
  \item the retained holographic boundary profile,
  \item and the nonredundant interface data surviving passage to the terminal contraction-stable readout object.
\end{enumerate}
\end{proposition}

\begin{proof}
The removed detailed block introduced exactly these ingredients in layered form. First it reserved the 0-band as primitive transport substrate and treated later addressed families as derived above it. Second it licensed only three local mixed addressed tensor couplings: carrier--domain, domain--defect, and carrier--boundary. Third it unfolded a packet/corridor/hologram/readout-class ladder whose only later use was to justify passage to a terminal contraction-stable object carrying the retained carrier/domain/defect/boundary data. Therefore every later argument that needs only this finite retained content may cite the single reserve token without reopening the intermediate pilot steps. No stronger global claim is being made.
\end{proof}

\begin{corollary}[Use rule for the compressed addressed crystal reserve surface]
\label{cor:use-rule-for-the-compressed-addressed-crystal-reserve-surface}
A later proof may remain entirely on the compressed addressed crystal reserve surface as long as it needs only the finite retained content listed in Proposition~\ref{prop:compressed-addressed-crystal-reserve-package}. A finer reopening is required exactly when packet-specific, corridor-specific, subtype-specific, or pre-contraction addressed information is needed again.
\end{corollary}

\begin{proof}
This is the operational restatement of Proposition~\ref{prop:compressed-addressed-crystal-reserve-package}.
\end{proof}

\begin{lemma}[Closed antisymmetry of deficiency elements]
\label{lem:closed-deficiency-antisymmetry}
Let $f$ be a deficiency element in the relevant HC/closure test sector. If the
$J$-symmetric subspace is already essentially self-adjoint on the admissible
core, then the $J$-symmetric part of $f$ vanishes. Equivalently,
\[
  f=f^-\qquad\text{and hence}\qquad Jf=-f.
\]
\end{lemma}

\begin{proof}
Write
\[
  f=f^+ + f^-,\qquad f^+=\tfrac12(f+Jf),\qquad f^- = \tfrac12(f-Jf).
\]
By construction, $f^+$ lies in the $J$-symmetric sector. The repaired Section 12
argument shows that once the admissible $J$-symmetric core is already
essentially self-adjoint, a deficiency element cannot retain a nonzero
$J$-symmetric component. Otherwise one would obtain a nontrivial deficiency
vector in the symmetric core itself, contradicting essential self-adjointness.
Hence $f^+=0$, so $f=f^-$ and therefore $Jf=-f$.
\end{proof}

\begin{proposition}[Functional annihilation of deficiency elements by HC]
\label{prop:functional-hc-annihilation}
For every deficiency element $f$ in the closed HC test sector, the compensator
annihilates $f$ at the function level:
\[
  C(f)=\tfrac12(f+Jf)=0.
\]
\end{proposition}

\begin{proof}
By Lemma~\ref{lem:closed-deficiency-antisymmetry}, every such deficiency element
satisfies $Jf=-f$. Substituting into the defining compensator formula gives
\[
  C(f)=\tfrac12(f+Jf)=\tfrac12(f-f)=0.
\]
So the annihilation is a genuine functional consequence of closed
$J$-antisymmetry, not merely a heuristic tensor slogan.
\end{proof}

\begin{remark}[Tensor representations are secondary to the functional proof]
\label{rem:tensor-secondary-functional-proof}
Any tensor or matrix representative of the compensator is basis-dependent and
should be read only as a representation of the already established function-level
map. In particular, the annihilation of deficiency elements is proved by the
closed antisymmetry route
\[
  f^+=0 \Longrightarrow Jf=-f \Longrightarrow C(f)=0,
\]
not by a standalone coefficient-level matrix argument.
\end{remark}

\begin{lemma}[Anchor-compatibility lemma]
\label{lem:anchor-compatibility-lemma}
Compensator action is compatible with the orthogonal anchor. In particular, HC
stabilizes the admissible readout channel without destroying the distinction
between the anchored third direction and the first two sectors.
\end{lemma}

\begin{proof}
By Lemma~\ref{lem:orthogonal-anchor-lemma}, the third sector is not decorative
but the stabilizing orthogonal anchor. By Theorem~\ref{thm:hc-forcing-theorem},
HC is the mechanism that re-identifies shifted descriptions in a common output
channel. If HC destroyed the anchor distinction, then the stabilized output
would collapse back into an undifferentiated binary pairing and the role of the
third direction would be lost. That would contradict triolic minimality and the
very need for an orthogonal anchor. Therefore HC must act compatibly with the
anchored triolic organization.
\end{proof}

\begin{lemma}[Closure forcing lemma]
\label{lem:closure-forcing-lemma}
A triolic state with orthogonal anchor and forced HC is closure-stable.
Equivalently,
\[
\mathrm{Triolic}(x)\wedge \mathrm{Anchored}(x)\wedge \mathrm{HCForced}(x)
\Longrightarrow \mathrm{ClosureStable}(x).
\]
\end{lemma}

\begin{proof}
We verify the three contributing conditions and their combination.

\smallskip
\noindent\textbf{Step 1: Triolic structure.}
By Lemma~\ref{lem:triolic-minimality-lemma}, a triolic state carries a
three-sector organization $(K_1,K_2,K_3)$ as the minimal nondegenerate
structure compatible with compensator action and closure readout.
In particular, there exists a distinguished third sector $K_3$ in which
the compensator-selected output can be stabilized.

\smallskip
\noindent\textbf{Step 2: Orthogonal anchor.}
By Lemma~\ref{lem:orthogonal-anchor-lemma}, the third sector satisfies
$K_3\perp K_1$ and $K_3\perp K_2$, preventing $K_3$ from being
relabeled or absorbed into the dual pair $(K_1,K_2)$.
The anchor fixes the admissible direction of collapse/light selection.

\smallskip
\noindent\textbf{Step 3: HC forcing.}
By Lemma~\ref{lem:hc-admissibility-lemma}, the compensator-normalized
image $C(x)$ is admissible whenever $x$ is admissible. Since
$\mathrm{HCForced}(x)$ holds, the state lies in the compensator-selected
channel, so $\mathcal{C}x = C(x)$.

\smallskip
\noindent\textbf{Step 4: Closure stability.}
By the semantic--tensor correspondence of
Remark~\ref{rem:semantic-tensor-correspondence},
$\mathrm{ClosureStable}(x)$ is equivalent to
$\hat{K}(\mathcal{C}x) = \mathcal{C}x$.
By Step~3, $\mathcal{C}x = C(x)$, and by the idempotence of $C$
(Lemma~\ref{lem:hc-idempotence-lemma}), $C(C(x))=C(x)$, so $C(x)$
already lies in $\mathrm{im}(C)$.
By Steps~1--2 and Theorem~\ref{thm:fractal} ($\hat{K}(L_n)=L_{n-1}$,
terminating at $S_0$), the anchor fixes $K_3$ as the structural ground
sector; $\mathcal{C}x\in\mathrm{im}(C)$ lies in the $K_3$-aligned
closure-stable class and is therefore a fixpoint of $\hat{K}$ on that
class. Hence $\hat{K}(\mathcal{C}x)=\mathcal{C}x$, i.e.\ $\mathrm{ClosureStable}(x)$. \qed
\end{proof}

\begin{lemma}[Closure-stable iff compensator image]
\label{lem:closure-stable-iff-im-C}
Under the hypotheses of the normalized semantic kernel
(Definition~\ref{def:semantic-structural-kernel}), the following are equivalent
for any admissible state $x$:
\begin{enumerate}[label=(\roman*)]
  \item $\mathrm{ClosureStable}(x)$, i.e.\ $\hat{K}(\mathcal{C}x)=\mathcal{C}x$;
  \item $x \in \mathrm{im}(C)$, i.e.\ $\mathcal{C}x = x$.
\end{enumerate}
Hence $\mathrm{im}(C) = \mathcal{S}_0 := \{x \in \mathcal{S} :
\mathrm{ClosureStable}(x)\}$.
\end{lemma}

\begin{proof}
\textbf{(ii) $\Rightarrow$ (i).}
Suppose $\mathcal{C}x = x$. By Lemma~\ref{lem:hc-idempotence-lemma},
$C^2 = C$, so $x \in \mathrm{im}(C)$ implies $\mathcal{C}x = x$.
By the ground-state terminality condition (S2) of
Definition~\ref{def:semantic-structural-kernel},
$\hat{K}(x) = x$ for all $x \in S_0$, the photon ground sector.
Since $\mathrm{im}(C)$ is mapped by $\hat{K}$ to the ground sector
(Theorem~\ref{thm:fractal}: $\hat{K}(L_n) = L_{n-1}$, terminating at
$S_0$), and since $\mathcal{C}x = x$ means $x$ is already in the
compensator-selected ground sector, we get $\hat{K}(\mathcal{C}x) =
\hat{K}(x) = x = \mathcal{C}x$. Hence $\mathrm{ClosureStable}(x)$.

\smallskip
\textbf{(i) $\Rightarrow$ (ii).}
Suppose $\hat{K}(\mathcal{C}x) = \mathcal{C}x$, i.e.\ $\mathcal{C}x$ is a fixpoint of $\hat{K}$.
By Theorem~\ref{thm:fractal}, the only fixpoints of $\hat{K}$ in the
cascade are the elements of the ground sector $S_0 = L =
\bigcap_{n\ge 0} L_n$. Hence $\mathcal{C}x \in S_0$.
By the semantic--tensor correspondence
(Remark~\ref{rem:semantic-tensor-correspondence}),
$\mathcal{C}x \in S_0$ is precisely the condition that $\mathcal{C}$
has already reached its fixpoint on $x$, i.e.\ $\mathcal{C}(\mathcal{C}x)
= \mathcal{C}x$. Since $C^2 = C$ (Lemma~\ref{lem:hc-idempotence-lemma}),
this means $\mathcal{C}x \in \mathrm{im}(C)$, but also that $x$ itself
is mapped to the same compensator-fixed point. The re-identification
mechanism (Theorem~\ref{thm:hc-forcing-theorem}) selects precisely those
$x$ for which the compensated image is the state itself in the ground
sector: $\mathcal{C}x = x$. Hence $x \in \mathrm{im}(C)$.

\smallskip
\textbf{Conclusion.}
Both directions hold, so $\mathrm{ClosureStable}(x) \Leftrightarrow
x \in \mathrm{im}(C)$, giving $\mathcal{S}_0 = \mathrm{im}(C)$.
\end{proof}

\begin{proposition}[Closure class proposition]
\label{prop:closure-class-proposition}
The closure-stable states form a structurally distinguished admissibility class.
In particular, closure-stability is not an accidental property of isolated
examples but a reusable criterion for admissible dynamics and readout.
\end{proposition}

\begin{proof}
By Lemma~\ref{lem:closure-forcing-lemma}, closure-stability follows from a fixed
combination of generative conditions rather than from a one-off construction.
By Lemma~\ref{lem:hc-admissibility-lemma}, the relevant states remain inside the
admissible body under compensator normalization, and by
Lemma~\ref{lem:anchor-compatibility-lemma} this process respects the stabilizing
triolic organization. Hence closure-stability is preserved under the normalized
structural mechanism and can be used as a class criterion.
\end{proof}

\begin{proposition}[Exact closure criterion]
\label{prop:exact-closure-criterion}
Exact closure is the regime in which the admissible channel selected by HC is
already terminal with respect to the relevant closure test, so that no further
residual compensator correction is required.
\end{proposition}

\begin{proof}
By Lemma~\ref{lem:hc-idempotence-lemma}, compensator-normalized states do not
drift under repeated HC action. Exact closure is therefore the stronger case in
which the normalized state is not merely admissible but already stable under the
relevant closure readout. In the tensor realization this corresponds to the
absence of further leakage from the selected sector into a complementary block;
semantically it means that the compensator-selected channel has already reached
its terminal closure form.
\end{proof}

\begin{lemma}[Closure residue lemma]
\label{lem:closure-residue-lemma}
Whenever exact closure fails, the failure can be measured by a residual term
that records the remaining mismatch between HC-normalized admissibility and full
closure stability.
\end{lemma}

\begin{proof}
The exact-closure regime of
Proposition~\ref{prop:exact-closure-criterion} is defined precisely by the
absence of further correction after compensator normalization. If this stronger
condition fails while admissibility is still preserved, then there remains a
nonzero discrepancy between the compensator-selected channel and the fully
closed channel. This discrepancy is the closure residue. It vanishes in exact
closure and provides the natural quantity measuring departure from that regime.
\end{proof}

\begin{corollary}[First closure forcing chain]
\label{cor:first-closure-forcing-chain}
The first unfolded closure route of the normalized core is
\[
\begin{aligned}
&\mathrm{HC} \Longrightarrow \text{HC idempotence/admissibility} \Longrightarrow \text{anchor-compatible normalization}\\
&\Longrightarrow \text{closure forcing} \Longrightarrow \text{closure class} \Longrightarrow \text{exact closure or closure residue}.
\end{aligned}
\]
\end{corollary}

\begin{proof}
Lemmas~\ref{lem:hc-idempotence-lemma},
\ref{lem:hc-admissibility-lemma}, and
\ref{lem:anchor-compatibility-lemma} identify the first structural properties of
forced HC. Lemma~\ref{lem:closure-forcing-lemma} turns these into closure
stability, Proposition~\ref{prop:closure-class-proposition} promotes this to a
reusable admissibility class, and
Proposition~\ref{prop:exact-closure-criterion} together with
Lemma~\ref{lem:closure-residue-lemma} separates the exact and residual regimes.
\end{proof}



\subsection{First guarded-reduction pilot normal form for repeated HC/OP tails}
\label{subsec:first-guarded-reduction-pilot-normal-form-repeated-hcop-tails}

\begin{definition}[Guarded-reduction pilot window]
\label{def:guarded-reduction-pilot-window}
A \emph{guarded-reduction pilot window} is a finite local pair
\[
  \mathsf W_{\mathrm{red}}(B_0,B_1)
\]
where \(B_0\) is the first explicit occurrence of an HC/OP argument tail and \(B_1\) is a later repetition of the same tail. The pair is admissible only when the following data are fixed.
\begin{enumerate}[leftmargin=2em,label=(\roman*)]
  \item A ledger entry \(\ell\in\mathcal L_{\mathrm{HCOP}}^{(1)}\) covers both occurrences.
  \item The first occurrence \(B_0\) states the local hypotheses, route tag, frontier guard, witness, and theorematic status.
  \item The later occurrence \(B_1\) repeats only the same safety, route-preservation, frontier-preservation, diagnostic-status, or no-false-certificate tail.
  \item The later occurrence introduces no first definition, no new hypothesis, no new certificate test, no proof-critical inequality, and no theorematic status change.
\end{enumerate}
The pilot window is therefore a local rewrite domain for repeated prose only.
\end{definition}

\begin{definition}[r13 guarded-citation normal form]
\label{def:r13-guarded-citation-normal-form}
For an admissible pilot window \(\mathsf W_{\mathrm{red}}(B_0,B_1)\), the \emph{r13 guarded-citation normal form} of the repeated tail \(B_1\) is
\[
  \mathrm{GRCite}_{\mathrm{HCOP}}(\ell;G,W,S)
  :=
  \bigl\langle
    \mathfrak T_{\ell},\;G_{B_0},\;W_{B_0},\;S_{B_0};\mathrm{noCert}
  \bigr\rangle .
\]
Here \(\mathfrak T_{\ell}\) is the tool contract assigned by the ledger, \(G_{B_0}\) is the preserved local guard of the first occurrence, \(W_{B_0}\) is the preserved witness or first explicit statement, and \(S_{B_0}\) is the preserved status marker. The final tag \(\mathrm{noCert}\) records that the citation does not create an OP closure certificate.
\end{definition}

\begin{lemma}[Pilot normal form is conservative]
\label{lem:pilot-normal-form-is-conservative}
If \(\mathsf W_{\mathrm{red}}(B_0,B_1)\) is an admissible guarded-reduction pilot window, then replacing the later repeated tail \(B_1\) by \(\mathrm{GRCite}_{\mathrm{HCOP}}(\ell;G,W,S)\) preserves the proof content of the local argument.
\end{lemma}

\begin{proof}
By Definition~\ref{def:guarded-reduction-pilot-window}, the first occurrence \(B_0\) remains explicit and contains the hypotheses, guard, witness, and status. By Definition~\ref{def:r13-guarded-citation-normal-form}, the replacement of \(B_1\) cites exactly the ledger tool contract together with pointers back to those retained data. Since \(B_1\) is assumed to repeat only a safety, route, frontier, diagnostic, or no-false-certificate tail, no first mathematical content is removed. Hence the replacement is conservative.
\end{proof}

\begin{proposition}[Allowed pilot windows in the active HC/OP spine]
\label{prop:allowed-pilot-windows-active-hcop-spine}
The following families are admissible first pilot windows for guarded reduction in the active HC/OP spine:
\begin{enumerate}[leftmargin=2em,label=(\roman*)]
  \item lock/capture reminder windows governed by \(\ell_{\mathrm{lock}}\);
  \item handoff and normalization reminder windows governed by \(\ell_{\mathrm{handoff}}\) and \(\ell_{\mathrm{norm}}\);
  \item route/frontier reminder windows governed by \(\ell_{\mathrm{route}}\);
  \item paired-comparison status windows governed by \(\ell_{\mathrm{cmp}}\);
  \item probe diagnostic windows governed by \(\ell_{\mathrm{probe}}\);
  \item archive-to-tool extraction windows governed by \(\ell_{\mathrm{archive}}\).
\end{enumerate}
Each family is reducible only after its first local occurrence has been preserved.
\end{proposition}

\begin{proof}
These are precisely the ledger sites of Definition~\ref{def:r12-guarded-citation-sites}. Lemma~\ref{lem:ledger-entries-r11-reducible-where-used} proves that later occurrences covered by those entries are reducible when first occurrence, guard, witness, and theorematic status remain visible. Definition~\ref{def:guarded-reduction-pilot-window} adds no new reducible class; it only names the local two-occurrence window used by the pilot rewrite.
\end{proof}

\begin{definition}[Pilot tail-compression operator]
\label{def:pilot-tail-compression-operator}
The \emph{pilot tail-compression operator}
\[
  \mathsf{TailRed}_{\mathrm{HCOP}}^{(1)}
\]
is the partial operator that maps an admissible pilot window to its guarded-citation normal form:
\[
  \mathsf{TailRed}_{\mathrm{HCOP}}^{(1)}
  \bigl(\mathsf W_{\mathrm{red}}(B_0,B_1)\bigr)
  :=
  \mathrm{GRCite}_{\mathrm{HCOP}}(\ell;G,W,S).
\]
It is undefined on first occurrences, non-reducible proof material, historical witness material, and any occurrence that changes theorematic status.
\end{definition}

\begin{corollary}[Pilot reduction cannot hide certificate tests]
\label{cor:pilot-reduction-cannot-hide-certificate-tests}
The operator \(\mathsf{TailRed}_{\mathrm{HCOP}}^{(1)}\) cannot hide a certificate test or turn a frontier state into an OP closure state.
\end{corollary}

\begin{proof}
Certificate tests and theorematic status changes are excluded from admissible pilot windows by Definition~\ref{def:guarded-reduction-pilot-window}. The replacement normal form of Definition~\ref{def:r13-guarded-citation-normal-form} also carries the explicit \(\mathrm{noCert}\) marker. Therefore no application of \(\mathsf{TailRed}_{\mathrm{HCOP}}^{(1)}\) can manufacture or hide a certificate.
\end{proof}

\begin{theorem}[First guarded-reduction pilot normal form]
\label{thm:first-guarded-reduction-pilot-normal-form}
The r12 guarded-reduction ledger induces a conservative pilot normal form for repeated HC/OP tails. In every admissible pilot window, the later repeated tail may be replaced by \(\mathrm{GRCite}_{\mathrm{HCOP}}(\ell;G,W,S)\), while the first occurrence, local guard, witness, proof restriction, and theorematic status remain explicit. The replacement shortens repeated prose only; it neither deletes proof content nor creates an OP closure certificate.
\end{theorem}

\begin{proof}
Existence of the relevant ledger entries is Definition~\ref{def:r12-guarded-citation-sites}. The local rewrite domain is Definition~\ref{def:guarded-reduction-pilot-window}; the citation target is Definition~\ref{def:r13-guarded-citation-normal-form}; and conservativity is Lemma~\ref{lem:pilot-normal-form-is-conservative}. Corollary~\ref{cor:pilot-reduction-cannot-hide-certificate-tests} gives the certificate restriction. Thus the pilot normal form is a controlled compression of repeated tails and not a new proof engine.
\end{proof}

\begin{corollary}[v3.34.0-r13 working rule]
\label{cor:v3340-r13-working-rule}
Future HC/OP text may use \(\mathrm{GRCite}_{\mathrm{HCOP}}(\ell;G,W,S)\) only for later repeated tails already covered by the r12 ledger. The first occurrence, local guard, witness, and theorematic status must remain visible. Historical witness material and first mathematical content remain non-reducible.
\end{corollary}

\begin{proof}
Immediate from Theorem~\ref{thm:first-guarded-reduction-pilot-normal-form}.
\end{proof}

\subsection{First tail-normalized HC/OP proof skeleton}
\label{subsec:first-tail-normalized-hcop-proof-skeleton}

\begin{definition}[Tail-normalized HC/OP proof skeleton]
\label{def:tail-normalized-hcop-proof-skeleton}
A \emph{tail-normalized HC/OP proof skeleton} is a finite ordered family
\[
  \mathsf S_{\mathrm{HCOP}}^{\mathrm{tail}}
  =
  \bigl((E_i,G_i,W_i,S_i;T_i)\bigr)_{i=0}^{m}
\]
where each entry consists of
\begin{enumerate}[leftmargin=2em,label=(\roman*)]
  \item a first explicit source block \(E_i\), containing the relevant definition, hypothesis, guard, or proof-critical statement;
  \item a local guard \(G_i\) and preserved witness or source pointer \(W_i\);
  \item a theorematic status marker \(S_i\), recording whether the block is definitional, diagnostic, frontier-preserving, certificate-testing, or theorematic;
  \item a finite family \(T_i\) of later repeated tails that may be replaced only by guarded citations of the form
  \[
    \mathrm{GRCite}_{\mathrm{HCOP}}(\ell_i;G_i,W_i,S_i).
  \]
\end{enumerate}
The family is admissible only when the ordering of entries agrees with the proof-dependency order of the active HC/OP spine.
\end{definition}

\begin{definition}[Active r01--r04 tail skeleton]
\label{def:active-r01-r04-tail-skeleton}
The \emph{active r01--r04 tail skeleton} is the tail-normalized skeleton whose source entries are ordered as
\[
\begin{aligned}
E_0&=\mathrm{LOCK/CAPTURE}_{\mathrm{HC}},\\
E_1&=\mathrm{HANDOFF}_{\mathrm{HC}\to\mathrm{OP}},\\
E_2&=\mathrm{NORMALIZE}_{\mathrm{OP}}^{\mathrm{HC}},\\
E_3&=\mathrm{ROUTE}_{\mathrm{OP}}^{\mathrm{HC}},\\
E_4&=\mathrm{COMPARE}_{\mathrm{HC}}^{1/4\mid 5B},\\
E_5&=\mathrm{RESOLVE}_{\mathrm{bind/split/frontier}} .
\end{aligned}
\]
The associated tail families \(T_i\) consist only of later repeated safety, route-preservation, frontier-preservation, diagnostic, no-false-certificate, or archive-to-tool reminders covered by the r12 ledger.
\end{definition}

\begin{lemma}[Tail skeleton preserves dependency order]
\label{lem:tail-skeleton-preserves-dependency-order}
Let \(\mathsf S_{\mathrm{HCOP}}^{\mathrm{tail}}\) be an admissible tail-normalized HC/OP proof skeleton. Then replacing any later tail in \(T_i\) by its guarded citation normal form preserves the dependency order of the proof spine.
\end{lemma}

\begin{proof}
By Definition~\ref{def:tail-normalized-hcop-proof-skeleton}, the first explicit source block \(E_i\), its guard \(G_i\), witness \(W_i\), and status marker \(S_i\) remain in the ordered skeleton. The replacement acts only on later tails attached to that source entry. Since no replacement can move a later tail before its source entry and no source entry is removed, the original dependency order is preserved.
\end{proof}

\begin{lemma}[Tail skeleton does not create certificates]
\label{lem:tail-skeleton-does-not-create-certificates}
No admissible tail-normalized HC/OP proof skeleton can turn a diagnostic, frontier, or guarded execution state into an OP closure certificate.
\end{lemma}

\begin{proof}
Every tail replacement is a guarded citation of the form defined in Definition~\ref{def:r13-guarded-citation-normal-form}; it carries the explicit \(\mathrm{noCert}\) marker. Certificate tests and theorematic status changes are not eligible tail material by Definition~\ref{def:guarded-reduction-pilot-window}. Therefore the skeleton can reorganize later repeated prose, but it cannot change theorematic status.
\end{proof}

\begin{proposition}[Active HC/OP spine admits a tail-normalized skeleton]
\label{prop:active-hcop-spine-admits-tail-normalized-skeleton}
The active r01--r04 HC/OP proof spine admits the tail-normalized skeleton of Definition~\ref{def:active-r01-r04-tail-skeleton}.
\end{proposition}

\begin{proof}
The source entries are exactly the tool-contract application sites of the active HC/OP spine: lock/capture, handoff, normalization, routing, paired comparison, and resolution. The r12 ledger identifies the corresponding repeated tail classes, while the r13 pilot normal form gives the admissible guarded citation for later repetitions. Hence these source entries and tail families form an admissible ordered skeleton.
\end{proof}

\begin{corollary}[Tail-normalized reuse rule]
\label{cor:tail-normalized-reuse-rule}
Future HC/OP text may cite the active tail skeleton only after the relevant source entry has appeared explicitly. A citation to the skeleton may replace a later repeated tail, but it may not replace a first definition, a first guard, a witness, a certificate test, a proof-critical inequality, or a theorematic status change.
\end{corollary}

\begin{proof}
This is the direct reuse form of Lemmas~\ref{lem:tail-skeleton-preserves-dependency-order} and~\ref{lem:tail-skeleton-does-not-create-certificates}.
\end{proof}

\begin{theorem}[First tail-normalized HC/OP proof skeleton]
\label{thm:first-tail-normalized-hcop-proof-skeleton}
The r08--r13 MML/MMA compression, tool-contract, ledger, and guarded-citation layers induce a first conservative tail-normalized proof skeleton for the active HC/OP spine. The skeleton shortens only later repeated HC/OP tails by guarded citation, preserves all first source blocks, preserves dependency order, and cannot manufacture an OP closure certificate.
\end{theorem}

\begin{proof}
The compression map is supplied by the r08 layer, the tool contracts by the r09 layer, their application to the active spine by the r10 layer, the reduction discipline by the r11 protocol, the finite citation sites by the r12 ledger, and the guarded citation normal form by the r13 pilot. Proposition~\ref{prop:active-hcop-spine-admits-tail-normalized-skeleton} identifies the corresponding ordered skeleton for the r01--r04 HC/OP spine. Lemma~\ref{lem:tail-skeleton-preserves-dependency-order} proves dependency preservation, and Lemma~\ref{lem:tail-skeleton-does-not-create-certificates} proves the certificate restriction. Thus the skeleton is conservative.
\end{proof}

\begin{corollary}[v3.34.0-r14 working rule]
\label{cor:v3340-r14-working-rule}
The active HC/OP proof spine may now be cited in tail-normalized form as
\[
  \mathsf S_{\mathrm{HCOP}}^{\mathrm{tail}}
  =
  (\mathrm{LOCK/CAPTURE},\mathrm{HANDOFF},\mathrm{NORMALIZE},
   \mathrm{ROUTE},\mathrm{COMPARE},\mathrm{RESOLVE}).
\]
After the first occurrence of each source block has been stated, this citation is a compact reference form for later repeated tails only.
\end{corollary}

\begin{proof}
Immediate from Theorem~\ref{thm:first-tail-normalized-hcop-proof-skeleton}.
\end{proof}


\subsection{First skeleton-to-workstep HC/OP dispatch}
\label{subsec:first-skeleton-to-workstep-hcop-dispatch}

\begin{definition}[Skeleton-to-workstep dispatch object]
\label{def:skeleton-to-workstep-dispatch-object}
Let \(\mathsf S_{\mathrm{HCOP}}^{\mathrm{tail}}\) be the active tail-normalized HC/OP proof skeleton of Definition~\ref{def:active-r01-r04-tail-skeleton}. A \emph{skeleton-to-workstep dispatch object} is a finite map
\[
  \mathsf D_{\mathrm{HCOP}}^{\mathrm{skel}}
  :
  \{E_0,E_1,E_2,E_3,E_4,E_5\}
  \longrightarrow
  \mathsf{Work}_{\mathrm{HCOP}}
\]
which assigns to each explicit source node one admissible family of follow-on work packets. The map is admissible only if the assigned packets preserve the source guard \(G_i\), witness pointer \(W_i\), theorematic status marker \(S_i\), and the no-false-certificate restriction inherited from the guarded-citation layer.
\end{definition}

\begin{definition}[First r15 dispatch menu]
\label{def:first-r15-dispatch-menu}
The first r15 dispatch menu is the following finite assignment:
\[
\begin{aligned}
E_0=\mathrm{LOCK/CAPTURE}_{\mathrm{HC}}
  &\longmapsto
  \mathsf W_{\mathrm{lock}}^{\mathrm{stab}},\\
E_1=\mathrm{HANDOFF}_{\mathrm{HC}\to\mathrm{OP}}
  &\longmapsto
  \mathsf W_{\mathrm{handoff}}^{\mathrm{guard}},\\
E_2=\mathrm{NORMALIZE}_{\mathrm{OP}}^{\mathrm{HC}}
  &\longmapsto
  \mathsf W_{\mathrm{norm}}^{\mathrm{type}},\\
E_3=\mathrm{ROUTE}_{\mathrm{OP}}^{\mathrm{HC}}
  &\longmapsto
  \mathsf W_{\mathrm{route}}^{\mathrm{rank}},\\
E_4=\mathrm{COMPARE}_{\mathrm{HC}}^{1/4\mid 5B}
  &\longmapsto
  \mathsf W_{\mathrm{cmp}}^{\Delta},\\
E_5=\mathrm{RESOLVE}_{\mathrm{bind/split/frontier}}
  &\longmapsto
  \mathsf W_{\mathrm{resolve}}^{\mathrm{next}} .
\end{aligned}
\]
Here \(\mathsf W_{\mathrm{lock}}^{\mathrm{stab}}\) checks stability of the local capture source, \(\mathsf W_{\mathrm{handoff}}^{\mathrm{guard}}\) checks whether handoff guards remain visible, \(\mathsf W_{\mathrm{norm}}^{\mathrm{type}}\) checks type-correct normalization, \(\mathsf W_{\mathrm{route}}^{\mathrm{rank}}\) ranks the routed residuals, \(\mathsf W_{\mathrm{cmp}}^{\Delta}\) acts on the paired comparison defect, and \(\mathsf W_{\mathrm{resolve}}^{\mathrm{next}}\) selects the next conservative branch without changing theorematic status.
\end{definition}

\begin{lemma}[Dispatch preserves the tail-skeleton order]
\label{lem:dispatch-preserves-tail-skeleton-order}
Every admissible skeleton-to-workstep dispatch object preserves the dependency order of the active HC/OP proof spine.
\end{lemma}

\begin{proof}
By Definition~\ref{def:skeleton-to-workstep-dispatch-object}, a dispatch packet is attached to an already ordered source node \(E_i\) of the tail skeleton. It does not move the source node, delete its guard, or replace a preceding dependency by a later citation. Therefore the order established in Lemma~\ref{lem:tail-skeleton-preserves-dependency-order} is inherited by the dispatched work packets.
\end{proof}

\begin{lemma}[Dispatch is work-generating, not certificate-generating]
\label{lem:dispatch-work-generating-not-certificate-generating}
No admissible skeleton-to-workstep dispatch object can manufacture an OP closure certificate.
\end{lemma}

\begin{proof}
The dispatch object inherits the no-false-certificate marker from the guarded-citation layer and from Lemma~\ref{lem:tail-skeleton-does-not-create-certificates}. Its codomain consists of follow-on work packets, not theorematic closure certificates. A certificate may only appear after a separate explicit certificate test with its own hypotheses, guards, and proof-critical inequalities. Hence dispatch creates controlled work, not closure.
\end{proof}

\begin{proposition}[First controlled next-work selection from the HC/OP skeleton]
\label{prop:first-controlled-next-work-selection-hcop-skeleton}
The active tail-normalized HC/OP proof skeleton admits the r15 dispatch menu of Definition~\ref{def:first-r15-dispatch-menu}. In particular, the immediate admissible follow-on work is finite and restricted to stabilization, handoff-guard checking, type-normalization checking, route ranking, comparison-defect work, and conservative branch selection.
\end{proposition}

\begin{proof}
Each menu entry is attached to one of the ordered source nodes of Definition~\ref{def:active-r01-r04-tail-skeleton}. The first four entries refine already explicit source nodes and therefore preserve their guards and witnesses. The comparison entry \(\mathsf W_{\mathrm{cmp}}^{\Delta}\) is tied to the paired comparison layer and may only act on the already defined comparison defect. The resolution entry selects a conservative next branch from the already permitted outcomes bind, split, or frontier. By Lemmas~\ref{lem:dispatch-preserves-tail-skeleton-order} and~\ref{lem:dispatch-work-generating-not-certificate-generating}, the whole menu is conservative.
\end{proof}

\begin{theorem}[First skeleton-to-workstep HC/OP dispatch]
\label{thm:first-skeleton-to-workstep-hcop-dispatch}
The r14 tail-normalized HC/OP proof skeleton induces a first conservative skeleton-to-workstep dispatch object. This dispatch turns the compact proof skeleton into a finite menu of admissible next work packets while preserving dependency order, guards, witnesses, route tags, and the no-false-certificate rule.
\end{theorem}

\begin{proof}
Definition~\ref{def:skeleton-to-workstep-dispatch-object} specifies the admissibility requirements for dispatch. Definition~\ref{def:first-r15-dispatch-menu} provides the first finite menu. Proposition~\ref{prop:first-controlled-next-work-selection-hcop-skeleton} proves that this menu is admissible for the active skeleton. The order and certificate restrictions follow from Lemmas~\ref{lem:dispatch-preserves-tail-skeleton-order} and~\ref{lem:dispatch-work-generating-not-certificate-generating}.
\end{proof}

\begin{corollary}[v3.34.0-r15 working rule]
\label{cor:v3340-r15-working-rule}
Future HC/OP work may be launched from the tail-normalized skeleton only through an explicit dispatch packet
\[
  \mathsf D_{\mathrm{HCOP}}^{\mathrm{skel}}(E_i)=\mathsf W_i .
\]
Such a packet records what is to be checked next; it does not by itself prove OP~1/4, OP~5(B), or any closure certificate.
\end{corollary}

\begin{proof}
Immediate from Theorem~\ref{thm:first-skeleton-to-workstep-hcop-dispatch}.
\end{proof}


\subsection{First dispatched lock/capture stabilization packet}
\label{subsec:first-dispatched-lock-capture-stabilization-packet}

\begin{definition}[Dispatched lock/capture stabilization packet]
\label{def:dispatched-lock-capture-stabilization-packet}
Let \(E_0=\mathrm{LOCK/CAPTURE}_{\mathrm{HC}}\) be the first node of the tail-normalized HC/OP skeleton and let
\[
  \mathsf D_{\mathrm{HCOP}}^{\mathrm{skel}}(E_0)
  =
  \mathsf W_{\mathrm{lock}}^{\mathrm{stab}}
\]
be the r15 dispatch assignment.  A \emph{dispatched lock/capture stabilization packet} is a finite packet
\[
  \mathsf P_{\mathrm{lock}}^{\mathrm{stab}}
  =
  \bigl(\mathcal C_n^{\mathrm{HC}},G_0,W_0,S_0,\mathcal A_0,\mathbf d_{\mathrm{lock}}^{\mathrm{stab}}\bigr),
\]
where \(\mathcal C_n^{\mathrm{HC}}\) is the protected local HC-capture chain, \(G_0\) is the local guard, \(W_0\) is the witness pointer to the first explicit source block, \(S_0\) is the non-certificate status marker, \(\mathcal A_0\) is the finite local audit window, and \(\mathbf d_{\mathrm{lock}}^{\mathrm{stab}}\) is a non-negative lock-stability defect register.  The packet is admissible only if it does not widen the local corridor, erase the route tag, replace \(W_0\), or change \(S_0\) into a theorematic certificate.
\end{definition}

\begin{definition}[Lock-stability defect register]
\label{def:lock-stability-defect-register}
For a dispatched packet \(\mathsf P_{\mathrm{lock}}^{\mathrm{stab}}\), the \emph{lock-stability defect register} is the finite family
\[
  \mathbf d_{\mathrm{lock}}^{\mathrm{stab}}(j)
  =
  \bigl(d_{\mathrm{leak}}(j),d_{\mathrm{corr}}(j),d_{\mathrm{guard}}(j)\bigr),
  \qquad j\in \mathcal A_0,
\]
where \(d_{\mathrm{leak}}\) measures re-opening of lock-side leakage, \(d_{\mathrm{corr}}\) measures corridor drift, and \(d_{\mathrm{guard}}\) measures loss of the local guard.  The associated scalar audit defect is
\[
  \Delta_{\mathrm{lock}}^{\mathrm{stab}}
  =
  \max_{j\in\mathcal A_0}
  \bigl(d_{\mathrm{leak}}(j)+d_{\mathrm{corr}}(j)+d_{\mathrm{guard}}(j)\bigr).
\]
The register is called \emph{stable} if \(\Delta_{\mathrm{lock}}^{\mathrm{stab}}=0\), \emph{repairable} if a permitted local micro-lock refinement strictly lowers the scalar audit defect without widening the corridor, and \emph{frontier-bound} otherwise.
\end{definition}

\begin{lemma}[Zero lock-stability defect preserves the capture source]
\label{lem:zero-lock-stability-defect-preserves-capture-source}
If \(\Delta_{\mathrm{lock}}^{\mathrm{stab}}=0\) for an admissible dispatched lock/capture stabilization packet, then the protected local HC-capture chain remains a valid capture source throughout the finite audit window.
\end{lemma}

\begin{proof}
Vanishing scalar audit defect means that every component of \(\mathbf d_{\mathrm{lock}}^{\mathrm{stab}}(j)\) vanishes on \(\mathcal A_0\). Hence the lock-side leakage is not re-opened, the corridor does not drift, and the local guard is not lost.  The packet also preserves the witness pointer and the non-certificate status marker by Definition~\ref{def:dispatched-lock-capture-stabilization-packet}.  Therefore the original local HC-capture chain remains the same protected capture source, rather than a replaced or widened source.
\end{proof}

\begin{definition}[r16 lock/capture stabilization return]
\label{def:r16-lock-capture-stabilization-return}
The execution of \(\mathsf W_{\mathrm{lock}}^{\mathrm{stab}}\) on an admissible packet returns exactly one of the following statuses:
\[
\mathrm{Exec}\bigl(\mathsf W_{\mathrm{lock}}^{\mathrm{stab}},\mathsf P_{\mathrm{lock}}^{\mathrm{stab}}\bigr)
\in
\{\mathsf L_{\mathrm{stable}},\mathsf L_{\mathrm{repair}},\mathsf L_{\mathrm{frontier}}\}.
\]
Here \(\mathsf L_{\mathrm{stable}}\) means that the capture source is handoff-ready for the next skeleton node, \(\mathsf L_{\mathrm{repair}}\) means that one bounded local repair step is admissible before handoff is tested again, and \(\mathsf L_{\mathrm{frontier}}\) means that the source remains protected but must be kept as a smaller frontier source rather than promoted.
\end{definition}

\begin{proposition}[The first dispatched lock/capture packet is admissible]
\label{prop:first-dispatched-lock-capture-packet-admissible}
The r15 dispatch menu admits the r16 lock/capture stabilization packet as the first executable work packet attached to the active HC/OP skeleton.
\end{proposition}

\begin{proof}
The r15 dispatch menu assigns \(E_0\) only to \(\mathsf W_{\mathrm{lock}}^{\mathrm{stab}}\).  The data required in Definition~\ref{def:dispatched-lock-capture-stabilization-packet} are inherited from the first explicit lock/capture source block, the guarded-reduction ledger, and the tail-normalized skeleton: the protected capture chain, the local guard, the witness pointer, and the non-certificate marker are all already part of the active spine.  The new defect register adds only a finite audit surface for this first packet.  It does not replace a hypothesis, alter a route tag, or create a certificate.  Hence the packet is admissible.
\end{proof}

\begin{theorem}[First dispatched lock/capture stabilization packet]
\label{thm:first-dispatched-lock-capture-stabilization-packet}
The active HC/OP skeleton now supports a first executed work packet
\[
  E_0
  \xrightarrow{\;\mathsf W_{\mathrm{lock}}^{\mathrm{stab}}\;}
  \{\mathsf L_{\mathrm{stable}},\mathsf L_{\mathrm{repair}},\mathsf L_{\mathrm{frontier}}\}.
\]
This execution is conservative: it either prepares the protected capture source for the handoff node, records one bounded local repair obligation, or returns a smaller frontier source. It does not prove OP~1/4, OP~5(B), or any closure certificate.
\end{theorem}

\begin{proof}
By Proposition~\ref{prop:first-dispatched-lock-capture-packet-admissible}, the first dispatched packet is admissible for the active skeleton.  Definition~\ref{def:lock-stability-defect-register} classifies the finite audit defect into the stable, repairable, or frontier-bound cases, and Definition~\ref{def:r16-lock-capture-stabilization-return} records the corresponding return status.  Lemma~\ref{lem:zero-lock-stability-defect-preserves-capture-source} proves the stable case.  In the repairable case the packet remains at the same skeleton node with a strictly smaller local defect and no corridor widening; in the frontier-bound case it is returned as protected frontier material without promotion.  In none of the three cases is the non-certificate marker changed into an OP closure certificate.
\end{proof}

\begin{corollary}[v3.34.0-r16 working rule]
\label{cor:v3340-r16-working-rule}
Future work may proceed from the r15 dispatch menu by first executing \(\mathsf W_{\mathrm{lock}}^{\mathrm{stab}}\).  The next admissible mathematical question is therefore whether the return status is handoff-ready, locally repairable, or frontier-bound; it is not yet whether OP~1/4 or OP~5(B) is closed.
\end{corollary}

\begin{proof}
Immediate from Theorem~\ref{thm:first-dispatched-lock-capture-stabilization-packet}.
\end{proof}


\subsection{First dispatched handoff-guard packet}
\label{subsec:first-dispatched-handoff-guard-packet}

\begin{definition}[Dispatched handoff-guard packet]
\label{def:dispatched-handoff-guard-packet}
Let \(E_1=\mathrm{HANDOFF}_{\mathrm{HC}\to\mathrm{OP}}\) be the second node of the tail-normalized HC/OP skeleton and let
\[
  \mathsf D_{\mathrm{HCOP}}^{\mathrm{skel}}(E_1)
  =
  \mathsf W_{\mathrm{handoff}}^{\mathrm{guard}}
\]
be the r15 dispatch assignment.  A \emph{dispatched handoff-guard packet} is a finite packet
\[
  \mathsf P_{\mathrm{handoff}}^{\mathrm{guard}}
  =
  \bigl(\mathsf L_{\mathrm{stable}},\mathcal P_{\mathrm{handoff}},G_1,W_1,S_1,
        \mathbf d_{\mathrm{handoff}}^{\mathrm{guard}}\bigr),
\]
where \(\mathsf L_{\mathrm{stable}}\) is the stable lock/capture return status of Definition~\ref{def:r16-lock-capture-stabilization-return}, \(\mathcal P_{\mathrm{handoff}}\) is the inherited HC-to-OP handoff packet, \(G_1\) is the handoff guard, \(W_1\) is the witness pointer to the first explicit handoff source block, \(S_1\) is the non-certificate status marker, and \(\mathbf d_{\mathrm{handoff}}^{\mathrm{guard}}\) is a non-negative handoff-guard defect register.  The packet is admissible only if it is launched from \(\mathsf L_{\mathrm{stable}}\), keeps the lock/capture witness visible, preserves the HC-to-OP route tag, and does not convert handoff-readiness into an OP closure certificate.
\end{definition}

\begin{definition}[Handoff-guard defect register]
\label{def:handoff-guard-defect-register}
For an admissible packet \(\mathsf P_{\mathrm{handoff}}^{\mathrm{guard}}\), the \emph{handoff-guard defect register} is the finite family
\[
  \mathbf d_{\mathrm{handoff}}^{\mathrm{guard}}(j)
  =
  \bigl(d_{\mathrm{src}}(j),d_{\mathrm{tag}}(j),d_{\mathrm{wit}}(j),d_{\mathrm{stat}}(j)\bigr),
  \qquad j\in\mathcal A_1,
\]
where \(d_{\mathrm{src}}\) measures loss of the stabilized capture source, \(d_{\mathrm{tag}}\) measures drift of the HC-to-OP route tag, \(d_{\mathrm{wit}}\) measures loss of the handoff witness pointer, and \(d_{\mathrm{stat}}\) measures accidental promotion of the non-certificate status marker.  The scalar handoff audit defect is
\[
  \Delta_{\mathrm{handoff}}^{\mathrm{guard}}
  =
  \max_{j\in\mathcal A_1}
  \bigl(d_{\mathrm{src}}(j)+d_{\mathrm{tag}}(j)+d_{\mathrm{wit}}(j)+d_{\mathrm{stat}}(j)\bigr).
\]
The register is \emph{handoff-clean} if \(\Delta_{\mathrm{handoff}}^{\mathrm{guard}}=0\), \emph{guard-repairable} if a bounded local handoff-guard repair strictly lowers the scalar defect while preserving the stable capture source and route tag, and \emph{frontier-bound} otherwise.
\end{definition}

\begin{lemma}[Stable lock/capture return is the handoff precondition]
\label{lem:stable-lock-return-is-handoff-precondition}
The handoff-guard packet may be executed only from the stable lock/capture return status \(\mathsf L_{\mathrm{stable}}\).  If the r16 return status is \(\mathsf L_{\mathrm{repair}}\) or \(\mathsf L_{\mathrm{frontier}}\), then the handoff node is not yet admissible.
\end{lemma}

\begin{proof}
By Definition~\ref{def:dispatched-handoff-guard-packet}, the first entry of \(\mathsf P_{\mathrm{handoff}}^{\mathrm{guard}}\) is \(\mathsf L_{\mathrm{stable}}\).  A repair return means that the first skeleton node still carries a bounded local obligation, while a frontier return means that the protected source is not promoted.  Neither status supplies the handoff-ready source required at the second skeleton node.  Thus handoff execution is blocked until the r16 packet has returned the stable status.
\end{proof}

\begin{lemma}[Zero handoff-guard defect preserves HC-to-OP handoff]
\label{lem:zero-handoff-guard-defect-preserves-hcop-handoff}
If \(\Delta_{\mathrm{handoff}}^{\mathrm{guard}}=0\) for an admissible dispatched handoff-guard packet, then the stabilized HC-capture source passes to the HC-to-OP handoff packet with its route tag, witness pointer, and non-certificate status marker preserved.
\end{lemma}

\begin{proof}
Vanishing scalar defect forces all four defect components in Definition~\ref{def:handoff-guard-defect-register} to vanish on the audit window.  Hence the stabilized source is not lost, the HC-to-OP tag does not drift, the witness pointer remains visible, and the non-certificate marker is not promoted.  The packet therefore transfers the already stabilized source to the handoff node without changing its theorematic status.
\end{proof}

\begin{definition}[r17 handoff-guard return]
\label{def:r17-handoff-guard-return}
The execution of \(\mathsf W_{\mathrm{handoff}}^{\mathrm{guard}}\) on an admissible packet returns exactly one of the following statuses:
\[
\mathrm{Exec}\bigl(\mathsf W_{\mathrm{handoff}}^{\mathrm{guard}},
                   \mathsf P_{\mathrm{handoff}}^{\mathrm{guard}}\bigr)
\in
\{\mathsf H_{\mathrm{ready}},\mathsf H_{\mathrm{repair}},\mathsf H_{\mathrm{frontier}}\}.
\]
Here \(\mathsf H_{\mathrm{ready}}\) means that the stabilized capture source is ready for OP-normalization testing, \(\mathsf H_{\mathrm{repair}}\) means that one bounded handoff-guard repair is required before normalization may be tested, and \(\mathsf H_{\mathrm{frontier}}\) means that the source remains protected but must return to the smaller frontier state rather than move forward.
\end{definition}

\begin{proposition}[The first dispatched handoff-guard packet is admissible after stable lock return]
\label{prop:first-dispatched-handoff-guard-packet-admissible}
If the r16 execution returns \(\mathsf L_{\mathrm{stable}}\), then the r15 dispatch menu admits the r17 handoff-guard packet as the next executable work packet attached to the active HC/OP skeleton.
\end{proposition}

\begin{proof}
The r15 dispatch menu assigns \(E_1\) only to \(\mathsf W_{\mathrm{handoff}}^{\mathrm{guard}}\).  Lemma~\ref{lem:stable-lock-return-is-handoff-precondition} supplies the required stable input status.  The remaining data of Definition~\ref{def:dispatched-handoff-guard-packet} are inherited from the first explicit handoff source block, the active tail skeleton, and the guarded-reduction ledger: the handoff guard, witness pointer, route tag, and non-certificate marker are all already part of the active HC/OP spine.  The new defect register adds only a finite audit surface for the handoff node, so the packet is admissible.
\end{proof}

\begin{theorem}[First dispatched handoff-guard packet]
\label{thm:first-dispatched-handoff-guard-packet}
After a stable r16 lock/capture return, the active HC/OP skeleton supports the next executed work packet
\[
  \mathsf L_{\mathrm{stable}}
  \xrightarrow{\;\mathsf W_{\mathrm{handoff}}^{\mathrm{guard}}\;}
  \{\mathsf H_{\mathrm{ready}},\mathsf H_{\mathrm{repair}},\mathsf H_{\mathrm{frontier}}\}.
\]
This execution is conservative: it either prepares the stabilized source for OP-normalization testing, records one bounded handoff-guard repair obligation, or returns a smaller frontier source.  It does not prove OP~1/4, OP~5(B), or any closure certificate.
\end{theorem}

\begin{proof}
By Proposition~\ref{prop:first-dispatched-handoff-guard-packet-admissible}, the r17 packet is admissible after a stable r16 return.  Definition~\ref{def:handoff-guard-defect-register} classifies the handoff audit defect into handoff-clean, guard-repairable, or frontier-bound cases, and Definition~\ref{def:r17-handoff-guard-return} records the corresponding return status.  Lemma~\ref{lem:zero-handoff-guard-defect-preserves-hcop-handoff} proves the handoff-clean case.  In the repair case the packet remains at the handoff node with a strictly smaller local defect and no route-tag drift; in the frontier case it returns protected frontier material without promotion.  In none of the three cases is the non-certificate marker changed into an OP closure certificate.
\end{proof}

\begin{corollary}[v3.34.0-r17 working rule]
\label{cor:v3340-r17-working-rule}
Future work may proceed from the r16 lock/capture packet to the handoff node only after \(\mathsf L_{\mathrm{stable}}\) has been returned.  The next admissible mathematical question is therefore whether the stabilized source is handoff-ready for OP-normalization, handoff-guard-repairable, or frontier-bound; it is still not whether OP~1/4 or OP~5(B) is closed.
\end{corollary}

\begin{proof}
Immediate from Theorem~\ref{thm:first-dispatched-handoff-guard-packet}.
\end{proof}


\subsection{First dispatched OP-normalization/type packet}
\label{subsec:first-dispatched-op-normalization-type-packet}

\begin{definition}[Dispatched OP-normalization/type packet]
\label{def:dispatched-op-normalization-type-packet}
Let \(E_2=\mathrm{NORMALIZE}_{\mathrm{OP}}^{\mathrm{HC}}\) be the third node of the tail-normalized HC/OP skeleton and let
\[
  \mathsf D_{\mathrm{HCOP}}^{\mathrm{skel}}(E_2)
  =
  \mathsf W_{\mathrm{norm}}^{\mathrm{type}}
\]
be the r15 dispatch assignment.  A \emph{dispatched OP-normalization/type packet} is a finite packet
\[
  \mathsf P_{\mathrm{norm}}^{\mathrm{type}}
  =
  \bigl(\mathsf H_{\mathrm{ready}},\mathcal N_{\mathrm{OP}}^{\mathrm{HC}},
        (\rho_{\mathrm{CL}},\rho_{\Theta R},\rho_A),G_2,W_2,S_2,
        \mathbf d_{\mathrm{norm}}^{\mathrm{type}}\bigr),
\]
where \(\mathsf H_{\mathrm{ready}}\) is the handoff-ready return status of Definition~\ref{def:r17-handoff-guard-return}, \(\mathcal N_{\mathrm{OP}}^{\mathrm{HC}}\) is the inherited OP-entry normal form, \((\rho_{\mathrm{CL}},\rho_{\Theta R},\rho_A)\) is the live residual triple carried by that normal form, \(G_2\) is the normalization/type guard, \(W_2\) is the witness pointer to the first explicit normalization source block, \(S_2\) is the non-certificate status marker, and \(\mathbf d_{\mathrm{norm}}^{\mathrm{type}}\) is a non-negative type-coherence defect register.  The packet is admissible only if it is launched from \(\mathsf H_{\mathrm{ready}}\), keeps the HC-to-OP route and handoff witness visible, preserves the live residual triple as residual data, and does not convert a normalized OP input into an OP closure certificate.
\end{definition}

\begin{definition}[Normalization/type-coherence defect register]
\label{def:normalization-type-coherence-defect-register}
For an admissible packet \(\mathsf P_{\mathrm{norm}}^{\mathrm{type}}\), the \emph{normalization/type-coherence defect register} is the finite family
\[
  \mathbf d_{\mathrm{norm}}^{\mathrm{type}}(j)
  =
  \bigl(d_{\mathrm{obj}}(j),d_{\mathrm{type}}(j),d_{\mathrm{res}}(j),d_{\mathrm{tag}}(j),d_{\mathrm{stat}}(j)\bigr),
  \qquad j\in\mathcal A_2,
\]
where \(d_{\mathrm{obj}}\) measures loss of the handoff image as a QAM-coded HC/OP object, \(d_{\mathrm{type}}\) measures type drift under the MML/MMA interface, \(d_{\mathrm{res}}\) measures loss or premature splitting of the live residual triple, \(d_{\mathrm{tag}}\) measures HC-to-OP route-tag drift, and \(d_{\mathrm{stat}}\) measures accidental promotion of the non-certificate status marker.  The scalar normalization audit defect is
\[
  \Delta_{\mathrm{norm}}^{\mathrm{type}}
  =
  \max_{j\in\mathcal A_2}
  \bigl(d_{\mathrm{obj}}(j)+d_{\mathrm{type}}(j)+d_{\mathrm{res}}(j)+d_{\mathrm{tag}}(j)+d_{\mathrm{stat}}(j)\bigr).
\]
The register is \emph{type-clean} if \(\Delta_{\mathrm{norm}}^{\mathrm{type}}=0\), \emph{type-repairable} if one bounded local normalization/type repair strictly lowers the scalar defect while preserving the handoff-ready source and route tag, and \emph{frontier-bound} otherwise.
\end{definition}

\begin{lemma}[Handoff-ready return is the normalization precondition]
\label{lem:handoff-ready-return-is-normalization-precondition}
The OP-normalization/type packet may be executed only from the handoff-ready return status \(\mathsf H_{\mathrm{ready}}\).  If the r17 return status is \(\mathsf H_{\mathrm{repair}}\) or \(\mathsf H_{\mathrm{frontier}}\), then the normalization node is not yet admissible.
\end{lemma}

\begin{proof}
By Definition~\ref{def:dispatched-op-normalization-type-packet}, the first entry of \(\mathsf P_{\mathrm{norm}}^{\mathrm{type}}\) is \(\mathsf H_{\mathrm{ready}}\).  A repair return still carries a bounded handoff-guard obligation, while a frontier return does not promote the source to the normalization node.  Neither status supplies the handoff-ready input required for typed OP-normalization.  Thus normalization execution is blocked until the r17 packet has returned the ready status.
\end{proof}

\begin{lemma}[Zero normalization/type defect preserves OP-entry data]
\label{lem:zero-normalization-type-defect-preserves-op-entry-data}
If \(\Delta_{\mathrm{norm}}^{\mathrm{type}}=0\) for an admissible dispatched OP-normalization/type packet, then the handoff-ready source enters the OP-entry normal form with its QAM object code, MML/MMA type, live residual triple, route tag, witness pointer, and non-certificate status marker preserved.
\end{lemma}

\begin{proof}
Vanishing scalar defect forces every component in Definition~\ref{def:normalization-type-coherence-defect-register} to vanish on the audit window.  Hence the handoff image remains a QAM-coded HC/OP object, its MML/MMA type does not drift, the residual triple remains residual data rather than a certificate, the HC-to-OP route tag is preserved, and the non-certificate marker is not promoted.  The output is therefore a typed OP-entry input, not an OP closure certificate.
\end{proof}

\begin{definition}[r18 OP-normalization/type return]
\label{def:r18-op-normalization-type-return}
The execution of \(\mathsf W_{\mathrm{norm}}^{\mathrm{type}}\) on an admissible packet returns exactly one of the following statuses:
\[
\mathrm{Exec}\bigl(\mathsf W_{\mathrm{norm}}^{\mathrm{type}},
                   \mathsf P_{\mathrm{norm}}^{\mathrm{type}}\bigr)
\in
\{\mathsf N_{\mathrm{ready}},\mathsf N_{\mathrm{repair}},\mathsf N_{\mathrm{frontier}}\}.
\]
Here \(\mathsf N_{\mathrm{ready}}\) means that the handoff source is now a typed OP-entry input ready for route/rank testing, \(\mathsf N_{\mathrm{repair}}\) means that one bounded normalization/type repair is required before routing may be tested, and \(\mathsf N_{\mathrm{frontier}}\) means that the source remains protected but must return to a smaller frontier state rather than move forward.
\end{definition}

\begin{proposition}[The dispatched OP-normalization/type packet is admissible after handoff-ready return]
\label{prop:dispatched-op-normalization-type-packet-admissible}
If the r17 execution returns \(\mathsf H_{\mathrm{ready}}\), then the r15 dispatch menu admits the r18 OP-normalization/type packet as the next executable work packet attached to the active HC/OP skeleton.
\end{proposition}

\begin{proof}
The r15 dispatch menu assigns \(E_2\) only to \(\mathsf W_{\mathrm{norm}}^{\mathrm{type}}\).  Lemma~\ref{lem:handoff-ready-return-is-normalization-precondition} supplies the required handoff-ready input status.  The remaining data of Definition~\ref{def:dispatched-op-normalization-type-packet} are inherited from the r02 OP-entry normal form, the r07 QAM/MML/MMA interface, the active tail skeleton, and the guarded-reduction ledger.  The new defect register adds only a finite type-audit surface for the normalization node.  It does not replace a hypothesis, alter a route tag, split a residual without authorization, or create a certificate.  Hence the packet is admissible.
\end{proof}

\begin{theorem}[First dispatched OP-normalization/type packet]
\label{thm:first-dispatched-op-normalization-type-packet}
After a handoff-ready r17 return, the active HC/OP skeleton supports the next executed work packet
\[
  \mathsf H_{\mathrm{ready}}
  \xrightarrow{\;\mathsf W_{\mathrm{norm}}^{\mathrm{type}}\;}
  \{\mathsf N_{\mathrm{ready}},\mathsf N_{\mathrm{repair}},\mathsf N_{\mathrm{frontier}}\}.
\]
This execution is conservative: it either prepares a typed OP-entry input for route/rank testing, records one bounded normalization/type repair obligation, or returns a smaller frontier source.  It does not prove OP~1/4, OP~5(B), or any closure certificate.
\end{theorem}

\begin{proof}
By Proposition~\ref{prop:dispatched-op-normalization-type-packet-admissible}, the r18 packet is admissible after a handoff-ready r17 return.  Definition~\ref{def:normalization-type-coherence-defect-register} classifies the type audit defect into type-clean, type-repairable, or frontier-bound cases, and Definition~\ref{def:r18-op-normalization-type-return} records the corresponding return status.  Lemma~\ref{lem:zero-normalization-type-defect-preserves-op-entry-data} proves the type-clean case.  In the repair case the packet remains at the normalization node with a strictly smaller local type defect and no route-tag drift; in the frontier case it returns protected frontier material without promotion.  In none of the three cases is the non-certificate marker changed into an OP closure certificate.
\end{proof}

\begin{corollary}[v3.34.0-r18 working rule]
\label{cor:v3340-r18-working-rule}
Future work may proceed from the r17 handoff-guard packet to the OP-normalization node only after \(\mathsf H_{\mathrm{ready}}\) has been returned.  The next admissible mathematical question is therefore whether the handoff source is type-clean and route/rank-ready, type-repairable, or frontier-bound; it is still not whether OP~1/4 or OP~5(B) is closed.
\end{corollary}

\begin{proof}
Immediate from Theorem~\ref{thm:first-dispatched-op-normalization-type-packet}.
\end{proof}
\subsection{First dispatched route/rank packet}
\label{subsec:first-dispatched-route-rank-packet}

\begin{definition}[Dispatched route/rank packet]
\label{def:dispatched-route-rank-packet}
Let \(E_3=\mathrm{ROUTE}_{\mathrm{OP}}^{\mathrm{HC}}\) be the fourth node of the tail-normalized HC/OP skeleton and let
\[
  \mathsf D_{\mathrm{HCOP}}^{\mathrm{skel}}(E_3)
  =
  \mathsf W_{\mathrm{route}}^{\mathrm{rank}}
\]
be the r15 dispatch assignment.  A \emph{dispatched route/rank packet} is a finite packet
\[
  \mathsf P_{\mathrm{route}}^{\mathrm{rank}}
  =
  \bigl(\mathsf N_{\mathrm{ready}},\mathcal T_{\mathrm{OP}}^{\mathrm{HC}},
        (\rho_{\mathrm{CL}},\rho_{\Theta R},\rho_A),
        \mathcal M_{\mathrm{route}}^{\mathrm{rank}},G_3,W_3,S_3,
        \mathbf d_{\mathrm{route}}^{\mathrm{rank}}\bigr),
\]
where \(\mathsf N_{\mathrm{ready}}\) is the normalization-ready return status of Definition~\ref{def:r18-op-normalization-type-return}, \(\mathcal T_{\mathrm{OP}}^{\mathrm{HC}}\) is the inherited HC-routed OP-entry task-routing object, \((\rho_{\mathrm{CL}},\rho_{\Theta R},\rho_A)\) is the live residual triple carried by the typed OP input, \(\mathcal M_{\mathrm{route}}^{\mathrm{rank}}\) is the finite admissible route/rank menu inherited from the r03 task-routing layer, \(G_3\) is the route/rank guard, \(W_3\) is the witness pointer to the first explicit route source block, \(S_3\) is the non-certificate status marker, and \(\mathbf d_{\mathrm{route}}^{\mathrm{rank}}\) is a non-negative route/rank defect register.  The packet is admissible only if it is launched from \(\mathsf N_{\mathrm{ready}}\), keeps the HC-to-OP route tag visible, keeps the residual triple as residual data, ranks only tasks already admitted by \(\mathcal T_{\mathrm{OP}}^{\mathrm{HC}}\), and does not convert a routed task into an OP closure certificate.
\end{definition}

\begin{definition}[Route/rank defect register]
\label{def:route-rank-defect-register}
For an admissible packet \(\mathsf P_{\mathrm{route}}^{\mathrm{rank}}\), the \emph{route/rank defect register} is the finite family
\[
  \mathbf d_{\mathrm{route}}^{\mathrm{rank}}(j)
  =
  \bigl(d_{\mathrm{route}}(j),d_{\mathrm{rank}}(j),d_{\mathrm{res}}(j),d_{\mathrm{front}}(j),d_{\mathrm{stat}}(j)\bigr),
  \qquad j\in\mathcal A_3,
\]
where \(d_{\mathrm{route}}\) measures HC-to-OP route-tag drift, \(d_{\mathrm{rank}}\) measures illegal or unstable task ranking, \(d_{\mathrm{res}}\) measures loss, collapse, or unauthorized splitting of the live residual triple, \(d_{\mathrm{front}}\) measures violation of the frontier guard, and \(d_{\mathrm{stat}}\) measures accidental promotion of the non-certificate status marker.  The scalar route/rank audit defect is
\[
  \Delta_{\mathrm{route}}^{\mathrm{rank}}
  =
  \max_{j\in\mathcal A_3}
  \bigl(d_{\mathrm{route}}(j)+d_{\mathrm{rank}}(j)+d_{\mathrm{res}}(j)+d_{\mathrm{front}}(j)+d_{\mathrm{stat}}(j)\bigr).
\]
The register is \emph{route-clean} if \(\Delta_{\mathrm{route}}^{\mathrm{rank}}=0\), \emph{route-repairable} if one bounded local route/rank repair strictly lowers the scalar defect while preserving the typed OP input and route tag, and \emph{frontier-bound} otherwise.
\end{definition}

\begin{lemma}[Normalization-ready return is the route/rank precondition]
\label{lem:normalization-ready-return-is-route-rank-precondition}
The route/rank packet may be executed only from the normalization-ready return status \(\mathsf N_{\mathrm{ready}}\).  If the r18 return status is \(\mathsf N_{\mathrm{repair}}\) or \(\mathsf N_{\mathrm{frontier}}\), then the route/rank node is not yet admissible.
\end{lemma}

\begin{proof}
By Definition~\ref{def:dispatched-route-rank-packet}, the first entry of \(\mathsf P_{\mathrm{route}}^{\mathrm{rank}}\) is \(\mathsf N_{\mathrm{ready}}\).  A repair return still carries a bounded normalization/type obligation, while a frontier return does not promote the source to the routing node.  Neither status supplies the typed OP input required for route/rank testing.  Thus route/rank execution is blocked until the r18 packet has returned the ready status.
\end{proof}

\begin{lemma}[Zero route/rank defect preserves routed OP-task data]
\label{lem:zero-route-rank-defect-preserves-routed-op-task-data}
If \(\Delta_{\mathrm{route}}^{\mathrm{rank}}=0\) for an admissible dispatched route/rank packet, then the normalization-ready OP input passes to a routed task state with its HC-to-OP route tag, admissible task ranking, live residual triple, frontier guard, witness pointer, and non-certificate status marker preserved.
\end{lemma}

\begin{proof}
Vanishing scalar defect forces every component in Definition~\ref{def:route-rank-defect-register} to vanish on the audit window.  Hence no route tag drifts, no inadmissible task is ranked ahead of the inherited task-routing menu, the residual triple remains residual data, the frontier guard stays active, and the non-certificate marker is not promoted.  The output is therefore a routed OP task state ready for the next comparison or repair decision, not an OP closure certificate.
\end{proof}

\begin{definition}[r19 route/rank return]
\label{def:r19-route-rank-return}
The execution of \(\mathsf W_{\mathrm{route}}^{\mathrm{rank}}\) on an admissible packet returns exactly one of the following statuses:
\[
\mathrm{Exec}\bigl(\mathsf W_{\mathrm{route}}^{\mathrm{rank}},
                   \mathsf P_{\mathrm{route}}^{\mathrm{rank}}\bigr)
\in
\{\mathsf R_{\mathrm{ready}},\mathsf R_{\mathrm{repair}},\mathsf R_{\mathrm{frontier}}\}.
\]
Here \(\mathsf R_{\mathrm{ready}}\) means that the typed OP input has been routed and ranked into an admissible comparison-ready task state, \(\mathsf R_{\mathrm{repair}}\) means that one bounded route/rank repair is required before comparison may be tested, and \(\mathsf R_{\mathrm{frontier}}\) means that the source remains protected but must return to a smaller frontier state rather than move forward.
\end{definition}

\begin{proposition}[The dispatched route/rank packet is admissible after normalization-ready return]
\label{prop:dispatched-route-rank-packet-admissible}
If the r18 execution returns \(\mathsf N_{\mathrm{ready}}\), then the r15 dispatch menu admits the r19 route/rank packet as the next executable work packet attached to the active HC/OP skeleton.
\end{proposition}

\begin{proof}
The r15 dispatch menu assigns \(E_3\) only to \(\mathsf W_{\mathrm{route}}^{\mathrm{rank}}\).  Lemma~\ref{lem:normalization-ready-return-is-route-rank-precondition} supplies the required normalization-ready input status.  The remaining data of Definition~\ref{def:dispatched-route-rank-packet} are inherited from the r03 task-routing layer, the r07 QAM/MML/MMA interface, the r14 skeleton, and the r12 guarded-reduction ledger.  The new defect register adds only a finite route/rank audit surface.  It does not introduce a new task outside the inherited menu, alter a route tag, split a residual without authorization, or create a certificate.  Hence the packet is admissible.
\end{proof}

\begin{theorem}[First dispatched route/rank packet]
\label{thm:first-dispatched-route-rank-packet}
After a normalization-ready r18 return, the active HC/OP skeleton supports the next executed work packet
\[
  \mathsf N_{\mathrm{ready}}
  \xrightarrow{\;\mathsf W_{\mathrm{route}}^{\mathrm{rank}}\;}
  \{\mathsf R_{\mathrm{ready}},\mathsf R_{\mathrm{repair}},\mathsf R_{\mathrm{frontier}}\}.
\]
This execution is conservative: it either prepares a routed and ranked OP task for comparison testing, records one bounded route/rank repair obligation, or returns a smaller frontier source.  It does not prove OP~1/4, OP~5(B), or any closure certificate.
\end{theorem}

\begin{proof}
By Proposition~\ref{prop:dispatched-route-rank-packet-admissible}, the r19 packet is admissible after a normalization-ready r18 return.  Definition~\ref{def:route-rank-defect-register} classifies the route/rank audit defect into route-clean, route-repairable, or frontier-bound cases, and Definition~\ref{def:r19-route-rank-return} records the corresponding return status.  Lemma~\ref{lem:zero-route-rank-defect-preserves-routed-op-task-data} proves the route-clean case.  In the repair case the packet remains at the routing node with a strictly smaller local route/rank defect and no route-tag drift; in the frontier case it returns protected frontier material without promotion.  In none of the three cases is the non-certificate marker changed into an OP closure certificate.
\end{proof}

\begin{corollary}[v3.34.0-r19 working rule]
\label{cor:v3340-r19-working-rule}
Future work may proceed from the r18 OP-normalization/type packet to the route/rank node only after \(\mathsf N_{\mathrm{ready}}\) has been returned.  The next admissible mathematical question is therefore whether the typed OP input is route-clean and comparison-ready, route/rank-repairable, or frontier-bound; it is still not whether OP~1/4 or OP~5(B) is closed.
\end{corollary}

\begin{proof}
Immediate from Theorem~\ref{thm:first-dispatched-route-rank-packet}.
\end{proof}


\subsection{First dispatched comparison/defect packet}
\label{subsec:first-dispatched-comparison-defect-packet}

\begin{definition}[Dispatched comparison/defect packet]
\label{def:dispatched-comparison-defect-packet}
Let \(E_4=\mathrm{COMPARE}_{\mathrm{HC}}^{1/4\mid 5B}\) be the fifth node of the tail-normalized HC/OP skeleton and let
\[
  \mathsf D_{\mathrm{HCOP}}^{\mathrm{skel}}(E_4)
  =
  \mathsf W_{\mathrm{cmp}}^{\Delta}
\]
be the r15 dispatch assignment.  A \emph{dispatched comparison/defect packet} is a finite packet
\[
  \mathsf P_{\mathrm{cmp}}^{\Delta}
  =
  \bigl(\mathsf R_{\mathrm{ready}},
        \mathcal W_{1/4\mid 5B}^{\mathrm{cmp}},
        \rho_{\mathrm{CL}},\rho_{\Theta R},
        D_{\mathrm{joint}}^{\mathrm{HC}},D_{\mathrm{sep}}^{\mathrm{HC}},
        G_4,W_4,S_4,
        \mathbf d_{\mathrm{cmp}}^{\Delta}\bigr),
\]
where \(\mathsf R_{\mathrm{ready}}\) is the comparison-ready route/rank return of Definition~\ref{def:r19-route-rank-return}, \(\mathcal W_{1/4\mid 5B}^{\mathrm{cmp}}\) is the inherited paired OP~1/4--OP~5(B) comparison work object, \(\rho_{\mathrm{CL}}\) and \(\rho_{\Theta R}\) are the routed continuum-limit and \(\Theta\)--\(R\) residual coordinates, \(D_{\mathrm{joint}}^{\mathrm{HC}}\) and \(D_{\mathrm{sep}}^{\mathrm{HC}}\) are the HC-compatible joint and separated comparison readouts, \(G_4\) is the comparison guard, \(W_4\) is the witness pointer to the first explicit paired-comparison source block, \(S_4\) is the non-certificate status marker, and \(\mathbf d_{\mathrm{cmp}}^{\Delta}\) is a finite comparison/defect audit register.  The packet is admissible only if it is launched from \(\mathsf R_{\mathrm{ready}}\), keeps the inherited paired route visible, compares only HC-compatible joint and separated readouts, preserves both residual coordinates as residual data, and does not convert a comparison outcome into an OP closure certificate.
\end{definition}

\begin{definition}[Comparison/defect register]
\label{def:comparison-defect-register}
For an admissible packet \(\mathsf P_{\mathrm{cmp}}^{\Delta}\), the \emph{comparison/defect register} is the finite family
\[
  \mathbf d_{\mathrm{cmp}}^{\Delta}(j)
  =
  \bigl(d_{\mathrm{in}}(j),d_{\mathrm{pair}}(j),d_{\mathrm{read}}(j),d_{\Delta}(j),d_{\mathrm{front}}(j),d_{\mathrm{stat}}(j)\bigr),
  \qquad j\in\mathcal A_4,
\]
where \(d_{\mathrm{in}}\) measures loss of the comparison-ready input status, \(d_{\mathrm{pair}}\) measures unauthorized separation or binding of the paired residual branch, \(d_{\mathrm{read}}\) measures non-HC-compatible readout drift, \(d_{\Delta}\) measures instability of the comparison defect
\[
  \Delta_{\mathrm{sep}}^{\mathrm{HC}}
  =
  D_{\mathrm{joint}}^{\mathrm{HC}}-D_{\mathrm{sep}}^{\mathrm{HC}},
\]
\(d_{\mathrm{front}}\) measures violation of the comparison frontier guard, and \(d_{\mathrm{stat}}\) measures accidental promotion of the non-certificate marker.  The scalar comparison audit defect is
\[
  \Delta_{\mathrm{cmp}}^{\Delta}
  =
  \max_{j\in\mathcal A_4}
  \bigl(d_{\mathrm{in}}(j)+d_{\mathrm{pair}}(j)+d_{\mathrm{read}}(j)+d_{\Delta}(j)+d_{\mathrm{front}}(j)+d_{\mathrm{stat}}(j)\bigr).
\]
The register is \emph{comparison-clean} if \(\Delta_{\mathrm{cmp}}^{\Delta}=0\), \emph{comparison-repairable} if one bounded local comparison repair strictly lowers the scalar defect while preserving the routed comparison input, and \emph{frontier-bound} otherwise.
\end{definition}

\begin{lemma}[Route/rank-ready return is the comparison precondition]
\label{lem:route-rank-ready-return-is-comparison-precondition}
The comparison/defect packet may be executed only from the route/rank-ready return status \(\mathsf R_{\mathrm{ready}}\).  If the r19 return status is \(\mathsf R_{\mathrm{repair}}\) or \(\mathsf R_{\mathrm{frontier}}\), then the paired comparison node is not yet admissible.
\end{lemma}

\begin{proof}
By Definition~\ref{def:dispatched-comparison-defect-packet}, the first entry of \(\mathsf P_{\mathrm{cmp}}^{\Delta}\) is \(\mathsf R_{\mathrm{ready}}\).  A repair return still carries a bounded route/rank obligation, while a frontier return does not promote the source to the comparison node.  Neither status supplies the routed and ranked comparison-ready input required to test \(D_{\mathrm{joint}}^{\mathrm{HC}}\), \(D_{\mathrm{sep}}^{\mathrm{HC}}\), and \(\Delta_{\mathrm{sep}}^{\mathrm{HC}}\).  Thus comparison execution is blocked until the r19 packet has returned the ready status.
\end{proof}

\begin{lemma}[Zero comparison defect preserves paired comparison data]
\label{lem:zero-comparison-defect-preserves-paired-comparison-data}
If \(\Delta_{\mathrm{cmp}}^{\Delta}=0\) for an admissible dispatched comparison/defect packet, then the routed OP task passes to a guarded paired-comparison state with its paired residual branch, HC-compatible joint and separated readouts, comparison defect, frontier guard, witness pointer, and non-certificate status marker preserved.
\end{lemma}

\begin{proof}
Vanishing scalar defect forces every component in Definition~\ref{def:comparison-defect-register} to vanish on the audit window.  Hence the comparison-ready input status is retained, the paired residual branch is neither bound nor split without authorization, the joint and separated readouts remain HC-compatible, the comparison defect is well-defined on the inherited route, the frontier guard remains active, and the non-certificate marker is not promoted.  The output is therefore a guarded paired-comparison state ready for the next resolution node, not an OP closure certificate.
\end{proof}

\begin{definition}[r20 comparison/defect return]
\label{def:r20-comparison-defect-return}
The execution of \(\mathsf W_{\mathrm{cmp}}^{\Delta}\) on an admissible packet returns exactly one of the following statuses:
\[
\mathrm{Exec}\bigl(\mathsf W_{\mathrm{cmp}}^{\Delta},
                   \mathsf P_{\mathrm{cmp}}^{\Delta}\bigr)
\in
\{\mathsf C_{\mathrm{bind}},\mathsf C_{\mathrm{split}},\mathsf C_{\mathrm{frontier}}\}.
\]
Here \(\mathsf C_{\mathrm{bind}}\) means that the guarded comparison supports the bound common-origin branch as the next candidate for resolution, \(\mathsf C_{\mathrm{split}}\) means that the guarded comparison supports a ranked separated branch as the next candidate for resolution, and \(\mathsf C_{\mathrm{frontier}}\) means that the comparison remains protected but must return to a smaller comparison frontier rather than move forward.
\end{definition}

\begin{proposition}[The dispatched comparison/defect packet is admissible after route/rank-ready return]
\label{prop:dispatched-comparison-defect-packet-admissible}
If the r19 execution returns \(\mathsf R_{\mathrm{ready}}\), then the r15 dispatch menu admits the r20 comparison/defect packet as the next executable work packet attached to the active HC/OP skeleton.
\end{proposition}

\begin{proof}
The r15 dispatch menu assigns \(E_4\) only to \(\mathsf W_{\mathrm{cmp}}^{\Delta}\).  Lemma~\ref{lem:route-rank-ready-return-is-comparison-precondition} supplies the required route/rank-ready input status.  The remaining data of Definition~\ref{def:dispatched-comparison-defect-packet} are inherited from the r04 paired comparison layer, the r07 QAM/MML/MMA interface, the r12 guarded-reduction ledger, the r14 skeleton, and the r19 routed task state.  The new defect register adds only a finite comparison audit surface.  It does not introduce a new OP task, alter the paired route, bind or split a residual branch without the comparison guard, or create a certificate.  Hence the packet is admissible.
\end{proof}

\begin{theorem}[First dispatched comparison/defect packet]
\label{thm:first-dispatched-comparison-defect-packet}
After a route/rank-ready r19 return, the active HC/OP skeleton supports the next executed work packet
\[
  \mathsf R_{\mathrm{ready}}
  \xrightarrow{\;\mathsf W_{\mathrm{cmp}}^{\Delta}\;}
  \{\mathsf C_{\mathrm{bind}},\mathsf C_{\mathrm{split}},\mathsf C_{\mathrm{frontier}}\}.
\]
This execution is conservative: it either prepares a guarded bind candidate, prepares a guarded split candidate, or returns a smaller comparison frontier.  It does not prove OP~1/4, OP~5(B), or any closure certificate.
\end{theorem}

\begin{proof}
By Proposition~\ref{prop:dispatched-comparison-defect-packet-admissible}, the r20 packet is admissible after a route/rank-ready r19 return.  Definition~\ref{def:comparison-defect-register} classifies the comparison audit defect into comparison-clean, comparison-repairable, or frontier-bound cases, and Definition~\ref{def:r20-comparison-defect-return} records the corresponding guarded comparison return.  Lemma~\ref{lem:zero-comparison-defect-preserves-paired-comparison-data} proves the comparison-clean preservation claim.  In the repair case the packet remains at the comparison node with a strictly smaller local comparison defect and no route drift; in the frontier case it returns protected comparison frontier material without promotion.  In none of the three cases is the non-certificate marker changed into an OP closure certificate.
\end{proof}

\begin{corollary}[v3.34.0-r20 working rule]
\label{cor:v3340-r20-working-rule}
Future work may proceed from the r19 route/rank packet to the paired comparison node only after \(\mathsf R_{\mathrm{ready}}\) has been returned.  The next admissible mathematical question is therefore how a guarded comparison return \(\mathsf C_{\mathrm{bind}}\), \(\mathsf C_{\mathrm{split}}\), or \(\mathsf C_{\mathrm{frontier}}\) should be resolved; it is still not whether OP~1/4 or OP~5(B) is closed.
\end{corollary}

\begin{proof}
Immediate from Theorem~\ref{thm:first-dispatched-comparison-defect-packet}.
\end{proof}

\subsection{First dispatched resolve/return packet}
\label{subsec:first-dispatched-resolve-return-packet}

\begin{definition}[Dispatched resolve/return packet]
\label{def:dispatched-resolve-return-packet}
Let \(E_5=\mathrm{RESOLVE}_{\mathrm{bind/split/frontier}}\) be the sixth node of the tail-normalized HC/OP skeleton and let
\[
  \mathsf D_{\mathrm{HCOP}}^{\mathrm{skel}}(E_5)
  =
  \mathsf W_{\mathrm{resolve}}^{\mathrm{next}}
\]
be the r15 dispatch assignment.  A \emph{dispatched resolve/return packet} is a finite packet
\[
  \mathsf P_{\mathrm{resolve}}^{\mathrm{next}}
  =
  \bigl(\mathsf C_{\eta},\eta,\Pi_{\eta},
        \rho_{\mathrm{CL}},\rho_{\Theta R},
        \Delta_{\mathrm{sep}}^{\mathrm{HC}},
        G_5,W_5,S_5,
        \mathbf d_{\mathrm{resolve}}^{\mathrm{next}}\bigr),
  \qquad
  \eta\in\{\mathrm{bind},\mathrm{split},\mathrm{frontier}\},
\]
where \(\mathsf C_{\eta}\) is the corresponding r20 comparison return, \(\Pi_{\eta}\) is the next finite continuation package associated with that return, \(\rho_{\mathrm{CL}}\) and \(\rho_{\Theta R}\) are the inherited routed residual coordinates, \(\Delta_{\mathrm{sep}}^{\mathrm{HC}}\) is the paired HC comparison defect, \(G_5\) is the resolution guard, \(W_5\) is the witness pointer to the first explicit comparison source, \(S_5\) is the non-certificate status marker, and \(\mathbf d_{\mathrm{resolve}}^{\mathrm{next}}\) is a finite resolution audit register.  The packet is admissible only if it is launched from one of \(\mathsf C_{\mathrm{bind}}\), \(\mathsf C_{\mathrm{split}}\), or \(\mathsf C_{\mathrm{frontier}}\), selects exactly the continuation package matching that return, keeps the inherited comparison data visible, and does not convert the selected continuation into an OP closure certificate.
\end{definition}

\begin{definition}[Resolve/return defect register]
\label{def:resolve-return-defect-register}
For an admissible packet \(\mathsf P_{\mathrm{resolve}}^{\mathrm{next}}\), the \emph{resolve/return defect register} is the finite family
\[
  \mathbf d_{\mathrm{resolve}}^{\mathrm{next}}(j)
  =
  \bigl(d_{\mathrm{in}}(j),d_{\mathrm{sel}}(j),d_{\mathrm{guard}}(j),
        d_{\mathrm{route}}(j),d_{\mathrm{wit}}(j),d_{\mathrm{stat}}(j)\bigr),
  \qquad j\in\mathcal A_5,
\]
where \(d_{\mathrm{in}}\) measures loss of the guarded r20 comparison return, \(d_{\mathrm{sel}}\) measures mismatch between the comparison return and the selected continuation package, \(d_{\mathrm{guard}}\) measures violation of the resolution guard, \(d_{\mathrm{route}}\) measures loss of the routed OP~1/4--OP~5(B) comparison data, \(d_{\mathrm{wit}}\) measures loss of the witness pointer, and \(d_{\mathrm{stat}}\) measures accidental promotion of the non-certificate marker.  The scalar resolve/return defect is
\[
  \Delta_{\mathrm{resolve}}^{\mathrm{next}}
  =
  \max_{j\in\mathcal A_5}
  \bigl(d_{\mathrm{in}}(j)+d_{\mathrm{sel}}(j)+d_{\mathrm{guard}}(j)+d_{\mathrm{route}}(j)+d_{\mathrm{wit}}(j)+d_{\mathrm{stat}}(j)\bigr).
\]
The register is \emph{resolve-clean} if \(\Delta_{\mathrm{resolve}}^{\mathrm{next}}=0\), \emph{resolve-repairable} if one bounded local resolution repair strictly lowers the scalar defect while preserving the guarded comparison return, and \emph{frontier-bound} otherwise.
\end{definition}

\begin{lemma}[Guarded comparison return is the resolve precondition]
\label{lem:guarded-comparison-return-is-resolve-precondition}
The resolve/return packet may be executed only from one of the r20 guarded comparison returns \(\mathsf C_{\mathrm{bind}}\), \(\mathsf C_{\mathrm{split}}\), or \(\mathsf C_{\mathrm{frontier}}\).  If the r20 comparison node has not returned one of these three guarded statuses, then the resolution node is not yet admissible.
\end{lemma}

\begin{proof}
By Definition~\ref{def:dispatched-resolve-return-packet}, the first entry of \(\mathsf P_{\mathrm{resolve}}^{\mathrm{next}}\) is exactly one guarded comparison return \(\mathsf C_{\eta}\).  The resolution node consumes a completed guarded comparison return, not a raw route/rank state and not an unresolved comparison audit surface.  Without such a return there is no admissible bind, split, or frontier branch for the packet to select.
\end{proof}

\begin{lemma}[Zero resolve/return defect preserves selected continuation data]
\label{lem:zero-resolve-return-defect-preserves-selected-continuation-data}
If \(\Delta_{\mathrm{resolve}}^{\mathrm{next}}=0\) for an admissible dispatched resolve/return packet, then the selected continuation package preserves the corresponding r20 comparison return, the routed residual data, the HC comparison defect, the resolution guard, the witness pointer, and the non-certificate status marker.
\end{lemma}

\begin{proof}
Vanishing scalar defect forces all components of Definition~\ref{def:resolve-return-defect-register} to vanish on the resolution audit window.  Hence the input comparison return is retained, the selected continuation package matches that return, the guard is not bypassed, the routed residual data and comparison defect remain attached to the selected branch, the witness pointer is preserved, and the non-certificate marker is not promoted.  Therefore the output is a guarded next work state, not an OP closure certificate.
\end{proof}

\begin{definition}[r21 resolve/return output]
\label{def:r21-resolve-return-output}
The execution of \(\mathsf W_{\mathrm{resolve}}^{\mathrm{next}}\) on an admissible packet returns exactly one of the following statuses:
\[
\mathrm{Exec}\bigl(\mathsf W_{\mathrm{resolve}}^{\mathrm{next}},
                   \mathsf P_{\mathrm{resolve}}^{\mathrm{next}}\bigr)
\in
\{\mathsf Q_{\mathrm{bind}},\mathsf Q_{\mathrm{split}},\mathsf Q_{\mathrm{frontier}}\}.
\]
Here \(\mathsf Q_{\mathrm{bind}}\) is a guarded common-origin continuation work state for the coupled OP~1/4--OP~5(B) branch, \(\mathsf Q_{\mathrm{split}}\) is a guarded ranked single-branch continuation work state, and \(\mathsf Q_{\mathrm{frontier}}\) is a protected smaller comparison-frontier descent state.  None of these statuses is an OP closure certificate.
\end{definition}

\begin{proposition}[The dispatched resolve/return packet is admissible after guarded comparison return]
\label{prop:dispatched-resolve-return-packet-admissible}
If the r20 execution returns one of \(\mathsf C_{\mathrm{bind}}\), \(\mathsf C_{\mathrm{split}}\), or \(\mathsf C_{\mathrm{frontier}}\), then the r15 dispatch menu admits the r21 resolve/return packet as the next executable work packet attached to the active HC/OP skeleton.
\end{proposition}

\begin{proof}
The r15 dispatch menu assigns \(E_5\) only to \(\mathsf W_{\mathrm{resolve}}^{\mathrm{next}}\).  Lemma~\ref{lem:guarded-comparison-return-is-resolve-precondition} supplies the required guarded comparison input status.  The remaining data of Definition~\ref{def:dispatched-resolve-return-packet} are inherited from the r04 paired comparison layer, the r07 QAM/MML/MMA interface, the r12 guarded-reduction ledger, the r14 skeleton, and the r20 comparison return.  The new defect register adds only a finite resolution audit surface.  It does not create a new certificate test, erase the residual route, suppress the witness pointer, or change the theorematic status of the selected continuation.  Hence the packet is admissible.
\end{proof}

\begin{theorem}[First dispatched resolve/return packet]
\label{thm:first-dispatched-resolve-return-packet}
After a guarded r20 comparison return, the active HC/OP skeleton supports the next executed work packet
\[
  \{\mathsf C_{\mathrm{bind}},\mathsf C_{\mathrm{split}},\mathsf C_{\mathrm{frontier}}\}
  \xrightarrow{\;\mathsf W_{\mathrm{resolve}}^{\mathrm{next}}\;}
  \{\mathsf Q_{\mathrm{bind}},\mathsf Q_{\mathrm{split}},\mathsf Q_{\mathrm{frontier}}\}.
\]
This execution is conservative: it selects a guarded common-origin continuation, a guarded ranked single-branch continuation, or a smaller comparison-frontier descent state.  It does not prove OP~1/4, OP~5(B), or any closure certificate.
\end{theorem}

\begin{proof}
By Proposition~\ref{prop:dispatched-resolve-return-packet-admissible}, the r21 packet is admissible after a guarded r20 comparison return.  Definition~\ref{def:resolve-return-defect-register} classifies the resolution audit defect into resolve-clean, resolve-repairable, or frontier-bound cases, and Definition~\ref{def:r21-resolve-return-output} records the corresponding guarded continuation output.  Lemma~\ref{lem:zero-resolve-return-defect-preserves-selected-continuation-data} proves the resolve-clean preservation claim.  In the repair case the packet remains at the resolution node with a strictly smaller local resolution defect and no route drift; in the frontier case it returns protected comparison-frontier material without promotion.  In none of the three cases is the non-certificate marker changed into an OP closure certificate.
\end{proof}

\begin{corollary}[v3.34.0-r21 working rule]
\label{cor:v3340-r21-working-rule}
Future work may proceed from the r20 guarded comparison returns to a selected next continuation only through \(\mathsf W_{\mathrm{resolve}}^{\mathrm{next}}\).  The next admissible mathematical question is therefore whether \(\mathsf Q_{\mathrm{bind}}\), \(\mathsf Q_{\mathrm{split}}\), or \(\mathsf Q_{\mathrm{frontier}}\) should seed the next finite work packet; it is still not whether OP~1/4 or OP~5(B) is closed.
\end{corollary}

\begin{proof}
Immediate from Theorem~\ref{thm:first-dispatched-resolve-return-packet}.
\end{proof}


\subsection{First continuation-queue packet after resolve/return}
\label{subsec:first-continuation-queue-packet-after-resolve-return}

\begin{definition}[Continuation-queue packet after resolve/return]
\label{def:continuation-queue-packet-after-resolve-return}
Let
\[
  \mathsf Q_{\eta}\in
  \{\mathsf Q_{\mathrm{bind}},\mathsf Q_{\mathrm{split}},\mathsf Q_{\mathrm{frontier}}\},
  \qquad
  \eta\in\{\mathrm{bind},\mathrm{split},\mathrm{frontier}\},
\]
be a guarded r21 resolve/return status.  A \emph{continuation-queue packet after resolve/return} is a finite packet
\[
  \mathsf P_{\mathrm{queue}}^{\mathrm{cont}}
  =
  \bigl(
    \mathsf Q_{\eta},
    \mathcal U_{\eta},
    \rho_{\mathrm{CL}},\rho_{\Theta R},\rho_A,
    \Delta_{\mathrm{sep}}^{\mathrm{HC}},
    G_6,W_6,S_6,
    \mathbf d_{\mathrm{queue}}^{\mathrm{cont}}
  \bigr),
\]
where \(\mathcal U_{\eta}\) is the finite guarded continuation queue selected by the r21 return, \((\rho_{\mathrm{CL}},\rho_{\Theta R},\rho_A)\) is the inherited live residual triple, \(\Delta_{\mathrm{sep}}^{\mathrm{HC}}\) is the inherited paired HC comparison defect, \(G_6\) is the queue guard, \(W_6\) is the witness pointer to the first explicit resolve/return source, \(S_6\) is the non-certificate status marker, and \(\mathbf d_{\mathrm{queue}}^{\mathrm{cont}}\) is a finite queue audit register.  The packet is admissible only if the queue type matches the selected r21 return, keeps the inherited residual and comparison data attached to the queue, and does not change a continuation queue into an OP closure certificate.
\end{definition}

\begin{definition}[Bind, split, and frontier continuation queues]
\label{def:bind-split-frontier-continuation-queues}
The continuation queue selected by Definition~\ref{def:continuation-queue-packet-after-resolve-return} is one of the following three finite guarded families:
\[
\begin{aligned}
  \mathcal U_{\mathrm{bind}}
  &=
  \bigl(
    \mathrm{commonOriginCheck},
    \mathrm{jointResidualRefine},
    \mathrm{pairedWitnessReturn}
  \bigr),\\
  \mathcal U_{\mathrm{split}}
  &=
  \bigl(
    \mathrm{branchRankCheck},
    \mathrm{singleResidualSelect},
    \mathrm{separatedWitnessReturn}
  \bigr),\\
  \mathcal U_{\mathrm{frontier}}
  &=
  \bigl(
    \mathrm{frontierDescentCheck},
    \mathrm{defectLocalize},
    \mathrm{smallerFrontierReturn}
  \bigr).
\end{aligned}
\]
The bind queue keeps OP~1/4 and OP~5(B) coupled under a common-origin guard; the split queue selects a ranked single-branch continuation under a separation guard; and the frontier queue performs only protected descent to a smaller comparison frontier.  None of these queues is a theorematic OP closure state.
\end{definition}

\begin{definition}[Continuation-queue defect register]
\label{def:continuation-queue-defect-register}
For an admissible packet \(\mathsf P_{\mathrm{queue}}^{\mathrm{cont}}\), the \emph{continuation-queue defect register} is the finite family
\[
  \mathbf d_{\mathrm{queue}}^{\mathrm{cont}}(j)
  =
  \bigl(d_{Q}(j),d_{\mathrm{type}}(j),d_{\mathrm{res}}(j),
        d_{\Delta}(j),d_{\mathrm{guard}}(j),d_{\mathrm{wit}}(j),d_{\mathrm{stat}}(j)\bigr),
  \qquad j\in\mathcal A_6,
\]
where \(d_Q\) measures loss of the r21 resolve/return status, \(d_{\mathrm{type}}\) measures mismatch between the selected status and the selected queue type, \(d_{\mathrm{res}}\) measures loss of the inherited live residual triple, \(d_{\Delta}\) measures loss of the paired comparison defect, \(d_{\mathrm{guard}}\) measures violation of the queue guard, \(d_{\mathrm{wit}}\) measures loss of the witness pointer, and \(d_{\mathrm{stat}}\) measures accidental certificate promotion.  The scalar queue defect is
\[
  \Delta_{\mathrm{queue}}^{\mathrm{cont}}
  =
  \max_{j\in\mathcal A_6}
  \bigl(d_Q(j)+d_{\mathrm{type}}(j)+d_{\mathrm{res}}(j)+d_{\Delta}(j)
        +d_{\mathrm{guard}}(j)+d_{\mathrm{wit}}(j)+d_{\mathrm{stat}}(j)\bigr).
\]
The register is \emph{queue-clean} if \(\Delta_{\mathrm{queue}}^{\mathrm{cont}}=0\), \emph{queue-repairable} if one bounded local queue repair strictly lowers the scalar defect while preserving the r21 return, and \emph{frontier-bound} otherwise.
\end{definition}

\begin{lemma}[Resolve/return output is the queue precondition]
\label{lem:resolve-return-output-is-queue-precondition}
The continuation-queue packet may be executed only from one of \(\mathsf Q_{\mathrm{bind}}\), \(\mathsf Q_{\mathrm{split}}\), or \(\mathsf Q_{\mathrm{frontier}}\).  If the r21 resolve/return node has not returned one of these three guarded statuses, then the continuation-queue node is not yet admissible.
\end{lemma}

\begin{proof}
By Definition~\ref{def:continuation-queue-packet-after-resolve-return}, the first entry of the packet is exactly one guarded r21 return \(\mathsf Q_{\eta}\).  The queue node schedules follow-on work from a completed resolve/return state; it does not consume an unresolved comparison state, a raw route/rank state, or a merely local HC-capture source.  Hence no continuation queue exists until the r21 node has selected one of the three guarded return statuses.
\end{proof}

\begin{lemma}[Zero continuation-queue defect preserves queue data]
\label{lem:zero-continuation-queue-defect-preserves-queue-data}
If \(\Delta_{\mathrm{queue}}^{\mathrm{cont}}=0\) for an admissible continuation-queue packet, then the selected queue preserves the r21 return status, the matching queue type, the inherited live residual triple, the paired HC comparison defect, the queue guard, the witness pointer, and the non-certificate status marker.
\end{lemma}

\begin{proof}
Vanishing scalar defect forces all components of Definition~\ref{def:continuation-queue-defect-register} to vanish on the queue audit window.  Thus the r21 return remains visible, the queue type matches the selected return, the live residual triple and paired comparison defect remain attached, the guard and witness pointer are preserved, and the non-certificate marker is not promoted.  Therefore the output is a finite guarded work queue, not an OP closure certificate.
\end{proof}

\begin{definition}[r22 continuation-queue output]
\label{def:r22-continuation-queue-output}
The execution of the continuation-queue packet returns exactly one of the following queued work states:
\[
\mathrm{Exec}\bigl(\mathsf W_{\mathrm{queue}}^{\mathrm{cont}},
                   \mathsf P_{\mathrm{queue}}^{\mathrm{cont}}\bigr)
\in
\{\mathsf U_{\mathrm{bind}}^{\mathrm{ready}},
  \mathsf U_{\mathrm{split}}^{\mathrm{ready}},
  \mathsf U_{\mathrm{frontier}}^{\mathrm{ready}}\}.
\]
Here \(\mathsf U_{\mathrm{bind}}^{\mathrm{ready}}\) is a finite common-origin work queue, \(\mathsf U_{\mathrm{split}}^{\mathrm{ready}}\) is a finite ranked single-branch work queue, and \(\mathsf U_{\mathrm{frontier}}^{\mathrm{ready}}\) is a finite protected frontier-descent work queue.  None of these statuses is an OP closure certificate.
\end{definition}

\begin{proposition}[The continuation-queue packet is admissible after r21 resolve/return]
\label{prop:continuation-queue-packet-admissible-after-r21-resolve-return}
If the r21 execution returns one of \(\mathsf Q_{\mathrm{bind}}\), \(\mathsf Q_{\mathrm{split}}\), or \(\mathsf Q_{\mathrm{frontier}}\), then the active HC/OP skeleton admits \(\mathsf W_{\mathrm{queue}}^{\mathrm{cont}}\) as the next executable work packet.
\end{proposition}

\begin{proof}
Lemma~\ref{lem:resolve-return-output-is-queue-precondition} supplies the required guarded r21 input.  Definitions~\ref{def:continuation-queue-packet-after-resolve-return} and~\ref{def:bind-split-frontier-continuation-queues} attach the matching finite queue to that input.  The remaining data are inherited from the r04 paired comparison layer, the r07 QAM/MML/MMA interface, the r15 dispatch, the r20 comparison return, and the r21 resolve/return output.  The queue defect register is finite and audits only preservation of input status, type, residual data, comparison defect, guard, witness, and non-certificate status.  Hence the packet is admissible as follow-on work and not as an OP closure claim.
\end{proof}

\begin{theorem}[First continuation-queue packet after resolve/return]
\label{thm:first-continuation-queue-packet-after-resolve-return}
After a guarded r21 resolve/return output, the active HC/OP skeleton supports the next executed work packet
\[
  \{\mathsf Q_{\mathrm{bind}},\mathsf Q_{\mathrm{split}},\mathsf Q_{\mathrm{frontier}}\}
  \xrightarrow{\;\mathsf W_{\mathrm{queue}}^{\mathrm{cont}}\;}
  \{\mathsf U_{\mathrm{bind}}^{\mathrm{ready}},
    \mathsf U_{\mathrm{split}}^{\mathrm{ready}},
    \mathsf U_{\mathrm{frontier}}^{\mathrm{ready}}\}.
\]
This execution is conservative: it turns the r21 return into a finite guarded continuation queue and does not prove OP~1/4, OP~5(B), or any closure certificate.
\end{theorem}

\begin{proof}
By Proposition~\ref{prop:continuation-queue-packet-admissible-after-r21-resolve-return}, the r22 packet is admissible after a guarded r21 return.  Definition~\ref{def:continuation-queue-defect-register} classifies the queue audit defect into queue-clean, queue-repairable, or frontier-bound cases, and Definition~\ref{def:r22-continuation-queue-output} records the corresponding queued work output.  Lemma~\ref{lem:zero-continuation-queue-defect-preserves-queue-data} proves the queue-clean preservation claim.  In the repair case the packet remains at the queue node with a strictly smaller local queue defect and no route drift; in the frontier case it returns protected smaller frontier material.  In none of the three cases is the non-certificate marker changed into an OP closure certificate.
\end{proof}

\begin{corollary}[v3.34.0-r22 working rule]
\label{cor:v3340-r22-working-rule}
Future work may proceed from the r21 guarded resolve/return status only by selecting the matching finite continuation queue.  The next admissible mathematical question is therefore which queued work family --- common-origin bind refinement, ranked split refinement, or protected frontier descent --- should be executed first; it is still not whether OP~1/4 or OP~5(B) is closed.
\end{corollary}

\begin{proof}
Immediate from Theorem~\ref{thm:first-continuation-queue-packet-after-resolve-return}.
\end{proof}


\subsection{First queue-head selection packet after continuation queue}
\label{subsec:first-queue-head-selection-packet-after-continuation-queue}

\begin{definition}[Queue-head selection packet after continuation queue]
\label{def:queue-head-selection-packet-after-continuation-queue}
Let
\[
  \mathsf U_{\eta}^{\mathrm{ready}}
  \in
  \{\mathsf U_{\mathrm{bind}}^{\mathrm{ready}},
    \mathsf U_{\mathrm{split}}^{\mathrm{ready}},
    \mathsf U_{\mathrm{frontier}}^{\mathrm{ready}}\},
  \qquad
  \eta\in\{\mathrm{bind},\mathrm{split},\mathrm{frontier}\},
\]
be a guarded r22 continuation-queue output.  A \emph{queue-head selection packet after continuation queue} is a finite packet
\[
  \mathsf P_{\mathrm{head}}^{\mathrm{select}}
  =
  \bigl(
    \mathsf U_{\eta}^{\mathrm{ready}},
    h_{\eta},
    \mathcal U_{\eta}^{\mathrm{tail}},
    \rho_{\mathrm{CL}},\rho_{\Theta R},\rho_A,
    \Delta_{\mathrm{sep}}^{\mathrm{HC}},
    G_7,W_7,S_7,
    \mathbf d_{\mathrm{head}}^{\mathrm{select}}
  \bigr),
\]
where \(h_{\eta}\) is the first guarded queue head selected from the matching r22 queue, \(\mathcal U_{\eta}^{\mathrm{tail}}\) is the remaining finite queue tail, \((\rho_{\mathrm{CL}},\rho_{\Theta R},\rho_A)\) is the inherited live residual triple, \(\Delta_{\mathrm{sep}}^{\mathrm{HC}}\) is the inherited paired HC comparison defect, \(G_7\) is the head-selection guard, \(W_7\) is the witness pointer to the first explicit continuation-queue source, \(S_7\) is the non-certificate status marker, and \(\mathbf d_{\mathrm{head}}^{\mathrm{select}}\) is a finite head-selection audit register.  The packet is admissible only if the selected head belongs to the matching r22 queue, the tail remains finite and guarded, the inherited residual and comparison data remain attached, and no selected head is promoted to an OP closure certificate.
\end{definition}

\begin{definition}[Continuation-queue heads and finite tails]
\label{def:continuation-queue-heads-and-finite-tails}
The matching queue head and tail in Definition~\ref{def:queue-head-selection-packet-after-continuation-queue} are fixed by the r22 queue type:
\[
\begin{aligned}
  h_{\mathrm{bind}}
  &= \mathrm{commonOriginCheck},
  &\mathcal U_{\mathrm{bind}}^{\mathrm{tail}}
  &= (\mathrm{jointResidualRefine},\mathrm{pairedWitnessReturn}),\\
  h_{\mathrm{split}}
  &= \mathrm{branchRankCheck},
  &\mathcal U_{\mathrm{split}}^{\mathrm{tail}}
  &= (\mathrm{singleResidualSelect},\mathrm{separatedWitnessReturn}),\\
  h_{\mathrm{frontier}}
  &= \mathrm{frontierDescentCheck},
  &\mathcal U_{\mathrm{frontier}}^{\mathrm{tail}}
  &= (\mathrm{defectLocalize},\mathrm{smallerFrontierReturn}).
\end{aligned}
\]
Thus the bind queue first tests whether the common-origin reading is still admissible, the split queue first tests whether a ranked branch can honestly be selected, and the frontier queue first tests whether a protected descent is available.  The selected head is a next work packet, not a theorematic OP closure state.
\end{definition}

\begin{definition}[Queue-head selection defect register]
\label{def:queue-head-selection-defect-register}
For an admissible packet \(\mathsf P_{\mathrm{head}}^{\mathrm{select}}\), the \emph{queue-head selection defect register} is the finite family
\[
  \mathbf d_{\mathrm{head}}^{\mathrm{select}}(j)
  =
  \bigl(d_U(j),d_h(j),d_{\mathrm{tail}}(j),d_{\mathrm{res}}(j),
        d_{\Delta}(j),d_{\mathrm{guard}}(j),d_{\mathrm{wit}}(j),d_{\mathrm{stat}}(j)\bigr),
  \qquad j\in\mathcal A_7,
\]
where \(d_U\) measures loss of the r22 queued work state, \(d_h\) measures mismatch between the queue type and selected head, \(d_{\mathrm{tail}}\) measures loss of the finite guarded queue tail, \(d_{\mathrm{res}}\) measures loss of the live residual triple, \(d_{\Delta}\) measures loss of the paired HC comparison defect, \(d_{\mathrm{guard}}\) measures guard drift, \(d_{\mathrm{wit}}\) measures witness-pointer drift, and \(d_{\mathrm{stat}}\) measures accidental certificate promotion.  The scalar head-selection defect is
\[
  \Delta_{\mathrm{head}}^{\mathrm{select}}
  =
  \max_{j\in\mathcal A_7}
  \bigl(d_U(j)+d_h(j)+d_{\mathrm{tail}}(j)+d_{\mathrm{res}}(j)+d_{\Delta}(j)
        +d_{\mathrm{guard}}(j)+d_{\mathrm{wit}}(j)+d_{\mathrm{stat}}(j)\bigr).
\]
The register is \emph{head-clean} if \(\Delta_{\mathrm{head}}^{\mathrm{select}}=0\), \emph{head-repairable} if one bounded local head-selection repair strictly lowers the scalar defect while preserving the r22 queue, and \emph{frontier-bound} otherwise.
\end{definition}

\begin{lemma}[Continuation-queue output is the head-selection precondition]
\label{lem:continuation-queue-output-is-head-selection-precondition}
The queue-head selection packet may be executed only from one of \(\mathsf U_{\mathrm{bind}}^{\mathrm{ready}}\), \(\mathsf U_{\mathrm{split}}^{\mathrm{ready}}\), or \(\mathsf U_{\mathrm{frontier}}^{\mathrm{ready}}\).  If the r22 continuation-queue node has not returned one of these three guarded ready queues, then the head-selection node is not yet admissible.
\end{lemma}

\begin{proof}
By Definition~\ref{def:queue-head-selection-packet-after-continuation-queue}, the first entry of the packet is exactly one r22 queued work state \(\mathsf U_{\eta}^{\mathrm{ready}}\).  Head selection consumes a completed finite queue and selects its first guarded work item; it does not consume a raw r21 return, an unresolved comparison return, or a merely local HC-capture source.  Hence no head-selection packet exists until the r22 node has returned a matching ready queue.
\end{proof}

\begin{lemma}[Zero queue-head selection defect preserves selected-head data]
\label{lem:zero-queue-head-selection-defect-preserves-selected-head-data}
If \(\Delta_{\mathrm{head}}^{\mathrm{select}}=0\) for an admissible queue-head selection packet, then the selected head preserves the r22 queued work state, the matching head type, the finite guarded tail, the inherited live residual triple, the paired HC comparison defect, the head-selection guard, the witness pointer, and the non-certificate status marker.
\end{lemma}

\begin{proof}
Vanishing scalar defect forces all components of Definition~\ref{def:queue-head-selection-defect-register} to vanish on the head-selection audit window.  Thus the r22 ready queue remains visible, the selected head matches that queue, the finite tail is retained for later execution, the live residual triple and paired comparison defect remain attached, the guard and witness pointer are preserved, and the non-certificate marker is not promoted.  Therefore the output is a guarded first work head together with a finite continuation tail, not an OP closure certificate.
\end{proof}

\begin{definition}[r23 queue-head selection output]
\label{def:r23-queue-head-selection-output}
The execution of the queue-head selection packet returns exactly one of the following selected-head work states:
\[
\mathrm{Exec}\bigl(\mathsf W_{\mathrm{head}}^{\mathrm{select}},
                   \mathsf P_{\mathrm{head}}^{\mathrm{select}}\bigr)
\in
\{\mathsf X_{\mathrm{bind}}^{\mathrm{head}},
  \mathsf X_{\mathrm{split}}^{\mathrm{head}},
  \mathsf X_{\mathrm{frontier}}^{\mathrm{head}}\}.
\]
Here \(\mathsf X_{\mathrm{bind}}^{\mathrm{head}}\) is the guarded common-origin-check head with its finite bind tail, \(\mathsf X_{\mathrm{split}}^{\mathrm{head}}\) is the guarded branch-rank-check head with its finite split tail, and \(\mathsf X_{\mathrm{frontier}}^{\mathrm{head}}\) is the guarded frontier-descent-check head with its finite frontier tail.  None of these statuses is an OP closure certificate.
\end{definition}

\begin{proposition}[The queue-head selection packet is admissible after r22 queue return]
\label{prop:queue-head-selection-packet-admissible-after-r22-queue-return}
If the r22 execution returns one of \(\mathsf U_{\mathrm{bind}}^{\mathrm{ready}}\), \(\mathsf U_{\mathrm{split}}^{\mathrm{ready}}\), or \(\mathsf U_{\mathrm{frontier}}^{\mathrm{ready}}\), then the active HC/OP skeleton admits \(\mathsf W_{\mathrm{head}}^{\mathrm{select}}\) as the next executable work packet.
\end{proposition}

\begin{proof}
Lemma~\ref{lem:continuation-queue-output-is-head-selection-precondition} supplies the required guarded r22 input.  Definition~\ref{def:continuation-queue-heads-and-finite-tails} supplies the matching head and finite tail for that input.  The remaining data are inherited from the r04 paired comparison layer, the r07 QAM/MML/MMA interface, the r15 dispatch, the r20 comparison return, the r21 resolve/return output, and the r22 continuation queue.  The head-selection defect register is finite and audits only preservation of input queue, head type, tail, residual data, comparison defect, guard, witness, and non-certificate status.  Hence the packet is admissible as follow-on work and not as an OP closure claim.
\end{proof}

\begin{theorem}[First queue-head selection packet after continuation queue]
\label{thm:first-queue-head-selection-packet-after-continuation-queue}
After a guarded r22 continuation-queue output, the active HC/OP skeleton supports the next executed work packet
\[
  \{\mathsf U_{\mathrm{bind}}^{\mathrm{ready}},
    \mathsf U_{\mathrm{split}}^{\mathrm{ready}},
    \mathsf U_{\mathrm{frontier}}^{\mathrm{ready}}\}
  \xrightarrow{\;\mathsf W_{\mathrm{head}}^{\mathrm{select}}\;}
  \{\mathsf X_{\mathrm{bind}}^{\mathrm{head}},
    \mathsf X_{\mathrm{split}}^{\mathrm{head}},
    \mathsf X_{\mathrm{frontier}}^{\mathrm{head}}\}.
\]
This execution is conservative: it selects the first guarded head of the matching finite continuation queue and keeps the remaining finite tail available for later work.  It does not prove OP~1/4, OP~5(B), or any closure certificate.
\end{theorem}

\begin{proof}
By Proposition~\ref{prop:queue-head-selection-packet-admissible-after-r22-queue-return}, the r23 packet is admissible after a guarded r22 ready queue.  Definition~\ref{def:queue-head-selection-defect-register} classifies the head-selection audit defect into head-clean, head-repairable, or frontier-bound cases, and Definition~\ref{def:r23-queue-head-selection-output} records the corresponding selected-head work output.  Lemma~\ref{lem:zero-queue-head-selection-defect-preserves-selected-head-data} proves the head-clean preservation claim.  In the repair case the packet remains at the head-selection node with a strictly smaller local head-selection defect and no route drift; in the frontier case it returns protected smaller frontier material.  In none of the three cases is the non-certificate marker changed into an OP closure certificate.
\end{proof}

\begin{corollary}[v3.34.0-r23 working rule]
\label{cor:v3340-r23-working-rule}
Future work may proceed from an r22 guarded ready queue only by selecting the matching first queue head while retaining the finite guarded tail.  The next admissible mathematical question is therefore whether to execute \(\mathrm{commonOriginCheck}\), \(\mathrm{branchRankCheck}\), or \(\mathrm{frontierDescentCheck}\) under its inherited guard; it is still not whether OP~1/4 or OP~5(B) is closed.
\end{corollary}

\begin{proof}
Immediate from Theorem~\ref{thm:first-queue-head-selection-packet-after-continuation-queue}.
\end{proof}


\subsection{First selected queue-head execution packet}
\label{subsec:first-selected-queue-head-execution-packet}

\begin{definition}[Selected queue-head execution packet]
\label{def:selected-queue-head-execution-packet}
Let
\[
  \mathsf X_{\eta}^{\mathrm{head}}
  \in
  \{\mathsf X_{\mathrm{bind}}^{\mathrm{head}},
    \mathsf X_{\mathrm{split}}^{\mathrm{head}},
    \mathsf X_{\mathrm{frontier}}^{\mathrm{head}}\},
  \qquad
  \eta\in\{\mathrm{bind},\mathrm{split},\mathrm{frontier}\},
\]
be a guarded r23 selected queue-head output.  A \emph{selected queue-head execution packet} is a finite packet
\[
  \mathsf P_{\mathrm{head}}^{\mathrm{exec}}
  =
  \bigl(
    \mathsf X_{\eta}^{\mathrm{head}},
    h_{\eta},
    \mathcal U_{\eta}^{\mathrm{tail}},
    \rho_{\mathrm{CL}},\rho_{\Theta R},\rho_A,
    \Delta_{\mathrm{sep}}^{\mathrm{HC}},
    \Delta_{\mathrm{cmp}}^{\Delta},
    G_8,W_8,S_8,
    \mathbf d_{\mathrm{head}}^{\mathrm{exec}}
  \bigr),
\]
where \(h_{\eta}\) is the selected r23 queue head, \(\mathcal U_{\eta}^{\mathrm{tail}}\) is the remaining finite guarded tail, \((\rho_{\mathrm{CL}},\rho_{\Theta R},\rho_A)\) is the inherited live residual triple, \(\Delta_{\mathrm{sep}}^{\mathrm{HC}}\) is the inherited paired HC separation defect, \(\Delta_{\mathrm{cmp}}^{\Delta}\) is the r20 comparison-defect witness, \(G_8\) is the execution guard, \(W_8\) is the witness pointer to the selected-head source, \(S_8\) is the non-certificate status marker, and \(\mathbf d_{\mathrm{head}}^{\mathrm{exec}}\) is a finite head-execution audit register.  The packet is admissible only if it executes the selected head belonging to the matching r23 output, preserves the finite tail for later continuation, keeps the paired OP~1/4--OP~5(B) data attached, and does not promote the executed head to an OP closure certificate.
\end{definition}

\begin{definition}[Executed queue-head tests]
\label{def:executed-queue-head-tests}
The selected queue-head execution in Definition~\ref{def:selected-queue-head-execution-packet} is fixed by the r23 head type:
\[
\begin{aligned}
  \mathsf X_{\mathrm{bind}}^{\mathrm{head}}
  &\leadsto
  \mathrm{ExecHead}(\mathrm{commonOriginCheck}),\\
  \mathsf X_{\mathrm{split}}^{\mathrm{head}}
  &\leadsto
  \mathrm{ExecHead}(\mathrm{branchRankCheck}),\\
  \mathsf X_{\mathrm{frontier}}^{\mathrm{head}}
  &\leadsto
  \mathrm{ExecHead}(\mathrm{frontierDescentCheck}).
\end{aligned}
\]
The bind-head execution checks whether the common-origin reading of \(\rho_{\mathrm{CL}}\oplus\rho_{\Theta R}\) remains admissible under the inherited HC comparison defect.  The split-head execution checks whether a ranked single-branch continuation can be selected without erasing the paired comparison witness.  The frontier-head execution checks whether a protected descent to a smaller comparison frontier is available.  Each head test is a work test, not a theorematic closure test.
\end{definition}

\begin{definition}[Selected queue-head execution defect register]
\label{def:selected-queue-head-execution-defect-register}
For an admissible packet \(\mathsf P_{\mathrm{head}}^{\mathrm{exec}}\), the \emph{selected queue-head execution defect register} is the finite family
\[
  \mathbf d_{\mathrm{head}}^{\mathrm{exec}}(j)
  =
  \bigl(d_X(j),d_h(j),d_{\mathrm{tail}}(j),d_{\mathrm{res}}(j),
        d_{\Delta}(j),d_{\mathrm{cmp}}(j),d_{\mathrm{guard}}(j),d_{\mathrm{wit}}(j),d_{\mathrm{stat}}(j)\bigr),
  \qquad j\in\mathcal A_8,
\]
where \(d_X\) measures loss of the r23 selected-head state, \(d_h\) measures mismatch between the selected head and the executed head test, \(d_{\mathrm{tail}}\) measures loss of the finite queue tail, \(d_{\mathrm{res}}\) measures loss of the live residual triple, \(d_{\Delta}\) measures loss of the inherited HC separation defect, \(d_{\mathrm{cmp}}\) measures loss of the r20 comparison-defect witness, \(d_{\mathrm{guard}}\) measures execution-guard drift, \(d_{\mathrm{wit}}\) measures witness-pointer drift, and \(d_{\mathrm{stat}}\) measures accidental certificate promotion.  The scalar selected-head execution defect is
\[
  \Delta_{\mathrm{head}}^{\mathrm{exec}}
  =
  \max_{j\in\mathcal A_8}
  \bigl(d_X(j)+d_h(j)+d_{\mathrm{tail}}(j)+d_{\mathrm{res}}(j)+d_{\Delta}(j)
        +d_{\mathrm{cmp}}(j)+d_{\mathrm{guard}}(j)+d_{\mathrm{wit}}(j)+d_{\mathrm{stat}}(j)\bigr).
\]
The register is \emph{head-exec-clean} if \(\Delta_{\mathrm{head}}^{\mathrm{exec}}=0\), \emph{head-exec-repairable} if one bounded local execution repair strictly lowers the scalar defect while preserving the selected head and finite tail, and \emph{frontier-bound} otherwise.
\end{definition}

\begin{lemma}[Selected head is the head-execution precondition]
\label{lem:selected-head-is-head-execution-precondition}
The selected queue-head execution packet may be executed only from one of \(\mathsf X_{\mathrm{bind}}^{\mathrm{head}}\), \(\mathsf X_{\mathrm{split}}^{\mathrm{head}}\), or \(\mathsf X_{\mathrm{frontier}}^{\mathrm{head}}\).  If the r23 head-selection node has not returned one of these three guarded head states, then the head-execution node is not yet admissible.
\end{lemma}

\begin{proof}
By Definition~\ref{def:selected-queue-head-execution-packet}, the first entry of the packet is exactly one selected r23 head state.  The execution node consumes a selected head, not a raw r22 queue, a raw r21 resolve return, or an earlier comparison candidate.  Hence the head-execution packet is undefined until the r23 node has selected the matching guarded queue head and retained the finite tail.
\end{proof}

\begin{lemma}[Zero selected-head execution defect preserves head-execution data]
\label{lem:zero-selected-head-execution-defect-preserves-head-execution-data}
If \(\Delta_{\mathrm{head}}^{\mathrm{exec}}=0\) for an admissible selected queue-head execution packet, then the executed head preserves the r23 selected-head state, the matching executed test, the finite guarded tail, the inherited live residual triple, the paired HC separation defect, the r20 comparison-defect witness, the execution guard, the witness pointer, and the non-certificate status marker.
\end{lemma}

\begin{proof}
Vanishing scalar defect forces all components of Definition~\ref{def:selected-queue-head-execution-defect-register} to vanish on the execution audit window.  Thus the selected head is retained, the executed test matches the selected head type, the finite tail remains available for follow-on work, the residual and comparison data remain attached, the guard and witness pointer do not drift, and the non-certificate marker is not promoted.  Therefore the output is an executed first queue head with a finite continuation tail, not an OP closure certificate.
\end{proof}

\begin{definition}[r24 selected queue-head execution output]
\label{def:r24-selected-queue-head-execution-output}
The execution of the selected queue-head packet returns exactly one of the following executed-head work states:
\[
\mathrm{Exec}\bigl(\mathsf W_{\mathrm{head}}^{\mathrm{exec}},
                   \mathsf P_{\mathrm{head}}^{\mathrm{exec}}\bigr)
\in
\{\mathsf Y_{\mathrm{bind}}^{\mathrm{exec}},
  \mathsf Y_{\mathrm{split}}^{\mathrm{exec}},
  \mathsf Y_{\mathrm{frontier}}^{\mathrm{exec}}\}.
\]
Here \(\mathsf Y_{\mathrm{bind}}^{\mathrm{exec}}\) is the executed common-origin-check surface with its finite bind tail, \(\mathsf Y_{\mathrm{split}}^{\mathrm{exec}}\) is the executed branch-rank-check surface with its finite split tail, and \(\mathsf Y_{\mathrm{frontier}}^{\mathrm{exec}}\) is the executed frontier-descent-check surface with its finite frontier tail.  None of these statuses is an OP closure certificate.
\end{definition}

\begin{proposition}[Selected queue-head execution is admissible after r23 head selection]
\label{prop:selected-queue-head-execution-admissible-after-r23-head-selection}
If the r23 execution returns one of \(\mathsf X_{\mathrm{bind}}^{\mathrm{head}}\), \(\mathsf X_{\mathrm{split}}^{\mathrm{head}}\), or \(\mathsf X_{\mathrm{frontier}}^{\mathrm{head}}\), then the active HC/OP skeleton admits \(\mathsf W_{\mathrm{head}}^{\mathrm{exec}}\) as the next executable work packet.
\end{proposition}

\begin{proof}
Lemma~\ref{lem:selected-head-is-head-execution-precondition} supplies the required guarded r23 input.  Definition~\ref{def:executed-queue-head-tests} supplies the matching head test.  The residual data, paired HC separation defect, comparison-defect witness, guard, witness pointer, finite tail, and non-certificate status are inherited from the r20 comparison packet, the r21 resolve return, the r22 continuation queue, and the r23 selected head.  The execution defect register is finite and audits only preservation of those inherited data.  Hence the packet is admissible as follow-on work and not as an OP closure claim.
\end{proof}

\begin{theorem}[First selected queue-head execution packet]
\label{thm:first-selected-queue-head-execution-packet}
After a guarded r23 queue-head selection output, the active HC/OP skeleton supports the next executed work packet
\[
  \{\mathsf X_{\mathrm{bind}}^{\mathrm{head}},
    \mathsf X_{\mathrm{split}}^{\mathrm{head}},
    \mathsf X_{\mathrm{frontier}}^{\mathrm{head}}\}
  \xrightarrow{\;\mathsf W_{\mathrm{head}}^{\mathrm{exec}}\;}
  \{\mathsf Y_{\mathrm{bind}}^{\mathrm{exec}},
    \mathsf Y_{\mathrm{split}}^{\mathrm{exec}},
    \mathsf Y_{\mathrm{frontier}}^{\mathrm{exec}}\}.
\]
This execution is conservative: it runs the first guarded head of the matching continuation queue, keeps the finite tail available for later work, and does not prove OP~1/4, OP~5(B), or any closure certificate.
\end{theorem}

\begin{proof}
By Proposition~\ref{prop:selected-queue-head-execution-admissible-after-r23-head-selection}, the r24 packet is admissible after a guarded r23 selected head.  Definition~\ref{def:selected-queue-head-execution-defect-register} classifies the head-execution audit defect into head-exec-clean, head-exec-repairable, or frontier-bound cases, and Definition~\ref{def:r24-selected-queue-head-execution-output} records the corresponding executed-head work output.  Lemma~\ref{lem:zero-selected-head-execution-defect-preserves-head-execution-data} proves the clean preservation claim.  In the repair case the packet remains at the selected-head execution node with a strictly smaller local execution defect and no route drift; in the frontier case it returns protected smaller frontier material.  In none of the three cases is the non-certificate marker changed into an OP closure certificate.
\end{proof}

\begin{corollary}[v3.34.0-r24 working rule]
\label{cor:v3340-r24-working-rule}
Future work may proceed from an r23 selected queue head only by executing the matching guarded head test while retaining the finite queue tail.  The next admissible mathematical question is therefore whether the executed bind, split, or frontier head can legally enter its corresponding tail task --- joint residual refinement, single residual selection, or defect localization --- still under the inherited no-false-certificate rule.
\end{corollary}

\begin{proof}
Immediate from Theorem~\ref{thm:first-selected-queue-head-execution-packet}.
\end{proof}
\subsection{First executed-head tail-entry packet}
\label{subsec:first-executed-head-tail-entry-packet}

\begin{definition}[Executed-head tail-entry packet]
\label{def:executed-head-tail-entry-packet}
Let
\[
  \mathsf Y_{\eta}^{\mathrm{exec}}
  \in
  \{\mathsf Y_{\mathrm{bind}}^{\mathrm{exec}},
    \mathsf Y_{\mathrm{split}}^{\mathrm{exec}},
    \mathsf Y_{\mathrm{frontier}}^{\mathrm{exec}}\},
  \qquad
  \eta\in\{\mathrm{bind},\mathrm{split},\mathrm{frontier}\},
\]
be a guarded executed-head output of Definition~\ref{def:r24-selected-queue-head-execution-output}.  An \emph{executed-head tail-entry packet} is a finite packet
\[
  \mathsf P_{\mathrm{tail}}^{\mathrm{entry}}
  =
  \bigl(
    \mathsf Y_{\eta}^{\mathrm{exec}},
    \mathcal U_{\eta}^{\mathrm{tail}},
    \rho_{\mathrm{CL}},\rho_{\Theta R},\rho_A,
    \Delta_{\mathrm{sep}}^{\mathrm{HC}},
    \Delta_{\mathrm{cmp}}^{\Delta},
    \Delta_{\mathrm{head}}^{\mathrm{exec}},
    G_9,W_9,S_9,
    \mathbf d_{\mathrm{tail}}^{\mathrm{entry}}
  \bigr),
\]
where \(\mathcal U_{\eta}^{\mathrm{tail}}\) is the finite guarded tail retained by the selected-head execution, \((\rho_{\mathrm{CL}},\rho_{\Theta R},\rho_A)\) is the inherited live residual triple, \(\Delta_{\mathrm{sep}}^{\mathrm{HC}}\) and \(\Delta_{\mathrm{cmp}}^{\Delta}\) are the inherited comparison witnesses, \(\Delta_{\mathrm{head}}^{\mathrm{exec}}\) is the selected-head execution defect, \(G_9\) is the tail-entry guard, \(W_9\) is the witness pointer to the executed-head source, \(S_9\) is the non-certificate status marker, and \(\mathbf d_{\mathrm{tail}}^{\mathrm{entry}}\) is a finite tail-entry audit register.  The packet is admissible only if it enters the tail matching the executed head, preserves the residual and comparison data, preserves the finite tail, and does not promote the tail entry to an OP closure certificate.
\end{definition}

\begin{definition}[Executed-head tail-entry defect]
\label{def:executed-head-tail-entry-defect}
For an admissible packet \(\mathsf P_{\mathrm{tail}}^{\mathrm{entry}}\), the \emph{tail-entry defect register} is the finite family
\[
  \mathbf d_{\mathrm{tail}}^{\mathrm{entry}}(j)
  =
  \bigl(d_Y(j),d_{\mathrm{tail}}(j),d_{\mathrm{res}}(j),d_{\Delta}(j),
        d_{\mathrm{head}}(j),d_{\mathrm{guard}}(j),d_{\mathrm{wit}}(j),d_{\mathrm{stat}}(j)\bigr),
  \qquad j\in\mathcal A_{\mathrm{tail}},
\]
where the components measure loss of the executed-head state, the finite tail, the residual triple, the comparison witnesses, the head-execution defect, the tail-entry guard, the witness pointer, and the non-certificate marker.  The scalar tail-entry defect is
\[
  \Delta_{\mathrm{tail}}^{\mathrm{entry}}
  =
  \max_{j\in\mathcal A_{\mathrm{tail}}}
  \bigl(d_Y(j)+d_{\mathrm{tail}}(j)+d_{\mathrm{res}}(j)+d_{\Delta}(j)+d_{\mathrm{head}}(j)
        +d_{\mathrm{guard}}(j)+d_{\mathrm{wit}}(j)+d_{\mathrm{stat}}(j)\bigr).
\]
The register is \emph{tail-entry-clean} if \(\Delta_{\mathrm{tail}}^{\mathrm{entry}}=0\), \emph{tail-entry-repairable} if one bounded local tail-entry repair strictly lowers the scalar defect while preserving the executed head and finite tail, and \emph{frontier-bound} otherwise.
\end{definition}

\begin{lemma}[Executed head is the tail-entry precondition]
\label{lem:executed-head-is-tail-entry-precondition}
The tail-entry packet may be launched only from one of \(\mathsf Y_{\mathrm{bind}}^{\mathrm{exec}}\), \(\mathsf Y_{\mathrm{split}}^{\mathrm{exec}}\), or \(\mathsf Y_{\mathrm{frontier}}^{\mathrm{exec}}\).  If the selected-head execution has not returned one of these three guarded executed-head states, then tail entry is not yet admissible.
\end{lemma}

\begin{proof}
Definition~\ref{def:executed-head-tail-entry-packet} takes the guarded executed-head state as the first datum of the packet.  A selected but unexecuted head, a ready queue, or a resolve return does not yet determine which finite tail is allowed to open.  Hence the tail-entry node is undefined until the selected-head execution has returned the matching executed-head state and retained its finite guarded tail.
\end{proof}

\begin{lemma}[Zero tail-entry defect preserves tail-entry data]
\label{lem:zero-tail-entry-defect-preserves-tail-entry-data}
If \(\Delta_{\mathrm{tail}}^{\mathrm{entry}}=0\) for an admissible executed-head tail-entry packet, then the packet preserves the executed-head state, the matching finite tail, the live residual triple, the inherited comparison witnesses, the selected-head execution defect, the tail-entry guard, the witness pointer, and the non-certificate status marker.
\end{lemma}

\begin{proof}
Vanishing scalar defect forces every component of Definition~\ref{def:executed-head-tail-entry-defect} to vanish on the tail-entry audit window.  Therefore no executed-head data, residual data, comparison witness, guard, witness pointer, or status marker is lost.  The output is a guarded tail-entry surface, not an OP closure certificate.
\end{proof}

\begin{definition}[Executed-head tail-entry output]
\label{def:executed-head-tail-entry-output}
The execution of the tail-entry packet returns exactly one of the following guarded tail-entry surfaces:
\[
\mathrm{Exec}\bigl(\mathsf W_{\mathrm{tail}}^{\mathrm{entry}},
                   \mathsf P_{\mathrm{tail}}^{\mathrm{entry}}\bigr)
\in
\{\mathsf Z_{\mathrm{bind}}^{\mathrm{tail}},
  \mathsf Z_{\mathrm{split}}^{\mathrm{tail}},
  \mathsf Z_{\mathrm{frontier}}^{\mathrm{tail}}\}.
\]
Here \(\mathsf Z_{\mathrm{bind}}^{\mathrm{tail}}\) enters the guarded joint-residual refinement tail, \(\mathsf Z_{\mathrm{split}}^{\mathrm{tail}}\) enters the guarded single-branch selection tail, and \(\mathsf Z_{\mathrm{frontier}}^{\mathrm{tail}}\) enters the guarded defect-localization tail.  None of these statuses is an OP closure certificate.
\end{definition}

\begin{proposition}[Tail entry is admissible after selected-head execution]
\label{prop:tail-entry-admissible-after-selected-head-execution}
If the selected-head execution returns one of \(\mathsf Y_{\mathrm{bind}}^{\mathrm{exec}}\), \(\mathsf Y_{\mathrm{split}}^{\mathrm{exec}}\), or \(\mathsf Y_{\mathrm{frontier}}^{\mathrm{exec}}\), then the active HC/OP skeleton admits \(\mathsf W_{\mathrm{tail}}^{\mathrm{entry}}\) as the next executable work packet.
\end{proposition}

\begin{proof}
Lemma~\ref{lem:executed-head-is-tail-entry-precondition} supplies the required executed-head input.  The matching finite tail, residual triple, comparison witnesses, head-execution defect, guard, witness pointer, and non-certificate marker are inherited from the selected-head execution packet.  The new defect register audits only preservation of those inherited data.  Thus the packet is admissible as follow-on work and not as an OP closure claim.
\end{proof}

\begin{theorem}[First executed-head tail-entry packet]
\label{thm:first-executed-head-tail-entry-packet}
After a guarded selected-head execution output, the active HC/OP skeleton supports the next tail-entry work packet
\[
  \{\mathsf Y_{\mathrm{bind}}^{\mathrm{exec}},
    \mathsf Y_{\mathrm{split}}^{\mathrm{exec}},
    \mathsf Y_{\mathrm{frontier}}^{\mathrm{exec}}\}
  \xrightarrow{\;\mathsf W_{\mathrm{tail}}^{\mathrm{entry}}\;}
  \{\mathsf Z_{\mathrm{bind}}^{\mathrm{tail}},
    \mathsf Z_{\mathrm{split}}^{\mathrm{tail}},
    \mathsf Z_{\mathrm{frontier}}^{\mathrm{tail}}\}.
\]
This execution is conservative: it opens the tail belonging to the executed head, retains all residual and comparison data, and does not prove OP~1/4, OP~5(B), or any closure certificate.
\end{theorem}

\begin{proof}
By Proposition~\ref{prop:tail-entry-admissible-after-selected-head-execution}, the packet is admissible after a guarded executed-head state.  Definition~\ref{def:executed-head-tail-entry-defect} classifies the tail-entry defect into clean, repairable, or frontier-bound cases, while Definition~\ref{def:executed-head-tail-entry-output} records the corresponding guarded tail-entry surface.  Lemma~\ref{lem:zero-tail-entry-defect-preserves-tail-entry-data} proves the clean preservation claim.  In the repair case the packet remains at the tail-entry node with a strictly smaller local tail-entry defect; in the frontier case it returns protected smaller frontier material.  In none of these cases is the non-certificate marker changed into an OP closure certificate.
\end{proof}

\begin{corollary}[Tail-entry working rule]
\label{cor:tail-entry-working-rule}
Future work may proceed from an executed queue head only by entering the matching finite guarded tail.  The next admissible mathematical question is therefore whether the bind, split, or frontier tail can be refined under its inherited guard; it is still not whether OP~1/4 or OP~5(B) is closed.
\end{corollary}

\begin{proof}
Immediate from Theorem~\ref{thm:first-executed-head-tail-entry-packet}.
\end{proof}

\subsection{First audit-freeze ledger}
\label{subsec:first-audit-freeze-ledger}

\begin{definition}[Audit-freeze ledger]
\label{def:first-audit-freeze-ledger}
The \emph{audit-freeze ledger} for the current HC/OP line is the finite record
\[
  F_{\mathrm{HCOP}}^{\mathrm{audit}}
  =
  \bigl(
    \mathcal B_{\mathrm{public}},
    \mathcal D_{\mathrm{doi}},
    \mathcal S_{\mathrm{spine}},
    \mathcal Z_{\mathrm{tail}},
    \mathcal O_{\mathrm{pres}},
    \mathcal C_{\mathrm{ref}},
    \mathcal B_{\mathrm{finalbib}},
    \mathcal H_{\mathrm{witness}},
    S_{\mathrm{noCert}}
  \bigr).
\]
Here \(\mathcal B_{\mathrm{public}}\) records the public baseline and concept-DOI status, \(\mathcal D_{\mathrm{doi}}\) records the version-specific Zenodo DOI assigned in the public release metadata, \(\mathcal S_{\mathrm{spine}}\) records the local HC/OP spine up to the tail-entry packet, \(\mathcal Z_{\mathrm{tail}}\) records the guarded tail-entry output, \(\mathcal O_{\mathrm{pres}}\) records preservation obligations, \(\mathcal C_{\mathrm{ref}}\) records compilation/reference checks as audit data, \(\mathcal B_{\mathrm{finalbib}}\) records terminal bibliography status, \(\mathcal H_{\mathrm{witness}}\) records preservation of the historical witness, and \(S_{\mathrm{noCert}}\) records that no OP closure certificate has been manufactured.
\end{definition}

\begin{definition}[Audit-freeze defect]
\label{def:audit-freeze-defect}
For a ledger \(F_{\mathrm{HCOP}}^{\mathrm{audit}}\), the \emph{audit-freeze defect} is
\[
  \Delta_{\mathrm{audit}}^{\mathrm{freeze}}
  =
  \max
  \{d_{\mathrm{base}},d_{\mathrm{doi}},d_{\mathrm{spine}},d_{\mathrm{tail}},
    d_{\mathrm{pres}},d_{\mathrm{ref}},d_{\mathrm{bib}},d_{\mathrm{hist}},d_{\mathrm{cert}}\}.
\]
The entries measure drift of the public baseline, DOI metadata, local spine, tail-entry output, preservation obligations, reference/compilation audit data, terminal bibliography, historical witness, and non-certificate status.  The ledger is \emph{audit-clean} if \(\Delta_{\mathrm{audit}}^{\mathrm{freeze}}=0\), \emph{audit-repairable} if only mechanical or metadata repair is needed, and \emph{not freeze-ready} otherwise.
\end{definition}

\begin{lemma}[Audit ledger is metadata-bearing, not theorem-generating]
\label{lem:audit-ledger-is-metadata-bearing-not-theorem-generating}
The audit-freeze ledger records status, preservation, and release-readiness data.  It cannot by itself create an OP closure certificate or change the theorematic status of OP~1/4 or OP~5(B).
\end{lemma}

\begin{proof}
Every component of Definition~\ref{def:first-audit-freeze-ledger} is either inherited mathematical state, preservation metadata, compilation/reference audit data, bibliography status, historical-witness status, or the explicit non-certificate marker.  None of these components is a certificate test for OP~1/4 or OP~5(B).  Therefore the ledger can freeze and audit the current line, but it cannot promote the current line into a theorematic closure.
\end{proof}

\begin{theorem}[First audit-freeze ledger]
\label{thm:first-audit-freeze-ledger}
After the executed-head tail-entry packet, the current HC/OP line admits the finite audit-freeze ledger \(F_{\mathrm{HCOP}}^{\mathrm{audit}}\) with audit-freeze defect \(\Delta_{\mathrm{audit}}^{\mathrm{freeze}}\).  The ledger records the public baseline, DOI state, local spine, tail-entry output, preservation obligations, reference/compilation audit data, terminal bibliography status, historical witness, and no-false-certificate marker without changing theorematic status.
\end{theorem}

\begin{proof}
The public baseline and DOI state are metadata inherited from the title and chronology records.  The local spine and tail-entry output are supplied by Theorem~\ref{thm:first-executed-head-tail-entry-packet} and its predecessors.  The preservation obligations, reference/compilation checks, terminal bibliography status, historical witness, and non-certificate marker are already explicit release-audit obligations of the working line.  Definition~\ref{def:first-audit-freeze-ledger} collects these data into one finite record, and Lemma~\ref{lem:audit-ledger-is-metadata-bearing-not-theorem-generating} prevents any promotion to OP closure.
\end{proof}

\subsection{Audit-freeze stabilization object}
\label{subsec:audit-freeze-stabilization-object}

\begin{definition}[Audit-freeze stabilization object]
\label{def:audit-freeze-stabilization-object}
The \emph{audit-freeze stabilization object} is
\[
  F_{\mathrm{freeze}}^{\mathrm{audit}}
  =
  \bigl(
    F_{\mathrm{HCOP}}^{\mathrm{audit}},
    \Delta_{\mathrm{audit}}^{\mathrm{freeze}},
    \Pi_{\mathrm{allowed}},
    S_{\mathrm{hold}}
  \bigr),
\]
where \(\Pi_{\mathrm{allowed}}\) is the finite permission set
\[
  \Pi_{\mathrm{allowed}}
  =
  \{\mathrm{audit},\mathrm{mechanicalRepair},\mathrm{metadataRepair},
    \mathrm{releasePackaging},\mathrm{zenodoDescription}\},
\]
and \(S_{\mathrm{hold}}\) is the hold marker blocking further theorem expansion before the release audit is completed.  The release label is metadata attached in prose and chronology, not a mathematical subscript or superscript of this object.
\end{definition}

\begin{lemma}[Freeze permissions exclude theorem expansion]
\label{lem:freeze-permissions-exclude-theorem-expansion}
While \(F_{\mathrm{freeze}}^{\mathrm{audit}}\) is active, the only admissible operations are audit, mechanical repair, metadata repair, release packaging, and Zenodo-description preparation.  New theorem expansion is not an admissible freeze operation.
\end{lemma}

\begin{proof}
The admissible operation set is exactly \(\Pi_{\mathrm{allowed}}\) in Definition~\ref{def:audit-freeze-stabilization-object}.  None of its entries is a proof-extension, closure-promotion, OP-claim, or certificate-generation operation.  Hence the active state may be repaired or packaged, but not mathematically expanded, until the release audit is discharged.
\end{proof}

\begin{theorem}[Audit-freeze stabilization]
\label{thm:audit-freeze-stabilization}
The current HC/OP working line is stabilized as an audit-freeze candidate by \(F_{\mathrm{freeze}}^{\mathrm{audit}}\).  In this state, the line is available for audit, mechanical repair, metadata repair, release packaging, and Zenodo-description preparation, while OP~1/4 and OP~5(B) remain open and no closure certificate is manufactured.
\end{theorem}

\begin{proof}
Theorem~\ref{thm:first-audit-freeze-ledger} supplies the finite ledger and audit-freeze defect.  Definition~\ref{def:audit-freeze-stabilization-object} adds only a finite permission set and a hold marker.  Lemma~\ref{lem:freeze-permissions-exclude-theorem-expansion} blocks theorem expansion during the freeze state.  Therefore the line is audit-stabilized without changing the open status of OP~1/4 or OP~5(B).
\end{proof}

\begin{corollary}[Audit-freeze working rule]
\label{cor:audit-freeze-working-rule}
Future work before the next public release may only audit, mechanically repair, repair metadata, package the release, or prepare the Zenodo description.  Any new mathematical theorem expansion must wait until the freeze is either released or explicitly reopened.
\end{corollary}

\begin{proof}
Immediate from Theorem~\ref{thm:audit-freeze-stabilization} and Lemma~\ref{lem:freeze-permissions-exclude-theorem-expansion}.
\end{proof}



\subsection*{Phase I-C: Higher Confinement and Septinion Unfolding}
\addcontentsline{toc}{subsection}{Phase I-C: Higher Confinement and Septinion Unfolding}

This third unfolding block turns closure stability into a higher confinement
reading and makes explicit why the septinion/seption interpretation is not an
external axiom but a structurally selected reading of the admissible closure
class.

\begin{lemma}[Higher confinement lemma]
\label{lem:higher-confinement-lemma}
If a structural class is closure-stable and remains invariant under repeated
compensator-compatible normalization, then it admits a higher confinement
reading.
\end{lemma}

\begin{proof}
Closure stability excludes uncontrolled drift out of the admissible class, while
repeated compensator-compatible normalization prevents the persistence of an
independent escaping sector. Once both conditions hold, the admissible class is
not merely locally normalized but globally retained under the structural return
mechanisms of the system. This is exactly the content of higher confinement:
the admissible states are no longer only corrected case by case but persist as a
retained confinement regime.
\end{proof}

\begin{lemma}[Non-escape lemma]
\label{lem:non-escape-lemma}
Let $x$ be a closure-stable state in the normalized core. Then no admissible
transition generated within the core sends $x$ into a structurally independent
external sector.
\end{lemma}

\begin{proof}
By closure stability, every admissible normalization step sends $x$ back into
the same closure class. If an admissible transition could produce a genuinely
independent external sector, then closure stability would fail exactly at that
point, because the state would no longer be returnable by the same structural
normalization. Hence admissible transitions may rearrange or re-express states
inside the confinement class, but they cannot export them into an independent
outside sector. This is the structural non-escape property.
\end{proof}

\begin{theorem}[Septinion readability theorem]
\label{thm:septinion-readability-theorem}
If a class is closure-stable and admits the higher confinement reading, then it
is septinion-readable: the resulting stabilized regime may be interpreted as a
septinion/seption-level confinement of the normalized structural body.
\end{theorem}

\begin{proof}
The septinion reading is not introduced as an arbitrary enlargement of the
formalism. Rather, it names the level at which the closure-stable and
non-escaping structural class is most naturally read as a higher unified
container of the admissible sectors. By Lemma~\ref{lem:higher-confinement-lemma},
closure stability already upgrades the system to a higher confinement regime,
and Lemma~\ref{lem:non-escape-lemma} shows that admissible states remain inside
that regime. The septinion/seption reading is therefore the natural structural
reading of this higher confinement container.
\end{proof}

\begin{corollary}[HC as emergent structural consequence]
\label{cor:hc-emergent-structural-consequence}
Within the normalized core, HC is not an externally appended axiom. It appears
as an emergent structural consequence of the one-body, frame-shift, and closure
requirements.
\end{corollary}

\begin{proof}
The generative chain of Phase~I-A forces re-identification and therefore HC,
while the closure chain of Phase~I-B shows that HC is precisely the mechanism
that stabilizes the admissible class. Phase~I-C now identifies the higher
confinement reading selected by this stabilization. Consequently HC is not an
optional addition but the structurally emergent mechanism by which the unified
body remains admissible across shifted readouts and under closure.
\end{proof}

\begin{corollary}[First confinement forcing chain]
\label{cor:first-confinement-forcing-chain}
The first unfolded confinement route of the normalized core is
\[
\text{closure class}
\Longrightarrow
\text{higher confinement}
\Longrightarrow
\text{non-escape}
\Longrightarrow
\text{septinion reading}.
\]
Together with the preceding phases, this yields
\[
\begin{aligned}
&\text{one body} \Longrightarrow \text{re-identification} \Longrightarrow \mathrm{HC}\\
&\Longrightarrow \text{closure} \Longrightarrow \text{higher confinement} \Longrightarrow \text{septinion/seption reading}.
\end{aligned}
\]
\end{corollary}

\begin{proof}
Proposition~\ref{prop:closure-class-proposition} provides the stabilized
closure class. Lemma~\ref{lem:higher-confinement-lemma} upgrades this class to a
higher confinement reading, Lemma~\ref{lem:non-escape-lemma} gives the internal
retention property of admissible transitions, and
Theorem~\ref{thm:septinion-readability-theorem} identifies the corresponding
septinion/seption interpretation. The final displayed chain is obtained by
concatenating this phase with the forcing chains of Phases~I-A and~I-B.
\end{proof}


\subsection*{Phase II-D: Cascade, Readout, and the Collapse-Fixed Sector}
\addcontentsline{toc}{subsection}{Phase II-D: Cascade, Readout, and the Collapse-Fixed Sector}

\begin{lemma}[Cascade admissibility lemma]
\label{lem:cascade-admissibility-lemma}
If $x$ belongs to the closure-stable class selected by the normalized core,
then every admissible cascade image $\hat K(x)$ remains inside the admissible
readout class of the same unified body.
\end{lemma}

\begin{proof}
The cascade is introduced as a composition of projection maps from higher
strata toward the ground sector, not as an ontological splitting of the body.
By Corollary~\ref{cor:first-closure-forcing-chain}, closure-stable states are
precisely those stabilized against drift under shifted readouts. Applying the
cascade to such a state therefore yields another admissible readout of the same
body rather than a new object. Hence admissibility is preserved along the
cascade image of any closure-stable input.
\end{proof}

\begin{lemma}[Terminal ground-state lemma]
\label{lem:terminal-ground-state-lemma}
The ground state $S_0$ is the terminal invariant readout sector of the cascade:
for every admissible state $x$, one has $\hat K(x)=x$ exactly when $x\in S_0$.
\end{lemma}

\begin{proof}
This is the normalized readout form of the ground-state condition already built
into the semantic kernel. The cascade terminates in $S_0$, and terminality
means precisely that no further nontrivial readout step occurs. Therefore the
fixed-point condition $\hat K(x)=x$ is equivalent to membership in the terminal
sector $S_0$.
\end{proof}

\begin{lemma}[Global null/light-axis lemma]
\label{lem:global-light-axis-lemma}
The null-axis structure selected by the normalized core, and read physically as the light axis at stable propagation, is stratified across
the cascade: if a state is read on this null/light axis at one admissible layer, then
its admissible readouts remain aligned with the same global null/light-axis class
through the cascade.
\end{lemma}

\begin{proof}
The Lorentz/null-axis structure, read physically as the light axis, is not introduced as a local afterthought but
as a stratified feature of the whole construction. The cascade only changes the
readout level, not the structural identity of the body. Since admissible
cascade images preserve the same unified body and the same closure-compatible
sector, the null/light-axis reading persists as a global class across the projected
layers.
\end{proof}

\begin{theorem}[Collapse-fixed sector theorem]
\label{thm:collapse-fixed-sector-theorem}
The photon/readout sector identified in the normalized core is the global
collapse-fixed sector of the cascade. In particular, the sector denoted by
$S_0$ is not merely a local null readout but the terminal collapse-fixed sector
selected by the full stratified structure.
\end{theorem}

\begin{proof}
Lemma~\ref{lem:cascade-admissibility-lemma} ensures that the cascade preserves
admissibility, while Lemma~\ref{lem:terminal-ground-state-lemma} identifies
$S_0$ as the terminal invariant sector. Lemma~\ref{lem:global-light-axis-lemma}
shows that the light-axis reading persists globally across admissible layers.
Therefore the photon/readout sector is not obtained from an isolated local
calculation alone; rather, it is the collapse-fixed terminal sector selected by
the entire cascade of admissible projections.
\end{proof}

\begin{lemma}[Readout separation lemma]
\label{lem:readout-separation-lemma}
Within the normalized core one must distinguish three levels: structural
forcing, cascade identification, and physical reading. A statement belongs to
exactly one of these levels unless an explicit bridge is given.
\end{lemma}

\begin{proof}
The structural forcing chain establishes what must hold inside the unified
body. The cascade then identifies how such a structurally forced state appears
at lower readout levels. A physical interpretation is an additional reading of
that identified sector. Conflating these levels obscures where a statement is
proved and where it is interpreted. The separation therefore serves as a
normalization rule for subsequent unfolding steps.
\end{proof}

\begin{corollary}[First readout forcing chain]
\label{cor:first-readout-forcing-chain}
The first unfolded readout route of the normalized core is
\[
\begin{aligned}
&\text{closure-stable class} \Longrightarrow \text{admissible cascade image}\\
&\Longrightarrow \text{terminal ground sector} \Longrightarrow \text{collapse-fixed photon/readout sector}.
\end{aligned}
\]
Taken together with Phases~I-A through~I-C, the normalized core now yields the
combined chain
\[
\begin{aligned}
&\text{one body} \Longrightarrow \text{re-identification} \Longrightarrow \mathrm{HC} \Longrightarrow \text{closure}\\
&\Longrightarrow \text{higher confinement} \Longrightarrow \text{septinion reading}\\
&\Longrightarrow \text{cascade-fixed readout sector}.
\end{aligned}
\]
\end{corollary}

\begin{proof}
The first displayed chain is the direct concatenation of
Lemma~\ref{lem:cascade-admissibility-lemma},
Lemma~\ref{lem:terminal-ground-state-lemma}, and
Theorem~\ref{thm:collapse-fixed-sector-theorem}. The second displayed chain is
obtained by appending this readout route to the forcing chain from
Corollary~\ref{cor:first-confinement-forcing-chain}.
\end{proof}


\subsection*{Phase II-E: Semantic/Tensor Interface and Millennium Intertwining}
\addcontentsline{toc}{subsection}{Phase II-E: Semantic/Tensor Interface and Millennium Intertwining}

\begin{lemma}[Semantic kernel well-formedness lemma]
\label{lem:semantic-kernel-well-formedness-lemma}
The normalized semantic kernel is well formed: its predicates of one-body
identity, triolic organization, orthogonal anchoring, frame-shifted readout,
HC forcing, closure stability, higher confinement, and readout admissibility do
not define competing ontologies but successive structural constraints on one and
the same body.
\end{lemma}

\begin{proof}
By construction, the normalized core begins from the one-body condition and then
adds only constraint predicates that regulate how the same body may be read,
stabilized, and projected. None of the predicates introduces a second carrier
space or an ontologically independent duplicate; each one refines the same
underlying body by imposing an admissibility condition. The semantic kernel is
therefore well formed as a constraint hierarchy rather than a family of
competing state spaces.
\end{proof}

\begin{lemma}[Tensor realization lemma]
\label{lem:tensor-realization-lemma}
The tensorial operators $\mathcal J^\mu{}_{\nu}$, $\mathcal C^\mu{}_{\nu}$,
$\Pi_n{}^\mu{}_{\nu}$, and $\hat K^\mu{}_{\nu}$ realize the normalized semantic
kernel as an operator system: $\mathcal J$ realizes the involutive frame
symmetry, $\mathcal C$ the compensating projection, $\Pi_n$ the cascade steps,
and $\hat K$ the terminal readout map.
\end{lemma}

\begin{proof}
The semantic kernel distinguishes involutive reflection, compensating return,
projection through admissible layers, and terminal ground readout. The tensorial
operators are introduced precisely to encode these roles in compact operator
notation. The semantic and tensorial descriptions therefore stand in a role-wise
correspondence: the operator system is not a second theory but a tensor
realization of the already normalized semantic structure.
\end{proof}

\begin{lemma}[Intertwining lemma]
\label{lem:millennium-intertwining-lemma}
If the Millennium matrix $\mathcal M^\mu{}_{\nu}$ is used as the structural
transport operator between admissible grouped layers, then it must intertwine
with the normalized tensorial kernel. In the minimal normalized form one
requires
\[
\mathcal M\mathcal J=\mathcal J\mathcal M,
\qquad
\mathcal M\mathcal C=\mathcal C\mathcal M,
\qquad
\mathcal M\hat K=\hat K\mathcal M
\]
whenever the corresponding involution, compensation, and cascade-readout
structures are to be preserved across the transport.
\end{lemma}

\begin{proof}
A transport operator that fails to commute with the involution would change the
frame logic of the transported state; one that fails to commute with the
compensator would alter admissibility by moving states in or out of the
compensated sector; and one that fails to commute with the terminal readout map
would not preserve the normalized collapse/readout architecture. Hence a
Millennium transport that is meant to be structure preserving must satisfy the
stated intertwining relations at least for the operators whose semantic role it
claims to preserve.
\end{proof}

\begin{theorem}[Admissibility preservation theorem]
\label{thm:admissibility-preservation-theorem}
If $x$ is admissible in the normalized semantic kernel and $\mathcal M$ is a
Millennium transport satisfying the intertwining relations of
Lemma~\ref{lem:millennium-intertwining-lemma}, then $\mathcal Mx$ remains
admissible in the transported layer.
\end{theorem}

\begin{proof}
Admissibility in the normalized core is carried by compatibility with the one-body
constraint, compensator normalization, and the cascade-compatible readout
structure. By Lemma~\ref{lem:millennium-intertwining-lemma}, the transport
operator preserves the tensorial realizations of exactly these mechanisms.
Therefore applying $\mathcal M$ cannot convert an admissible state into an
inadmissible one merely by transport; it yields an admissible image in the next
layer provided the transport is itself taken within the normalized grouped
architecture.
\end{proof}

\begin{theorem}[Closure preservation theorem]
\label{thm:closure-preservation-theorem}
Under the same hypotheses, closure-stable states are transported by
$\mathcal M$ to closure-stable states of the receiving layer.
\end{theorem}

\begin{proof}
Closure stability is the semantic statement that the admissible state does not
drift out of the normalized class once HC forcing and anchoring are taken into
account. Since admissibility is preserved by
Theorem~\ref{thm:admissibility-preservation-theorem} and the intertwining
relations preserve the compensator and readout structures that stabilize the
closure class, the transported image cannot lose closure merely by layer
transport. Thus closure stability is preserved across admissible Millennium
transport.
\end{proof}

\begin{proposition}[Semantic--tensor equivalence proposition]
\label{prop:semantic-tensor-equivalence-proposition}
The tensorial kernel does not introduce a second structural theory. It is a
parallel notation for the normalized semantic kernel, and the Millennium matrix
is admissible only insofar as its tensorial action represents the same
constraint chain already present on the semantic side.
\end{proposition}

\begin{proof}
By Lemma~\ref{lem:semantic-kernel-well-formedness-lemma}, the semantic side is a
single constrained body. By Lemma~\ref{lem:tensor-realization-lemma}, the
operator side realizes these constraints in tensorial notation. The intertwining
and preservation results show that $\mathcal M$ is not free to create a new
structural doctrine; it is licensed only as a transport operator respecting the
same normalized constraint chain. Hence semantic kernel and tensor kernel are
parallel realizations of one structure.
\end{proof}

\begin{corollary}[First semantic--tensor forcing chain]
\label{cor:first-semantic-tensor-forcing-chain}
The normalized semantic/tensor interface yields the chain
\[
\begin{aligned}
&\text{semantic kernel} \Longrightarrow \text{tensor realization} \Longrightarrow \text{Millennium intertwining}\\
&\Longrightarrow \text{admissibility preservation} \Longrightarrow \text{closure preservation}.
\end{aligned}
\]
\end{corollary}

\begin{proof}
Lemma~\ref{lem:semantic-kernel-well-formedness-lemma} establishes the semantic
kernel, Lemma~\ref{lem:tensor-realization-lemma} gives its tensor realization,
Lemma~\ref{lem:millennium-intertwining-lemma} fixes the admissible transport
relations for $\mathcal M$, and Theorems~\ref{thm:admissibility-preservation-theorem}
and~\ref{thm:closure-preservation-theorem} provide the preservation steps.
\end{proof}


\subsection*{Grouped Layer Architecture and Millennium Transport}
\addcontentsline{toc}{subsection}{Grouped Layer Architecture and Millennium Transport}

The semantic/tensor kernel admits a second normalization in which the document's
many local cores, layers, and readout blocks are collected into three grouped
structural classes. This reduces duplication and makes the role of the
Millennium matrix explicit as a transport operator between admissible groups.

\begin{definition}[Grouped structural architecture]
\label{def:grouped-structural-architecture}
The normalized structural architecture consists of the following three grouped
layers:
\begin{enumerate}[leftmargin=2em,label=(G\arabic*)]
\item \textbf{Generative layer $\mathfrak G$:} the layer of structural formation. It
contains the one-body condition, triolic minimality, orthogonal anchoring,
frame-shifted observation, the ground condition, HC forcing, closure
admissibility, and the higher-confinement/septinion reading.
\item \textbf{Projection/translation layer $\mathfrak P$:} the layer of cascade readout,
sector selection, projection rules, collapse-fixed identifications, null-/light-axis
readings, photon-sector readings, and admissible translation between equivalent
structural descriptions.
\item \textbf{Field/probe realization layer $\mathfrak F$:} the layer of probe states,
two-probe geometry, closure-preserving field equations, interaction rules, and
physical realization of the admissible structure.
\end{enumerate}
These three grouped layers are not independent ontologies. They are different
organizational regimes of the same underlying body and must therefore remain
compatible under admissible transport.
\end{definition}

\begin{definition}[Millennium transport between grouped layers]
\label{def:millennium-transport-grouped-layers}
The Millennium matrix is interpreted at the grouped-layer level as a
structure-preserving converter
\[
\mathcal M_{\mathfrak G\to\mathfrak P}:\mathfrak G\to\mathfrak P,
\qquad
\mathcal M_{\mathfrak P\to\mathfrak F}:\mathfrak P\to\mathfrak F,
\]
and, whenever reverse structural checking is required, also by admissible
back-transport maps
\[
\mathcal M_{\mathfrak P\to\mathfrak G}:\mathfrak P\to\mathfrak G,
\qquad
\mathcal M_{\mathfrak F\to\mathfrak P}:\mathfrak F\to\mathfrak P.
\]
At this level $\mathcal M$ is not merely a notational relabeling but a
converter which must preserve structural admissibility across layer changes.
\end{definition}

\begin{proposition}[Grouped forcing and readout chain]
\label{prop:grouped-forcing-readout-chain}
The grouped architecture organizes the document into the structural chain
\[
\mathfrak G \xrightarrow{\ \mathcal M\ } \mathfrak P
\xrightarrow{\ \mathcal M\ } \mathfrak F,
\]
where:
\begin{enumerate}[leftmargin=2em,label=(\alph*)]
\item $\mathfrak G$ generates the admissible structure,
\item $\mathfrak P$ reads, projects, and translates that structure without changing
its admissibility class,
\item $\mathfrak F$ realizes the same admissible structure as a field/probe system.
\end{enumerate}
In particular, the field/probe layer is not a second origin of the theory; it
is a realization layer of structure first generated in $\mathfrak G$ and made
readable in $\mathfrak P$.
\end{proposition}

\begin{proof}
The generative layer contains the forcing data established above: one-body,
triolic organization, orthogonal anchor, frame shift, HC, and closure. The
projection/translation layer does not create new admissibility; it converts the
generative core into cascade-compatible, sector-resolved, and physically
legible readout. The field/probe layer then realizes those admissible readouts
as explicit probes, field equations, or interaction geometries. Since the same
body is being followed through all three layers, admissible transport by
$\mathcal M$ must preserve rather than replace the structural class.
\end{proof}

\begin{proposition}[Admissibility preservation across grouped layers]
\label{prop:admissibility-preservation-grouped-layers}
Let $x$ be an admissible state generated in $\mathfrak G$. If $\mathcal M$ acts as
an admissible grouped-layer converter, then
\[
\mathrm{Admissible}(x)
\Longrightarrow
\mathrm{Admissible}(\mathcal M_{\mathfrak G\to\mathfrak P}x)
\Longrightarrow
\mathrm{Admissible}(\mathcal M_{\mathfrak P\to\mathfrak F}\mathcal M_{\mathfrak G\to\mathfrak P}x).
\]
Moreover, when closure stability is part of the admissibility data, it is to be
preserved throughout this transport chain.
\end{proposition}

\begin{proof}
By Definition~\ref{def:millennium-transport-grouped-layers}, grouped-layer
transport is only allowed when it preserves structural admissibility. Hence an
admissible state generated in $\mathfrak G$ remains admissible after conversion
into the projection/translation layer, and likewise after conversion into the
field/probe layer. If closure stability were lost during transport, the result
would no longer be a valid readout or realization of the original structure but
a structural rupture. Therefore closure stability must be preserved whenever it
belongs to the originating admissibility class.
\end{proof}

\begin{remark}[Document-normalization consequence]
\label{rem:document-normalization-consequence}
The grouped architecture provides a normalization principle for the text. Local
``cores'' or ``integration blocks'' should be read either as part of
$\mathfrak G$, as part of $\mathfrak P$, or as part of $\mathfrak F$, rather than as
independent foundational centers. This removes duplication: the normative
foundation lies in $\mathfrak G$, the translation/readout logic in $\mathfrak P$,
and the field/probe realization in $\mathfrak F$.
\end{remark}

\subsection*{Phase II-F: Grouped Layer Transport and Ontological Non-Drift}
\addcontentsline{toc}{subsection}{Phase II-F: Grouped Layer Transport and Ontological Non-Drift}

This unfolding block upgrades the grouped architecture from a normalization note
into an explicit transport calculus. The goal is to show that grouped layers do
not introduce new origins of structure, but only different admissible regimes of
one and the same underlying body.

\begin{lemma}[Group decomposition lemma]
\label{lem:group-decomposition}
Every active structural block of the normalized Companion is dominantly
assignable to exactly one of the grouped layers $\mathfrak G$, $\mathfrak P$, or
$\mathfrak F$, even when secondary cross-references to the other layers remain.
\end{lemma}

\begin{proof}
By Definition~\ref{def:grouped-structural-architecture}, the grouped layers are
distinguished by role rather than by separate ontology: generation,
projection/translation, and realization. Any active block therefore has one
primary structural function. Cross-references may connect a block to the other
layers, but they do not erase its dominant role. Hence each active block is
assignable to one grouped layer as its principal location.
\end{proof}

\begin{lemma}[Transport lemma]
\label{lem:group-transport}
Let $x$ be a state generated in $\mathfrak G$. If $\mathcal M_{\mathfrak G\to\mathfrak P}$
is an admissible grouped-layer converter, then its image in $\mathfrak P$ is a
projection/translation of the same admissible structure and not a newly created
structural body.
\end{lemma}

\begin{proof}
By Definition~\ref{def:millennium-transport-grouped-layers}, admissible grouped
transport preserves structural admissibility. Proposition~\ref{prop:grouped-forcing-readout-chain}
states that $\mathfrak P$ reads and translates what $\mathfrak G$ has already
generated. Therefore the image of $x$ under $\mathcal M_{\mathfrak G\to\mathfrak P}$
remains the same admissible structure under a projection/translation regime.
It is not a second origin because admissibility is transported, not recreated.
\end{proof}

\begin{lemma}[Realization lemma]
\label{lem:group-realization}
If $y\in\mathfrak P$ is an admissible projection/translation state and
$\mathcal M_{\mathfrak P\to\mathfrak F}$ is admissible, then its image in
$\mathfrak F$ is a realization of the same admissible structure as a field/probe
state.
\end{lemma}

\begin{proof}
The role of $\mathfrak F$ is realization rather than generation by
Proposition~\ref{prop:grouped-forcing-readout-chain}. Since admissibility is
preserved across grouped transport by Proposition~\ref{prop:admissibility-preservation-grouped-layers},
$\mathcal M_{\mathfrak P\to\mathfrak F}(y)$ remains within the same
admissibility class. Thus the field/probe layer realizes the translated
structure without changing its underlying body.
\end{proof}

\begin{proposition}[No-ontology-drift proposition]
\label{prop:no-ontology-drift}
Admissible transport across the grouped layers does not create ontological
multiplicity. If
\[
x \in \mathfrak G,
\qquad
x_P := \mathcal M_{\mathfrak G\to\mathfrak P}(x),
\qquad
x_F := \mathcal M_{\mathfrak P\to\mathfrak F}(x_P),
\]
then $x$, $x_P$, and $x_F$ are three admissible organizational regimes of one
and the same structural body.
\end{proposition}

\begin{proof}
By Lemma~\ref{lem:group-transport}, passage from $\mathfrak G$ to $\mathfrak P$
produces a translation/readout of the same admissible body. By
Lemma~\ref{lem:group-realization}, passage from $\mathfrak P$ to $\mathfrak F$
produces a field/probe realization of that same admissible body. Since both
steps preserve admissibility rather than replacing it, no new ontology is
introduced along the chain. Therefore $x$, $x_P$, and $x_F$ are structurally
identical in body and differ only by organizational regime.
\end{proof}

\begin{corollary}[First grouped transport forcing chain]
\label{cor:first-grouped-transport-forcing-chain}
The normalized grouped architecture supports the forcing chain
\[
\begin{aligned}
&\mathfrak G \Longrightarrow \text{admissible transport to }\mathfrak P\\
&\Longrightarrow \text{admissible realization in }\mathfrak F \Longrightarrow \text{no ontology drift across grouped layers}.
\end{aligned}
\]
\end{corollary}

\begin{proof}
Combine Lemma~\ref{lem:group-decomposition}, Lemma~\ref{lem:group-transport},
Lemma~\ref{lem:group-realization}, and
Proposition~\ref{prop:no-ontology-drift}.
\end{proof}


\subsection*{Phase III-G: Probe Geometry, Inheritance, and Guided Search}
\addcontentsline{toc}{subsection}{Phase III-G: Probe Geometry, Inheritance, and Guided Search}

This unfolding block upgrades the probe sector from an interpretive add-on to a
structurally inherited realization of the normalized core. The goal is to show
that probes do not introduce a foreign search mechanism, but inherit triolic,
anchored, and closure-compatible structure from the generative and grouped
layers.

\begin{lemma}[Probe inheritance lemma]
\label{lem:probe-inheritance}
If the normalized structural body is triolic, anchored, and closure-admissible,
then every admissible probe realization inherited through $\mathfrak F$ carries a
corresponding triolic and closure-compatible organization.
\end{lemma}

\begin{proof}
By Proposition~\ref{prop:no-ontology-drift}, admissible realization in
$\mathfrak F$ does not create a second ontology, but realizes the same
underlying structural body under a field/probe regime. Hence structural features
that belong to the admissible body---in particular triolic organization and
anchor-compatible closure discipline---remain available in the probe image as
inherited structure rather than externally attached decoration.
\end{proof}

\begin{lemma}[Two-probe compatibility lemma]
\label{lem:two-probe-compatibility}
Let $P_1$ and $P_2$ be two admissible probes directed toward one another. If
both are inherited from the same normalized structural body, then their mutual
configuration is compatible with one common closure-governed geometry.
\end{lemma}

\begin{proof}
By Lemma~\ref{lem:probe-inheritance}, each probe inherits the same triolic and
closure-compatible structural discipline. Since neither probe introduces a new
ontology by Proposition~\ref{prop:no-ontology-drift}, their mutual interaction
must be read inside one common admissibility class. Therefore the two-probe
configuration is not a collision of unrelated structures but a common
closure-governed geometry of one inherited body.
\end{proof}

\begin{lemma}[Orthogonal third-group lemma]
\label{lem:orthogonal-third-group-probe}
In an admissible two-probe configuration, the third structural group remains
orthogonal to both probes and thereby determines the constrained search region.
\end{lemma}

\begin{proof}
The orthogonal third direction is already forced at the generative level by the
orthogonal-anchor principle and remains preserved under admissible realization.
Hence, when two probes are realized as inherited triolic structures, the third
group cannot collapse into either probe direction without violating anchor
compatibility. It therefore remains orthogonal to both and serves as the
stabilizing transversal direction that bounds the admissible search region.
\end{proof}

\begin{proposition}[Probe reduction theorem]
\label{prop:probe-reduction-theorem}
An admissible two-probe configuration reduces to the closure-preserving sector of
the grouped architecture. In particular, the probe interaction may be analyzed
inside the admissible closure class without introducing a second search
ontology.
\end{proposition}

\begin{proof}
By Lemma~\ref{lem:two-probe-compatibility}, the two probes share one common
closure-governed geometry. By Lemma~\ref{lem:orthogonal-third-group-probe}, the
third group fixes the transversal admissible search region. Since closure and
ontology are preserved across grouped realization by
Proposition~\ref{prop:no-ontology-drift}, the probe interaction reduces to the
closure-preserving sector of the same normalized architecture.
\end{proof}

\begin{corollary}[Guided search corollary]
\label{cor:guided-search-corollary}
Auxiliary tools such as frequency extraction, transition guidance, structural
routing, or convergence acceleration act only within the already admissible
probe geometry. They do not generate admissibility, but exploit a search region
that has already been restricted by inheritance, closure, and orthogonal
anchoring.
\end{corollary}

\begin{proof}
By Proposition~\ref{prop:probe-reduction-theorem}, admissible probe analysis is
already confined to the closure-preserving sector. Therefore any auxiliary tool
operates on a domain whose admissibility has been fixed before the tool is
applied. The tool may compress or guide the search, but it does not create the
search geometry itself.
\end{proof}

\begin{corollary}[First probe forcing chain]
\label{cor:first-probe-forcing-chain}
The normalized probe sector supports the forcing chain
\[
\begin{aligned}
&\text{inherited probe structure}
  \Longrightarrow \text{two-probe compatibility}\\
&\Longrightarrow \text{orthogonal third-group search region}\\
&\Longrightarrow \text{closure-preserving probe reduction}\\
&\Longrightarrow \text{guided search inside an already admissible geometry}.
\end{aligned}
\]
\end{corollary}

\begin{proof}
Combine Lemma~\ref{lem:probe-inheritance},
Lemma~\ref{lem:two-probe-compatibility},
Lemma~\ref{lem:orthogonal-third-group-probe},
Proposition~\ref{prop:probe-reduction-theorem}, and
Corollary~\ref{cor:guided-search-corollary}.
\end{proof}



\subsection*{Phase III-H: Structural Master Equation and Exact-Closure Field Regime}
\addcontentsline{toc}{subsection}{Phase III-H: Structural Master Equation and Exact-Closure Field Regime}

This unfolding block upgrades the previously compressed master-equation language
into an explicit field-level forcing chain. The aim is not yet to derive every
physical limit in full detail, but to show that the structural master equation,
its variational reading, the closure-preserving field regime, and the exact-
closure simplifications belong to one continuous argument rather than to loosely
connected interpretive remarks.

\begin{lemma}[Structural master-equation lemma]
\label{lem:structural-master-equation}
The structural master equation is the compressed field-level summary of the
normalized forcing chain from one-body structure through re-identification, HC,
closure, higher confinement, and admissible readout.
\end{lemma}

\begin{proof}
By Corollary~\ref{cor:first-generative-forcing-chain}, the normalized body
forces re-identification and HC. By Corollary~\ref{cor:first-closure-forcing-chain},
this induces a closure class with either exact closure or a controlled residue.
By Corollary~\ref{cor:first-confinement-forcing-chain}, closure yields higher
confinement and the septinion reading. By Corollary~\ref{cor:first-readout-forcing-chain}
and Corollary~\ref{cor:first-semantic-tensor-forcing-chain}, admissible readout
and tensorial transport preserve this structure. Hence the master equation is not
an additional postulate but the compressed field-level expression of the forcing
chain already established.
\end{proof}

\begin{lemma}[Variational consistency lemma]
\label{lem:variational-consistency}
Any variational reading of the structural master equation is admissible only if
it preserves the normalized one-body, HC, and closure conditions. In
particular, admissible variation may unfold the structural equation but may not
replace its forcing chain by an independent ontology.
\end{lemma}

\begin{proof}
By Proposition~\ref{prop:semantic-tensor-equivalence-proposition}, tensorial formulation is
only a parallel realization of the same semantic kernel. Therefore any
variational formulation that is meant to belong to the present architecture must
preserve the same semantic admissibility data. If the variation were to destroy
one-body identity, HC-forcing, or closure preservation, it would cease to be a
variation of the normalized structural equation and would instead define a
second independent system. Hence admissible variation is constrained by the
forcing chain and remains structurally internal to it.
\end{proof}

\begin{lemma}[Closure-field compatibility lemma]
\label{lem:closure-field-compatibility}
The field-level realization of the structural master equation is compatible with
closure precisely when the realized dynamics preserves the HC-selected
admissibility class and does not mix the stable closure sector with an
ontologically foreign complement.
\end{lemma}

\begin{proof}
By Theorem~\ref{thm:closure-preservation-theorem}, admissible transport through
the Millennium interface preserves closure classes. By Proposition~\ref{prop:no-ontology-drift},
realization in the field/probe layer introduces no second ontology. Therefore a
field realization is closure-compatible exactly when it remains inside the
HC-selected admissibility class rather than coupling the closure-stable sector
to a foreign complement. This is the field-level restatement of normalized
closure preservation.
\end{proof}

\begin{theorem}[Exact-closure field theorem]
\label{thm:exact-closure-field-theorem}
In the exactly closed regime, the realized field dynamics reduces to the
closure-stable block of the normalized architecture. Equivalently, the active
field transport becomes block-compatible with the exact closure class.
\end{theorem}

\begin{proof}
By Proposition~\ref{prop:exact-closure-criterion}, exact closure is equivalent
to the disappearance of the nontrivial closure defect. By
Lemma~\ref{lem:closure-field-compatibility}, closure-compatible field dynamics
may then act only on the already stabilized admissibility class. Since no
ontology drift occurs across grouped transport, the realized field dynamics
reduces to the closure-stable block itself. Thus exact closure yields a
block-compatible field regime.
\end{proof}

\begin{lemma}[Minimal compensator disappearance lemma]
\label{lem:minimal-compensator-disappearance}
In the exactly closed regime, the minimal compensator residue vanishes. What
remains is not the disappearance of structural compensation as such, but the
fact that no additional compensator correction is required beyond the exact
closure block.
\end{lemma}

\begin{proof}
By Lemma~\ref{lem:closure-residue-lemma}, the non-exact regime is marked
by a residual deviation from exact closure. By
Theorem~\ref{thm:exact-closure-field-theorem}, exact closure confines the field
transport completely to the closure-stable block. Therefore the additional
minimal compensator term that measures residual off-block correction must
vanish. Structural compensation remains implicit in the exact block, but no
separate compensator residue survives.
\end{proof}

\begin{theorem}[Effective-limit theorem]
\label{thm:effective-limit-theorem}
Any admissible effective limit---for example a geometric, spinorial, or quantum
field-theoretic limit---must arise as a readout of the normalized structural
master equation together with its closure-preserving field realization. Such a
limit may specialize the realization, but cannot bypass the forcing chain that
produces it.
\end{theorem}

\begin{proof}
By Lemma~\ref{lem:structural-master-equation}, the structural master equation is
the compressed summary of the normalized forcing chain. By
Lemma~\ref{lem:variational-consistency} and Lemma~\ref{lem:closure-field-compatibility},
any admissible field realization and any admissible variation remain internal to
that chain. By Theorem~\ref{thm:exact-closure-field-theorem} and
Lemma~\ref{lem:minimal-compensator-disappearance}, the exact regime yields a
pure closure block without additional corrective residue. Hence any effective
limit must be read as a specialization or projection of the normalized
structural field realization rather than as an independent starting point.
\end{proof}

\begin{corollary}[First field forcing chain]
\label{cor:first-field-forcing-chain}
The normalized field sector supports the forcing chain
\[
\begin{aligned}
&\text{structural master equation} \Longrightarrow \text{variational consistency}\\
&\Longrightarrow \text{closure-preserving field compatibility} \Longrightarrow \text{exact closure block regime}\\
&\Longrightarrow \text{vanishing minimal compensator residue} \Longrightarrow \text{admissible effective limit readout}.
\end{aligned}
\]
\end{corollary}

\begin{proof}
Combine the structural master equation, variational consistency, closure-field compatibility, the exact closure-field theorem, minimal compensator disappearance, and the effective limit theorem.
\end{proof}

\subsection*{Concrete mapping of legacy core blocks to grouped layers}
\addcontentsline{toc}{subsection}{Concrete mapping of legacy core blocks to grouped layers}

The grouped normalization can be made explicit by assigning the surviving legacy
core blocks to the three grouped layers $\mathfrak G$, $\mathfrak P$, and
$\mathfrak F$. This turns older version-labelled cores into either active parts
of the normalized architecture or historical/supporting blocks.

\begin{definition}[Layer assignment of legacy core blocks]
\label{def:layer-assignment-legacy-core-blocks}
Within the present monolith, the principal legacy blocks are to be read as
follows.
\begin{enumerate}[leftmargin=2em,label=(L\arabic*)]
\item \textbf{$\mathfrak G$-dominant blocks (generative layer):}
\begin{enumerate}[leftmargin=2.5em,label=(\alph*)]
\item the normalized \emph{Structural Generation Core},
\item \emph{Structural Core and Closure Principles},
\item the generative part of \emph{v1.2.0 Consolidation Core}, especially the
structural reading discipline together with the ground-condition/HC/closure and
higher-confinement logic,
\item the generative part of \emph{v1.4.0 Derivation and Appendix Integration Core},
especially the closure ladder and the micro-to-macro derivation chain.
\end{enumerate}
\item \textbf{$\mathfrak P$-dominant blocks (projection/translation layer):}
\begin{enumerate}[leftmargin=2.5em,label=(\alph*)]
\item \emph{v1.5.0 Projection and Appendix Table Integration Core},
\item the appendix-coupling and projection-summary parts of \emph{v1.4.0 Derivation and Appendix Integration Core},
\item the collapse, light-axis, photon-sector, and cascade-readout parts of the later main sections,
\item all readout, table, projection, and sector-translation interfaces whose
role is to make one admissible structure legible without changing its
admissibility class.
\end{enumerate}
\item \textbf{$\mathfrak F$-dominant blocks (field/probe realization layer):}
\begin{enumerate}[leftmargin=2.5em,label=(\alph*)]
\item the probe-sector and two-probe geometry blocks,
\item the field-equation, interaction, and closure-preserving dynamics sections,
\item the probe-reduction and guided-search realizations insofar as they are
used as operational realizations of previously fixed admissible structure.
\end{enumerate}
\item \textbf{Historical/supporting blocks:} readiness criteria, version-status
passages, and older local consolidation summaries are not to be read as new
foundational centers. They are historical normalization records unless they are
explicitly re-exported into one of $\mathfrak G$, $\mathfrak P$, or
$\mathfrak F$.
\end{enumerate}
\end{definition}

\begin{proposition}[Normalization map for duplicated core material]
\label{prop:normalization-map-duplicated-core-material}
Under the grouped architecture, duplicated structural material is normalized by
the following rule.
\begin{enumerate}[leftmargin=2em,label=(\roman*)]
\item When a legacy block states formation conditions, forcing rules, or
closure-generating principles, it is normalized into $\mathfrak G$.
\item When a legacy block states projection rules, cascade-compatible readout,
sector identifications, appendix transport, or summary-table logic, it is
normalized into $\mathfrak P$.
\item When a legacy block states probes, field equations, interaction channels,
or closure-preserving realizations, it is normalized into $\mathfrak F$.
\item If a block contains mixed material, it is to be read by dominant function,
with the remaining content treated as support, historical context, or material
that should later be split explicitly.
\end{enumerate}
In particular, \emph{v1.2.0 Consolidation Core} and the early derivation cores no
longer define independent foundations beside the normalized structural core;
they contribute either to $\mathfrak G$, to $\mathfrak P$, or to historical
context.
\end{proposition}

\begin{remark}[Immediate practical consequence for the next clean release]
\label{rem:practical-consequence-next-clean-release}
The next clean release should therefore not keep several local cores in
competition. Instead, older version-labelled cores are to be annotated and,
where necessary, shortened according to their dominant grouped-layer function:
formation material belongs to $\mathfrak G$, translation/readout material to
$\mathfrak P$, field/probe realization material to $\mathfrak F$, and pure
status/history material to archival support.
\end{remark}


